% Appendix C, part 1: deferred proofs of the numbered results of Exposés 0, I and II.
% Sources: theory/propositions.json (repaired proofs, public statements), theory/framework.md (Exposés I and II),
% theory/round2_verdicts.json (suggested fixes, applied), and the final statements in sections/01-slice.tex,
% 02a-strata.tex, 02b-erlangen.tex and 02c-cut.tex. External results are Auxiliary results B.1--B.39 of Appendix B.
\section{Deferred proofs}\label{app:proofs}

This appendix proves the numbered results of the main text in the order in which they appear there, one subsection per
result. Each subsection first restates the result in its final form, the form recorded after the two rounds of review,
and then proves it. External theorems enter through Appendix~\ref{app:auxiliary}, cited as Auxiliary results B.1--B.39,
and background notions through Appendix~\ref{app:background}. Exposé~\ref{sec:thesis} states its headline results
without numbers; each of them is proved below or in the later parts of this appendix under its number. Definitions I.1
and I.2 need no proof, since $\Para(\mathcal C)$ is a bicategory \citep{fong2017backprop,cruttwell2021categorical} and
$\Meth(\Theta,\theta_0)$ is a slice of a coslice; Remark I.3 and the two definitions that contain a claim, Definitions
I.5 and II.14, receive short proofs. Unless we say otherwise, parameter spaces are open definable subsets of a vector
space and the maps $\rho$ are semialgebraic or definable in $\R_{\exp}$ (\bgref{C.17}), dimensions are those of
\auxref{B.9}, ``generic'' has the meaning of \auxref{B.10}, and LoRA on one $m\times n$ box is
$\rho(B,A)=W_{\mathrm{res}}+sBA$ with $s\ne0$, $B\in\R^{m\times r}$ and $A\in\R^{r\times n}$. For subspaces
$V\subseteq\R^n$ and $U\subseteq\R^m$ we write $\mathcal R(V)=\{Z\in\R^{m\times n}:\operatorname{row}Z\subseteq V\}$ and
$\mathcal C(U)=\{Z:\operatorname{col}Z\subseteq U\}$, and $x\wedge y=xy\T-yx\T$. Values marked [Num] in a statement are
checked numerically and not proved; Appendix~\ref{app:repro} documents them.

% ======================================================================================================================
\subsection{Remark I.3}\label{proof:I.3}

\begin{restate}{Remark}{I.3}{orientation, and what Para adds}
Each $\rho$ is a 2-cell $F\Rightarrow F^\rho:=(Q,f\circ(\rho\times X))$ \emph{out of} $F$. A simulation $h:M\to N$ is a
2-cell $F^{\rho_N}\Rightarrow F^{\rho_M}$, since $\rho_N\circ h=\rho_M$. Hence $\Meth(\Theta,\theta_0)$ is the opposite
of the pointed coslice of reparametrisation 2-cells under $F$. In particular $\Meth(\Theta,\theta_0)$ depends only on
$(\Theta,\theta_0)$ and not on $f$.
\end{restate}

\begin{proof}
Let $\mathcal P$ be the category whose objects are the 1-cells $X\to Y$ of $\Para(\Diff)$ and whose arrows are the
2-cells, composed vertically. By \thmref{Definition}{I.1} a 2-cell $(P,f)\Rightarrow(Q,g)$ is a smooth map $r:Q\to P$ with
$g=f\circ(r\times X)$. The vertical composite of $r:(P,f)\Rightarrow(Q,g)$ and $r':(Q,g)\Rightarrow(Q',g')$ is
$r\circ r':Q'\to P$, a 2-cell because $g'=g\circ(r'\times X)=f\circ((r\circ r')\times X)$.

An object of the coslice $F/\mathcal P$ is a 2-cell out of $F=(\Theta,f)$. A 2-cell from $F$ to $(Q,g)$ is a smooth
$\rho:Q\to\Theta$ with $g=f\circ(\rho\times X)$, so its target is forced to be $F^\rho$ and the object is determined by the
pair $(Q,\rho)$. An arrow from $\rho_M:F\Rightarrow F^{\rho_M}$ to $\rho_N:F\Rightarrow F^{\rho_N}$ is a 2-cell
$k:F^{\rho_M}\Rightarrow F^{\rho_N}$ whose vertical composite with $\rho_M$ is $\rho_N$. Concretely, it is a smooth map
$k:Q_N\to Q_M$ with
\begin{equation*}
  f\circ(\rho_N\times X)=f\circ(\rho_M\times X)\circ(k\times X)\qquad\text{and}\qquad\rho_M\circ k=\rho_N .
\end{equation*}
The second condition implies the first. So the arrows $(Q_M,\rho_M)\to(Q_N,\rho_N)$ of $F/\mathcal P$ are exactly the
smooth maps $k:Q_N\to Q_M$ over $\Theta$, which are the arrows $N\to M$ of $\Diff/\Theta$, and composition is reversed.
Restricting to pointed parameter spaces with $\rho(q_0)=\theta_0$ and to pointed $k$ gives the pointed coslice, which is
therefore $\Meth(\Theta,\theta_0)^{\mathrm{op}}$. The network $f$ entered only through the first condition, which is
implied by the second for every $f$, so $\Meth(\Theta,\theta_0)$ depends only on $(\Theta,\theta_0)$.
\end{proof}

% ======================================================================================================================
\subsection{Observation I.4}\label{proof:I.4}

\begin{restate}{Observation}{I.4}{universal properties}
In $\Meth(\Theta,\theta_0)$:

(a) the frozen model $(\mathrm{pt},*,\theta_0)$ is initial, and full fine-tuning $(\Theta,\theta_0,\id)$ is terminal. The
unique morphism $M\to{}$Full FT is $\rho_M$ itself.

(b) If the fibre product $Q_M\times_\Theta Q_N$ exists in $\Diff$, it is the product $M\times N$. It exists, for instance,
when the set-level fibre product is an embedded submanifold of $Q_M\times Q_N$: a clean intersection, or linear or affine
$\rho$'s as in \thmref{Proposition}{I.8}(a). Transversality $\rho_M\pitchfork\rho_N$ is never the reason when
$\lvert M\rvert+\lvert N\rvert<\dim\Theta$, since at $(q_0,q_0')$ it would force $\lvert M\rvert+\lvert
N\rvert\ge\dim\Theta$. In $\Set$ the fibre product always exists and $\Im(M\times N)=\Im M\cap\Im N$. Whenever the product
exists in $\Meth$, its underlying set is this set-level fibre product, so the image identity holds in every case. The
product can fail to exist. On $\Theta=\R$ ($1\times1$ matrices), $\mathrm{LoRA}_1\times\text{Frozen}$ would have the cross
$\{BA=0\}\subset\R^2$ as underlying set. No manifold structure on the cross has the universal property, because both axes
would have to lift smoothly through the point over the origin, which forces rank $2$ there.

(c) $\Im$ is a functor to the poset of subsets of $\Theta$ containing $\theta_0$. Conversely, $\Im M\subseteq\Im N$ yields
a morphism of the set-level slice $\Set_*/(\Theta,\theta_0)$, though not necessarily a smooth one. Take $\rho_M(t)=t$ and
$\rho_N(u)=u^3$ on $\R$, pointed at $0$. The images agree, yet no smooth $h:\R\to\R$ satisfies $\rho_N\circ h=\id$, that
is, $h(t)^3=t$.
\end{restate}

\begin{proof}
Recall that $\Meth(\Theta,\theta_0)=\Diff_*/(\Theta,\theta_0)$ (\thmref{Definition}{I.2}).

\emph{(a)} Let $M=(Q,q_0,\rho)$. An arrow from the frozen model to $M$ is a pointed smooth map $h:\mathrm{pt}\to Q$ with
$\rho\circ h=\theta_0$. Pointedness forces $h(*)=q_0$, and the condition over $\Theta$ then holds because
$\rho(q_0)=\theta_0$. So there is exactly one such arrow. An arrow $h:M\to(\Theta,\theta_0,\id)$ is a pointed smooth map
$h:Q\to\Theta$ with $\id\circ h=\rho$. The only candidate is $h=\rho$, which is smooth and pointed.

\emph{(b)} By \auxref{B.3}(a), a product of $M$ and $N$ in the slice $\Diff_*/(\Theta,\theta_0)$ is the same as a pullback
of $\rho_M$ and $\rho_N$ in $\Diff_*$. Suppose the pullback $P$ of the underlying maps exists in $\Diff$, with projections
$\pi_M,\pi_N$. The constant maps $\mathrm{pt}\to Q_M$ and $\mathrm{pt}\to Q_N$ at $q_0$ and $q_0'$ agree over $\Theta$,
both composing to $\theta_0$, so they factor through a unique point $p_0\in P$. With base point $p_0$, $P$ is the pullback
in $\Diff_*$. For a pointed test object $T$ and pointed maps $h_M,h_N$ that agree over $\Theta$, the induced map $T\to P$
sends the base point to a point whose projections are $q_0$ and $q_0'$, and that point is $p_0$ by uniqueness. This is
the creation of limits by $\Diff_*\to\Diff$ (\auxref{B.3}(b)).

If the set-level fibre product $S=\{(q,q'):\rho_M(q)=\rho_N(q')\}$ is an embedded submanifold of $Q_M\times Q_N$, then $S$
with its two projections is the pullback in $\Diff$ (\auxref{B.2}(c)). Indeed, two smooth maps $T\to Q_M$ and $T\to Q_N$
that agree over $\Theta$ give a smooth map $T\to Q_M\times Q_N$ with image in $S$, which is smooth into $S$, and the
factorisation is unique because $S\to Q_M\times Q_N$ is injective. A clean intersection has, by definition, an embedded
submanifold $S$ (\bgref{C.7}), and for affine $\rho_M$ and $\rho_N$ on vector spaces $S$ is an affine subspace.

Transversality at $(q_0,q_0')$ asks $d\rho_M(T_{q_0}Q_M)+d\rho_N(T_{q_0'}Q_N)=T_{\theta_0}\Theta$. The left side has
dimension at most $\lvert M\rvert+\lvert N\rvert$, so transversality fails at the base point whenever
$\lvert M\rvert+\lvert N\rvert<\dim\Theta$ (\auxref{B.2}(b)).

In $\Set$ the fibre product is $S$, and its image in $\Theta$ is $\Im M\cap\Im N$: a pair $(q,q')\in S$ lies over
$\rho_M(q)=\rho_N(q')$, which belongs to both images, and a point $\theta$ of both images has preimages $q$ and $q'$ with
$(q,q')\in S$. Now suppose that a product $P$ exists in $\Meth$, and test it with the pointed $0$-manifold
$T=\{*,t\}$, pointed at $*$. A pointed map $T\to P$ is the choice of the image of $t$, that is, a point of $P$. A pair of
pointed maps $T\to Q_M$ and $T\to Q_N$ that agree over $\Theta$ is a point $(q,q')$ of $S$. The universal property makes
these correspond bijectively through $p\mapsto(\pi_M(p),\pi_N(p))$. So the canonical map $P\to Q_M\times Q_N$ is a bijection
onto $S$, and $\Im(M\times N)=\rho_M(\pi_M(P))=\Im M\cap\Im N$.

Take $\Theta=\R$ and $\mathrm{LoRA}_1:(B,A)\mapsto\theta_0+sBA$ on $\R^2$, with $s\ne0$ and frozen residual $\theta_0$,
so that its pointings are the $(B_0,A_0)$ with $B_0A_0=0$. Its set-level fibre product with the frozen model is the cross
$\mathcal X=\{(B,A):BA=0\}$. Suppose a product $P$ exists, and let $\iota:P\to\R^2$ be its projection to the parameters
of $\mathrm{LoRA}_1$. The forgetful functor $\Diff_*\to\Diff$ has a left adjoint, so it preserves this pullback
(\auxref{B.3}(c)), and $P$ is also the pullback in $\Diff$ of $\rho_{\mathrm{LoRA}_1}$ and the inclusion of $\theta_0$.
Hence every smooth curve $c:\R\to\R^2$ with $\rho\circ c\equiv\theta_0$ has a unique smooth lift $\tilde c:\R\to P$ with
$\iota\circ\tilde c=c$, and by the previous paragraph $\iota$ is a bijection onto $\mathcal X$. Let $p^*$ be the point of
$P$ over $(0,0)$. The axis curves $c_1(t)=(t,0)$ and $c_2(t)=(0,t)$ lie in $\mathcal X$, so they lift to smooth curves
through $p^*$, and $d\iota_{p^*}$ maps their velocities to $c_1'(0)=(1,0)$ and $c_2'(0)=(0,1)$. Hence
$\rk d\iota_{p^*}=2$. The rank cannot exceed $2$ and cannot drop near $p^*$ (\auxref{B.1}(a)), so $\iota$ is a submersion
on a neighbourhood $V$ of $p^*$, and $\iota(V)$ contains a nonempty open subset of $\R^2$ (\auxref{B.1}(c)). But
$\iota(V)\subseteq\mathcal X$, which contains no nonempty open set. So no product exists. The argument used only the
point over the origin, so it covers every pointing, including the standard one $(0,a_0)$ with $a_0\ne0$, whose base
point is a smooth point of $\mathcal X$. (A split pointing, with $B_0A_0\ne0$ and residual $\theta_0-sB_0A_0$, has a
hyperbola over $\theta_0$, an embedded submanifold, and there the product exists.)

\emph{(c)} If $h:M\to N$ is an arrow, then $\rho_M=\rho_N\circ h$ gives $\Im M=\rho_N(h(Q_M))\subseteq\Im N$, and both
images contain $\theta_0$; functoriality is automatic because the target is a poset. Conversely, let
$\Im M\subseteq\Im N$. Choose, by the axiom of choice, a preimage $s(\theta)\in\rho_N^{-1}(\theta)$ for each
$\theta\in\Im M$, with $s(\theta_0)=q_0'$, and put $h=s\circ\rho_M$. Then $\rho_N\circ h=\rho_M$ and $h(q_0)=q_0'$, so $h$
is an arrow of $\Set_*/(\Theta,\theta_0)$. For $\rho_M(t)=t$ and $\rho_N(u)=u^3$ both images are $\R$. A map with
$\rho_N\circ h=\id$ satisfies $h(t)^3=t$, so $h(t)=t^{1/3}$, the unique real cube root, whose difference quotient
$t^{-2/3}$ at $0$ is unbounded. So no $C^1$ arrow, and in particular no smooth one, exists.
\end{proof}

% ======================================================================================================================
\subsection{Definition I.5}\label{proof:I.5}

\begin{restate}{Definition}{I.5}{invariants of an object}
For $M=(Q,q_0,\rho)$:
\begin{itemize}
  \item the image $\Im M=\rho(Q)$;
  \item the \textbf{expressive dimension} $d(M)=\dim\Im M=\max_q\rk d\rho_q$;
  \item the \textbf{gauge waste} $g(M)=\lvert M\rvert-d(M)$;
  \item the \textbf{first-order image} $S_1(M)=d\rho_{q_0}(T_{q_0}Q)$, and the \textbf{first-order deficit}
    $d(M)-\dim S_1(M)$;
  \item the \textbf{gauge group} $\{\varphi\in\mathrm{Diff}(Q):\rho\circ\varphi=\rho\}$, which is unpointed, since a
    gauge transformation moves $q_0$ inside its fibre. In practice we exhibit a Lie group acting on $Q$ that preserves
    $\rho$ and is locally transitive on generic fibres, and call it the \textbf{generic gauge}.
\end{itemize}
The base point $\theta_0$ is \textbf{regular} for $M$ if the deficit is $0$, and an \textbf{apex} if the deficit is
positive. We say \textbf{vertex} for the subset-level notion, the cone point of an image germ, and keep \textbf{apex} for
the base-point notion. The two differ in general. At a crossing of two smooth branches the image is singular, yet the
deficit can be $0$. In every case below with a positive deficit, the image germ is (diffeomorphic to) a cone with vertex
$\theta_0$; the exception is DoRA, which inherits LoRA's deficit through its torus factor. A cone germ that is a $C^1$
manifold at its vertex equals its tangent space there, so it is linear. Hence the vertex of a non-linear cone is a
singular point.
\end{restate}

The definition contains three claims: the two descriptions of $d(M)$ agree, the deficit is never negative (so that
regular and apex exhaust the cases), and a cone germ that is $C^1$ at its vertex is linear.

\begin{proof}
\emph{$d(M)$ is well defined.} Let $Q$ be an open definable subset of $\R^N$ and $\rho$ definable and smooth, and let
$k=\max_q\rk d\rho_q$. The set where the rank equals $k$ is open and nonempty, because a nonzero $k\times k$ minor of the
Jacobian stays nonzero nearby (\auxref{B.1}(a)). At a point $p$ of it the rank is constant nearby, so by the rank theorem
(\auxref{B.1}(b)) a small open ball $V\ni p$, a definable set, has image $\rho(V)$ an embedded $k$-dimensional
submanifold. It is definable, of dimension $k$, and contained in $\Im M$, so $\dim\Im M\ge k$ (\auxref{B.9}(b)).
Conversely, let $e=\dim\Im M$. By \auxref{B.9}(b), $\Im M$ contains a definable $C^1$ cell $Y$ of dimension $e$. By
definable choice $\rho$ has a definable section $\sigma:Y\to Q$, and by generic smoothness $\sigma$ is $C^1$ on a nonempty
open subset $Y'$ of $Y$ (\auxref{B.9}(c)). Differentiating $\rho\circ\sigma=\id_{Y'}$ at $y\in Y'$ gives
$d\rho_{\sigma(y)}\circ d\sigma_y=\id_{T_yY}$, so $\rk d\rho_{\sigma(y)}\ge e$, and $k\ge e$. Hence
$\dim\Im M=\max_q\rk d\rho_q$, which is \auxref{B.9}(d). The parameter spaces of the paper are open definable sets, for
instance $\{(u_1,\dots,u_r):u_i\ne0\}$ for HRA, and the maps are built from polynomials, quotients, square roots and, for
softmax, the exponential (\bgref{C.17}).

\emph{The deficit is nonnegative.} $\dim S_1(M)=\rk d\rho_{q_0}\le\max_q\rk d\rho_q=d(M)$. Likewise
$d(M)\le\dim Q=\lvert M\rvert$, so $g(M)\ge0$.

\emph{Cones.} Let $X\subseteq\R^D$ agree near $v$ with a cone $K$ with vertex $v$, meaning $v\in K$ and
$v+t(x-v)\in K$ for all $x\in K$ and $t\ge0$, and suppose $X$ is a $C^1$ embedded submanifold near $v$, with tangent space
$T$ at $v$. Translate so that $v=0$. Near $0$, $X$ is the graph of a $C^1$ map $\varphi$ from a neighbourhood of $0$ in
$T$ to $T^\perp$ with $\varphi(0)=0$ and $d\varphi_0=0$. For $x\in K$ and small $t>0$ the point $tx$ lies in $K$ near $0$,
hence in $X$, so $tP_{T^\perp}x=\varphi(tP_Tx)=o(t)$, and therefore $P_{T^\perp}x=0$. Thus $K\subseteq T$. Near $0$, $X$
is then a graph over $T$ that lies in $T$, so $\varphi=0$ and $X$ contains a neighbourhood $W$ of $0$ in $T$. For $y\in
T$ and small $\varepsilon>0$ we have $\varepsilon y\in W\subseteq K$, so $y=\varepsilon^{-1}(\varepsilon y)\in K$. Hence
$K=T$, and the germ of $X$ at $v$ is the germ of the affine space $v+T$. If $K$ is not an affine subspace, $X$ cannot be a
$C^1$ manifold at $v$, so $v$ is a singular point (compare \auxref{B.6}). Finally, a figure-eight curve through its
crossing with nonzero velocity has deficit $0$ there and a singular image, so vertex and apex are different notions.
\end{proof}

% ======================================================================================================================
\subsection{Observation I.6}\label{proof:I.6}

\begin{restate}{Observation}{I.6}{$S_1$ is functorial}
If $M\to N$ is a (pointed) morphism, then $S_1(M)\subseteq S_1(N)$. In particular, isomorphic methods have equal $S_1$ and
equal first-order deficit. Monotonicity holds only along pointed morphisms. Along a weak morphism, such as a gauge map
moving $q_0$ inside its fibre, $S_1$ can change.
\end{restate}

\begin{proof}
Let $h:M\to N$ be a pointed arrow, so $\rho_N\circ h=\rho_M$ and $h(q_0)=q_0'$. The chain rule at $q_0$ gives
$d(\rho_N)_{q_0'}\circ dh_{q_0}=d(\rho_M)_{q_0}$, so
$S_1(M)=d(\rho_N)_{q_0'}\big(dh_{q_0}(T_{q_0}Q_M)\big)\subseteq d(\rho_N)_{q_0'}(T_{q_0'}Q_N)=S_1(N)$. If $h$ is an
isomorphism, $h^{-1}$ is a pointed arrow $N\to M$, so $S_1(N)\subseteq S_1(M)$ as well. Moreover $\rho_M=\rho_N\circ h$
with $h$ bijective gives $\Im M=\Im N$, so $d(M)=d(N)$ (equivalently, $\rk d(\rho_M)_q=\rk d(\rho_N)_{h(q)}$ since $dh_q$ is
invertible), and the deficits agree.

Without $h(q_0)=q_0'$ the chain rule compares $S_1(M)$ with the first-order image of $N$ at $h(q_0)$, which need not be
the base point of $N$. For instance, take LoRA with frozen residual $\theta_0$ and $1\le r<\min(m,n)$, and let $M$ and
$N$ be its pointings at $(0,A_0)$ and at $(B_0,0)$, with $A_0$ and $B_0$ of full rank; both points lie in the fibre over
$\theta_0$. The identity of $Q$ is a weak morphism $M\to N$, yet $S_1(M)=\{XA_0\}=\mathcal R(\operatorname{row}A_0)$ and
$S_1(N)=\{B_0Y\}=\mathcal C(\operatorname{col}B_0)$ are not nested: $\mathcal R(\operatorname{row}A_0)$ contains $e_iv\T$
for every $i$ and every $v\in\operatorname{row}A_0$, which is not in $\mathcal C(\operatorname{col}B_0)$ when
$e_i\notin\operatorname{col}B_0$, and symmetrically.
\end{proof}

% ======================================================================================================================
\subsection{Observation I.7}\label{proof:I.7}

\begin{restate}{Observation}{I.7}{additive monoidal structure}
$\Theta$ is a vector space, so every object of $\Meth(\Theta,\theta_0)$ can be written $\rho=\theta_0+\delta$ with
$\delta:=\rho-\theta_0$ and $\delta(q_0)=0$. Put
\begin{equation*}
  M\oplus N:=\big(Q_M\times Q_N,\ (q_0,q_0'),\ \theta_0+\delta_M+\delta_N\big).
\end{equation*}
This is a symmetric monoidal structure with the frozen model as unit. Its image is the Minkowski sum
$\theta_0+(\Im\delta_M+\Im\delta_N)$. It differs from the categorical product, whose image, when the product exists, is
the intersection $\Im M\cap\Im N$. ``Additive'' in \thmref{Definition}{I.2} ($\delta$ independent of $\theta_0$) is a
condition on \emph{families} $\theta_0\mapsto M(\theta_0)$; it matters for base change and plays no role here.

Write $\mathrm{LoRA}_k$ for LoRA with frozen residual $\theta_0$; for split pointings the sum below has residual
$\theta_0-P-P'$. Concatenation $h(B_1,A_1,B_2,A_2)=([B_1\ B_2],[A_1;A_2])$ is an isomorphism
$\mathrm{LoRA}_r\oplus\mathrm{LoRA}_s\cong\mathrm{LoRA}_{r+s}$ of unpointed objects. It is a pointed isomorphism onto
$\mathrm{LoRA}_{r+s}$ pointed at $([B_0\ B_0'],[A_0;A_0'])$. For $\mathrm T_A$ pointings ($B_0=0$, $\rk A_0=r$), the
target pointing has $\rk[A_0;A_0']=r+s$ if and only if $\operatorname{row}A_0\cap\operatorname{row}A_0'=0$; it lies in the
stratum $\mathrm T_A$ of \thmref{Theorem}{II.2} when moreover $r+s<\min(m,n)$.

Standard pointings ($\mathrm T_A$, or symmetrically $\mathrm T_B$) exist only for $r<\min(m,n)$, so the monoidal question
is pointwise. If $F(r)$ and $F(2r)$ are standard ($2r<\min(m,n)$), then $F(r)\oplus F(r)\not\cong F(2r)$ as pointed
objects; for $2r\ge\min(m,n)$ no standard $F(2r)$ exists. Hence no assignment $r\mapsto\mathrm{LoRA}_r$ with $F(1)$ and
$F(2)$ standard is \textbf{strong monoidal} into pointed objects. For $r\neq s$ and standard pointings, $F(r)\oplus
F(s)\cong F(r+s)$ as pointed objects exactly when the row space of the target's $A_0$ is the direct sum of the two source
row spaces. The assignment is strong monoidal on unpointed objects, up to weak morphisms, and on pointed objects with
degenerate pointings, namely the origin $(B_0,A_0)=(0,0)$, where $S_1=0$, or the stacked pointings $F(r)=F(1)^{\oplus r}$.
Concatenation is strictly associative and unital, so discrete $(\N,+)$ suffices for a plain strong monoidal functor.
\emph{Symmetric} strong monoidality needs the permutation groupoid, since swapping and then concatenating permutes the
blocks.

With per-adapter scales $s_r$ ($\alpha/r$ for LoRA, $\alpha/\sqrt r$ for rsLoRA) and any target scale $s_*$, the map
$h=([B_1\ B_2],[c_1A_1;c_2A_2])$ with $c_i=s_{r_i}/s_*$, where $r_1=r$ and $r_2=s$, is an isomorphism over $\Theta$. It is
a linear bijection and preserves row spaces, so the pointing criterion is unchanged.
\end{restate}

\begin{proof}
\emph{The monoidal structure.} $M\oplus N$ is an object: $Q_M\times Q_N$ is a manifold, $\theta_0+\delta_M+\delta_N$ is
smooth and takes the value $\theta_0$ at $(q_0,q_0')$. For arrows $h:M\to M'$ and $k:N\to N'$, the pointed map $h\times k$
lies over $\Theta$ because $\delta_{M'}\circ h=\rho_{M'}\circ h-\theta_0=\delta_M$ and likewise for $k$; this makes $\oplus$
a functor. The associator $(Q_1\times Q_2)\times Q_3\cong Q_1\times(Q_2\times Q_3)$, the unitors
$\mathrm{pt}\times Q\cong Q\cong Q\times\mathrm{pt}$ and the symmetry $Q_M\times Q_N\cong Q_N\times Q_M$ are the
structure maps of the cartesian monoidal structure of $\Diff_*$. They lie over $\Theta$ because addition is associative and
commutative and the frozen model has $\delta=0$. Their coherence (pentagon, triangle, hexagon and the symmetry being an
involution) holds in $\Diff_*$, and the forgetful functor $\Meth(\Theta,\theta_0)\to\Diff_*$ is faithful, so it holds in
$\Meth(\Theta,\theta_0)$.

\emph{Image.} $\{\theta_0+\delta_M(q)+\delta_N(q')\}=\theta_0+(\Im\delta_M+\Im\delta_N)$. The categorical product has
image $\Im M\cap\Im N$ when it exists (\thmref{Observation}{I.4}(b)), and the two differ in general: for zero-init
$\mathrm{LoRA}_1$ on a $2\times2$ box, $M\oplus M$ has image $\theta_0+\MM_{\le2}$, while any product of $M$ with itself
has image $\theta_0+\MM_{\le1}$.

\emph{Concatenation.} The map $h$ only regroups coordinates, so it is a linear bijection between the parameter spaces,
and $\theta_0+s[B_1\ B_2][A_1;A_2]=\theta_0+sB_1A_1+sB_2A_2$, so it lies over $\Theta$. It is a pointed isomorphism onto
$\mathrm{LoRA}_{r+s}$ pointed at the image of $(q_0,q_0')$, the concatenated point. For split pointings with
$P=sB_0A_0$ and $P'=sB_0'A_0'$, the sum is $(\theta_0-P-P')+s[B_1\ B_2][A_1;A_2]$, which is $\mathrm{LoRA}_{r+s}$ with
residual $\theta_0-P-P'$. For $\mathrm T_A$ pointings, the concatenated pointing has $B=0$ and
$\rk[A_0;A_0']=\dim(\operatorname{row}A_0+\operatorname{row}A_0')=r+s-\dim(\operatorname{row}A_0\cap\operatorname{row}A_0')$,
which is $r+s$ exactly when the intersection is $0$. The stratum $\mathrm T_A$ of rank $r+s$ is defined for
$r+s<\min(m,n)$, which gives the last condition.

\emph{Row and column types.} We use two elementary facts about the spaces $\mathcal R(V)$ and $\mathcal C(U)$ of the
introduction, which return in \thmref{Theorem}{II.2}. First, $\mathcal R(V)=\{XA:X\in\R^{m\times k}\}$ for any $A$ of full
row rank with $\operatorname{row}A=V$, since $Z=ZA^+A$ when $\operatorname{row}Z\subseteq V$; so $\dim\mathcal
R(V)=m\dim V$, and $\mathcal R(V)$ determines $V$ as the sum of the row spaces of its elements. Likewise
$\mathcal C(U)=\{BY\}$ for $\operatorname{col}B=U$, of dimension $n\dim U$, determines $U$. Second, if $V\ne0$ and
$\mathcal R(V)=\mathcal C(U)$, then $U=\R^m$, because $\mathcal R(V)$ contains $e_iv\T$ for every $i$ and every nonzero
$v\in V$, and $e_iv\T\in\mathcal C(U)$ forces $e_i\in U$.

\emph{The negative claim.} Standard pointings are the strata $\mathrm T_A$ and $\mathrm T_B$ of \thmref{Theorem}{II.2},
which are defined only for ranks below $\min(m,n)$. Let $F(r)$ and $F(2r)$ be standard, so $2r<\min(m,n)$. If $F(r)$ is
pointed at $(0,A_0)$, then $F(r)\oplus F(r)$ is pointed at $((0,A_0),(0,A_0))$ and its first-order image is
$\{X_1A_0+X_2A_0\}=\mathcal R(\operatorname{row}A_0)$, of dimension $mr$. If $F(r)$ is in $\mathrm T_B$, it is
$\mathcal C(\operatorname{col}B_0)$, of dimension $nr$. A standard $F(2r)$ has first-order image $\mathcal R(V'')$ with
$\dim V''=2r$ or $\mathcal C(U'')$ with $\dim U''=2r$. A pointed isomorphism would make the two first-order images equal
(\thmref{Observation}{I.6}). Two row-type spaces determine their subspaces, of dimensions $r\ne2r$; two column-type spaces
likewise; and a row-type space with $V\ne0$ equals a column-type space only if the column subspace is all of $\R^m$,
impossible for $U''$ of dimension $2r<m$ (and symmetrically). So $F(r)\oplus F(r)\not\cong F(2r)$. A strong monoidal
functor $F$ would supply a pointed isomorphism $F(1)\oplus F(1)\cong F(2)$, so none exists when $F(1)$ and $F(2)$ are
standard.

\emph{The case $r\ne s$.} Let $F(r)$, $F(s)$ and $F(r+s)$ be standard of type $\mathrm T_A$, pointed at $(0,A_0)$,
$(0,A_0')$ and $(0,A_0'')$; the type $\mathrm T_B$ is the transpose. Then
$S_1(F(r)\oplus F(s))=\{XA_0+X'A_0'\}=\mathcal R(\operatorname{row}A_0+\operatorname{row}A_0')$ and
$S_1(F(r+s))=\mathcal R(\operatorname{row}A_0'')$, with $\dim\operatorname{row}A_0''=r+s$. A pointed isomorphism forces
$\operatorname{row}A_0''=\operatorname{row}A_0+\operatorname{row}A_0'$, and the dimension count makes the sum direct.
Conversely, if $\operatorname{row}A_0''=\operatorname{row}A_0\oplus\operatorname{row}A_0'$, then $[A_0;A_0']$ has full row
rank and the row space of $A_0''$, so $A_0''=g[A_0;A_0']$ for a unique $g\in\GL_{r+s}$. Concatenation followed by the gauge
$(B,A)\mapsto(Bg^{-1},gA)$, which preserves $BA$, is a pointed isomorphism $F(r)\oplus F(s)\to F(r+s)$.

\emph{Strong monoidality.} The concatenation maps $\phi_{r,s}=h$ satisfy
$\phi_{r+s,t}\circ(\phi_{r,s}\times\id)=\phi_{r,s+t}\circ(\id\times\phi_{s,t})\circ a$, both sides sending
$(B_1,A_1,B_2,A_2,B_3,A_3)$ to $([B_1\ B_2\ B_3],[A_1;A_2;A_3])$, and $\phi_{0,r}$, $\phi_{r,0}$ are the unitors, with
$\mathrm{LoRA}_0$ the frozen model. Discrete $(\N,+)$ has only identity arrows, so naturality is vacuous, and $r\mapsto
\mathrm{LoRA}_r$ is strong monoidal on unpointed objects. As maps between chosen pointings, the $\phi_{r,s}$ are weak
morphisms. They are pointed isomorphisms when every $F(r)$ is pointed at the origin, since $h$ maps origin to origin, and
when $F(r)$ is pointed at the concatenation of $r$ copies of the base point of $F(1)$, since $h$ maps concatenations to
concatenations. The symmetry of discrete $(\N,+)$ is the identity, so symmetry would require $\phi_{r,s}=\phi_{s,r}\circ
\sigma$, where $\sigma$ swaps the factors. The left side sends $(B_1,A_1,B_2,A_2)$ to $([B_1\ B_2],[A_1;A_2])$ and the
right side to $([B_2\ B_1],[A_2;A_1])$, which differ for $r,s\ge1$. On the permutation groupoid, a permutation $\pi\in S_r$
acts on $\mathrm{LoRA}_r$ by the gauge $(B,A)\mapsto(BP_\pi^{-1},P_\pi A)$, which is functorial and lies over $\Theta$, the
maps $\phi_{r,s}$ are natural for block sums of permutations, and with the block transposition $\tau_{r,s}\in S_{r+s}$ as
symmetry, $F(\tau_{r,s})\circ\phi_{r,s}=\phi_{s,r}\circ\sigma$, because permuting the channel blocks of
$([B_1\ B_2],[A_1;A_2])$ gives $([B_2\ B_1],[A_2;A_1])$. The origin and the stacked pointings are fixed by channel
permutations, so the same holds for them.

\emph{Rescaled concatenation.} $s_*[B_1\ B_2][c_1A_1;c_2A_2]=s_*c_1B_1A_1+s_*c_2B_2A_2=s_rB_1A_1+s_sB_2A_2$. The scales
are nonzero, so $h$ is a linear bijection, and $\operatorname{row}(c_iA_i)=\operatorname{row}A_i$, so the criterion for
the concatenated pointing is the same as before.
\end{proof}

% ======================================================================================================================
\subsection{Proposition I.8}\label{proof:I.8}

\begin{restate}{Proposition}{I.8}{a product that produces a named method; an image-level meet}
(a) Let $A_0\in\R^{r\times n}$ have full row rank and $B_0\in\R^{m\times r}$ full column rank. Define the one-sided methods
\begin{equation*}
  \mathrm{FA}(A_0):B\mapsto W_0+BA_0,\qquad \mathrm{FB}(B_0):A\mapsto W_0+B_0A,
\end{equation*}
both pointed at $0$. Then
\begin{equation*}
  \mathrm{FA}(A_0)\times\mathrm{FB}(B_0)\;\cong\;\big(\R^{r\times r},\ 0,\ R\mapsto W_0+B_0RA_0\big).
\end{equation*}
The product supplies this sandwich for \emph{any} frames. Assume $\rk W_0\ge r$ ($\sigma_r>0$) and take the SVD frames
$B_0=U_r$, $A_0=\Sigma_rV_r\T$. The result is \textbf{LoRA-XS pointed at $R=0$}.

(b) Let $W_0$ have full row rank (so $m\le n$), and let $\mathrm{LoRA}_r$ carry its standard pointing ($B_0=0$, image
$W_0+\Mr$). The image-level meet of $\mathrm{LoRA}_r$ with the output-scaling method $\ell\mapsto\diag(\ell)W_0$ is
$\{\diag(\ell)W_0:\#\{i:\ell_i\neq1\}\le r\}$, that is, ``rescale at most $r$ output neurons''.

\upshape The LoRA-XS paper writes $xW+xARB$ in row-vector convention with $A=U_r\Sigma_r$ and $B=V_r\T$. Transposed, the
update is $B_0R\T A_0$, so the isomorphism is $R\mapsto\alpha R\T$. LoRA-XS initialises $R\sim\mathcal N(0,\sigma^2)$
with $\sigma=10^{-5}$, a negligible pointing defect. The output-scaling method in (b) is the key/value part of (IA)$^3$.
On the FFN, (IA)$^3$ scales the \emph{input} of $W_{\rm down}$, $W\mapsto W\diag(\ell)$; the mirror statement then needs
full column rank, or holds generically for $r<d$, where $d$ is the number of rows of $W_{\rm down}$.
\end{restate}

\begin{proof}
\emph{(a)} Write $B_0^+=(B_0\T B_0)^{-1}B_0\T$ and $A_0^+=A_0\T(A_0A_0\T)^{-1}$, so that $B_0^+B_0=I_r=A_0A_0^+$. A point of
the set-level fibre product is a pair $(B,A)$ with $BA_0=B_0A=:D$. Then $\operatorname{row}D\subseteq\operatorname{row}A_0$
and $\operatorname{col}D\subseteq\operatorname{col}B_0$, so $D=B_0B_0^+DA_0^+A_0=B_0RA_0$ with $R=B_0^+DA_0^+$, and $R$ is
unique because $B_0$ is injective and $A_0$ surjective. From $BA_0=B_0RA_0$ and $A_0A_0^+=I$ we get $B=B_0R$, and from
$B_0A=B_0RA_0$ we get $A=RA_0$. So the fibre product is the $r^2$-dimensional linear subspace $\{(B_0R,RA_0)\}$ of
$\R^{m\times r}\times\R^{r\times n}$, and $j:R\mapsto(B_0R,RA_0)$ is a linear isomorphism onto it, with inverse
$(B,A)\mapsto B_0^+B$. Both composites $R\mapsto W_0+B_0RA_0$ agree, and $j(0)=(0,0)$ is the pair of base points.

The universal property: let $T$ be a pointed manifold and $h_1:T\to\R^{m\times r}$, $h_2:T\to\R^{r\times n}$ pointed smooth
maps with $W_0+h_1(t)A_0=W_0+B_0h_2(t)$. By the previous paragraph $(h_1(t),h_2(t))=j(R(t))$ with $R(t)=B_0^+h_1(t)$, a
pointed smooth map, and it is the only map with this property because $j$ is injective. So $(\R^{r\times r},0,R\mapsto
W_0+B_0RA_0)$ with the projections $R\mapsto B_0R$ and $R\mapsto RA_0$ is the product in $\Meth(\Theta,W_0)$
(\auxref{B.2}(c), \auxref{B.3}). This holds for any frames of full rank. With $B_0=U_r$ and $A_0=\Sigma_rV_r\T$ from an SVD
of $W_0$ (\auxref{B.11}(a)), the frames have full rank exactly when $\sigma_r>0$; if $\rk W_0<r$, $\Sigma_rV_r\T$ is
rank-deficient and the hypothesis fails.

The row-vector convention: LoRA-XS acts on row vectors by $x\mapsto x(W_{\rm row}+\alpha ARB)$, where $W_{\rm row}=W_0\T$
has SVD $U'\Sigma V'^{\top}$, $A=U'_r\Sigma_r$ and $B=V_r'^{\top}$. In our convention the update is the transpose,
$\alpha V'_rR\T\Sigma_rU_r'^{\top}$. Since $W_0=V'\Sigma U'^{\top}$, the left and right singular frames of $W_0$ are
$U_r=V'_r$ and $V_r=U'_r$, so the update is $U_r(\alpha R\T)\Sigma_rV_r\T=B_0(\alpha R\T)A_0$. Hence LoRA-XS is the product
composed with the linear isomorphism $R\mapsto\alpha R\T$, which fixes $0$. A random initial $R$ makes
$\rho(R_0)-W_0=\alpha U_rR_0\T\Sigma_rV_r\T$ nonzero, of the size of $\sigma$, a pointing defect.

\emph{(b)} The image of $\mathrm{LoRA}_r$ with its standard pointing is $W_0+\Mr$ (\thmref{Theorem}{II.2}(c)), and the
output-scaling method has image $\{\diag(\ell)W_0\}$. A point $\diag(\ell)W_0$ lies in $W_0+\Mr$ iff
$\rk\big(\diag(\ell-\mathbf 1)W_0\big)\le r$. The rows of $\diag(\ell-\mathbf 1)W_0$ are $(\ell_i-1)w_i$, and the nonzero ones
are linearly independent because the rows of $W_0$ are, so the rank is $\#\{i:\ell_i\ne1\}$. The meet in the poset of
subsets is the intersection, which is the stated set.

For the input side, $W_0\diag(\ell)-W_0=W_0\diag(\ell-\mathbf 1)$ consists of the columns $(\ell_j-1)c_j$ of $W_0$, so its
rank is the rank of the columns of $W_0$ indexed by $J=\{j:\ell_j\ne1\}$. If $W_0$ has full column rank this is $\lvert
J\rvert$, and the meet is $\{W_0\diag(\ell):\lvert J\rvert\le r\}$. For $W_{\rm down}\in\R^{d\times d_{\rm ff}}$ with
$d<d_{\rm ff}$, full column rank is impossible. For generic $W_0$ every set of at most $d$ columns is independent (the
failure is the vanishing of all maximal minors of the chosen columns), so the rank is $\min(\lvert J\rvert,d)$, and when
$r<d$ the condition $\min(\lvert J\rvert,d)\le r$ is again $\lvert J\rvert\le r$.

\emph{No product in (b).} The commentary after the proposition says that, for $n\ge2$, $\mathrm{LoRA}_r$ and the
output-scaling method have no product; the argument is that of \thmref{Observation}{I.4}(b). Their set-level fibre product
is $\mathcal F=\{(B,A,\ell):W_0+sBA=\diag(\ell)W_0\}$, whatever the pointing of $\mathrm{LoRA}_r$ (with residual $W_0$).
Suppose a product $P$ exists. By the proof of \thmref{Observation}{I.4}(b) its canonical map $\iota:P\to\R^{m\times
r}\times\R^{r\times n}\times\R^m$ is a bijection onto $\mathcal F$, $P$ is also the pullback in $\Diff$, and smooth curves
in $\mathcal F$ lift. Let $p^*$ be the point over $(0,0,\mathbf 1)$. For every $B'$ and $A'$ the curves $t\mapsto(tB',0,\mathbf
1)$ and $t\mapsto(0,tA',\mathbf 1)$ lie in $\mathcal F$ and lift through $p^*$, so with $\mathrm{pr}$ the projection to the
LoRA parameters, $d(\mathrm{pr}\circ\iota)_{p^*}$ is onto $\R^{m\times r}\times\R^{r\times n}$. By \auxref{B.1}(a),(c),
$\mathrm{pr}\circ\iota$ is a submersion near $p^*$ and its image contains a nonempty open set of pairs $(B,A)$. But every pair
in $\mathrm{pr}(\mathcal F)$ has $BA=\diag(\ell-\mathbf 1)W_0/s$, so the first row of $BA$ is parallel to the first row
$w_{0,1}$ of $W_0$: the $2\times2$ minors of the $2\times n$ matrix with rows $e_1\T BA$ and $w_{0,1}\T$ vanish. For
$n\ge2$ one of these minors is a polynomial in $(B,A)$ that is not identically zero, so its zero set has empty interior
(\auxref{B.10}(a)), a contradiction. (For $m=n=1$ the fibre product is the graph of $\ell=1+sBA/W_0$, and the product
exists.)
\end{proof}

% ======================================================================================================================
\subsection{Theorem II.1}\label{proof:II.1}

\begin{restate}{Theorem}{II.1}{pointings of an additive method}
Let $M=(Q,q_0,\theta_0+\delta)$ be additive with shape $C=\delta(Q)$. Write $-v+C$ for a translate, and
$C-C:=\{c-c':c,c'\in C\}$ for the Minkowski difference.

(a) For each $q\in Q$, the re-pointed method $M_q:=(Q,q,\theta_0-\delta(q)+\delta(\cdot))$, with frozen residual
$\theta_0-\delta(q)$, is pointed at $\theta_0$, and $\Im M_q=\theta_0+(-\delta(q)+C)$.

(b) $\bigcup_q\Im M_q=\theta_0+(C-C)$.

(c) For LoRA$_r$ on a single linear box with $s\ne0$, $C-C=\MM_{\le\min(2r,m,n)}$; for several boxes, $C-C$ is the
product of the per-box varieties $\MM_{\le\min(2r,m_i,n_i)}$. Hence \textbf{no exact pointing of LoRA$_r$ (frozen
residual $\theta_0-\delta(q)$), whether zero-init or split (PiSSA, OLoRA, CorDA, LoRA-GA), makes an update
$\theta-\theta_0$ of rank $>2r$ reachable in any adapted matrix.} More robustly, from \emph{any} starting point $q_0$ of
LoRA$_r$, pointed or not, $\rho(q)-\rho(q_0)=\delta(q)-\delta(q_0)$ has rank $\le\min(2r,m,n)$. Quantised-residual
schemes (QLoRA, LoftQ, QPiSSA) are pointed only up to a defect that depends on the order of operations ($e$ for QLoRA;
\thmref{Proposition}{V.2}(b) for the others). Measured from $\theta_0$ rather than from their starting weight, their
updates contain this defect, whose rank is generically at least $\min(m,n)-r$. The bound $2r$ is attained by split
pointings when $2r\le\min(m,n)$.
\end{restate}

\begin{proof}
\emph{(a)} $M_q$ has parameter space $Q$ and the smooth map $\rho_q=\theta_0-\delta(q)+\delta(\cdot)$, with
$\rho_q(q)=\theta_0$, so it is pointed at $\theta_0$, and $\Im M_q=\theta_0-\delta(q)+\delta(Q)=\theta_0+(-\delta(q)+C)$.

\emph{(b)} As $q$ ranges over $Q$, $\delta(q)$ ranges over $C$, so
$\bigcup_q\Im M_q=\theta_0+\bigcup_{c'\in C}(C-c')=\theta_0+(C-C)$.

\emph{(c)} On one box $C=\{sBA\}$. Since $s\ne0$ and every matrix of rank at most $\min(r,m,n)$ factors as $BA$ with inner
dimension $r$ (\auxref{B.11}(a)), $C=\MM_{\le\min(r,m,n)}$, which we write $\Mr$. For $X,Y\in C$,
$\rk(X-Y)\le\rk X+\rk Y\le2r$ and $\rk(X-Y)\le\min(m,n)$ (\auxref{B.12}(a)). Conversely, let $Z$ have rank
$k\le\min(2r,m,n)$ and SVD $Z=\sum_{i\le k}\sigma_iu_iv_i\T$ (\auxref{B.11}(a)). Put $Z_1=\sum_{i\le\min(k,r)}\sigma_iu_iv_i\T$
and $Z_2=Z-Z_1$, both of rank at most $r$. The cone $C$ is closed under negation, since $-BA=(-B)A$, so
$Z=Z_1-(-Z_2)\in C-C$. On several boxes, each with its own factors, $\delta$ is the product of the per-box maps, so $C$ is
the product of the per-box shapes and $C-C$ the product of the per-box differences.

An exact pointing of $\mathrm{LoRA}_r$ is a re-pointed method $M_q$ of (a), so its image lies in $\theta_0+(C-C)$, and in
every box the update $\theta-\theta_0$ has rank at most $\min(2r,m_b,n_b)$. For an arbitrary starting point $q_0$ and an
arbitrary residual $W_{\rm res}$, $\rho(q)-\rho(q_0)=\delta(q)-\delta(q_0)\in C-C$ with no condition at all.

For a pipeline with pointing defect $e_{\rm pt}=\rho(q_0)-\theta_0$, the update measured from $\theta_0$ is
\begin{equation*}
  \theta-\theta_0=\big(\rho(q)-\rho(q_0)\big)+e_{\rm pt},
\end{equation*}
a matrix of rank at most $2r$ plus the defect. The defects are those of \thmref{Proposition}{V.2}: $e=\kappa\theta_0-
\theta_0$ for QLoRA, $e-s\,\mathrm{SVD}_r(e)$ for one-step LoftQ, and $\kappa(W_{\rm res})-W_{\rm res}$ for QPiSSA. Each is
a quantisation error, or a quantisation error minus a matrix of rank at most $r$, so it has rank at least $\min(m,n)-r$
whenever the quantisation error involved has full rank (\auxref{B.12}(a)). Full rank is the generic case recorded in
\thmref{Proposition}{V.2}(a). For instance, a round-to-nearest quantiser with fixed scales is constant on each open cell of
constant codes, so there $e(W)=c-W$, which is rank-deficient only on a translate of $\MM_{\le\min(m,n)-1}$, a proper
algebraic subset of the cell; the cells cover all weights outside a null set (\auxref{B.10}(b)).

Attainment: let $2r\le\min(m,n)$, and take a split pointing with $P=sB_0A_0$ of rank $r$. Choose $Y=sBA$ of rank $r$ with
$\operatorname{col}Y\cap\operatorname{col}P=0$ and $\operatorname{row}Y\cap\operatorname{row}P=0$, possible since
$2r\le\min(m,n)$. Writing $Y=U_YV_Y\T$ and $P=U_PV_P\T$ with factors of full column rank $r$, we have
$Y-P=[U_Y\ U_P]\,[V_Y\ {-V_P}]\T$, and both $m\times2r$ and $n\times2r$ factors have independent columns, so
$\rk(Y-P)=2r$. The weight $\theta_0-P+Y$ lies in the image of the split pointing, with update $Y-P$ of rank $2r$.
\end{proof}

% ======================================================================================================================
\subsection{Theorem II.2}\label{proof:II.2}

\begin{restate}{Theorem}{II.2}{three strata of LoRA initialisations}
Let $\delta(B,A)=BA$ with $1\le r<\min(m,n)$. Consider pointed LoRA objects in three distinguished strata of full-rank
pointings (pointings whose nonzero factors have full rank). Only $\mathrm S$ is open and dense in the space of pointings;
$\mathrm T_A$ and $\mathrm T_B$ are nowhere dense.
\begin{itemize}
  \item $(\mathrm T_A)$: $B_0=0$ and $\rk A_0=r$;
  \item $(\mathrm T_B)$: $A_0=0$ and $\rk B_0=r$;
  \item $(\mathrm S)$: $\rk B_0A_0=r$, with frozen residual $\theta_0-B_0A_0$.
\end{itemize}

(a) \textbf{Classification.} Within a stratum, $M_{q_0}\cong M_{q_0'}$ if and only if, respectively,
$\operatorname{row}A_0=\operatorname{row}A_0'$; $\operatorname{col}B_0=\operatorname{col}B_0'$; or $B_0A_0=B_0'A_0'$.
Objects in different strata are never isomorphic. As a \emph{set}, the isomorphism classes of the three strata therefore
form $\Gr(r,\R^n)\sqcup\Gr(r,\R^m)\sqcup\MM_r(m\times n)$.

(b) \textbf{First order.} $\dim S_1$ is $mr$, $nr$ and $r(m+n-r)$ respectively, while $d=r(m+n-r)$ in every case. The
deficits are $r(n-r)$, $r(m-r)$ and $0$.

(c) \textbf{Geometry.} In $\mathrm T_A$ and $\mathrm T_B$, $\Im=\theta_0+\Mr$, whose vertex $\theta_0$ is an
\textbf{apex}. In $\mathrm S$, $\Im=(\theta_0-P)+\Mr$ with $P=B_0A_0$, and $\theta_0$ is a \textbf{smooth point}.

(d) \textbf{Incomparability.} The images of $\mathrm S$-type and $\mathrm T$-type objects are incomparable.

(e) \textbf{Scope.} The three strata do not exhaust the pointings of LoRA$_r$. Outside them lie: $B_0=0$ with
$\rk A_0<r$; $A_0=0$ with $\rk B_0<r$; $B_0=A_0=0$; $0<\rk B_0A_0<r$; and pointings with $B_0,A_0$ both nonzero and
$B_0A_0=0$. HF PEFT ships one of the last kind: \texttt{init\_lora\_\allowbreak weights=\allowbreak"orthogonal"} (which
requires $r$ even) lies in the stratum $(\operatorname{rk}B_0,\operatorname{rk}A_0,\operatorname{rk}B_0A_0)=(r/2,r/2,0)$,
an apex with $\dim S_1=\tfrac r2(m+n-\tfrac r2)$. It is isomorphic to no object of the three strata. Its image
$\theta_0+\Mr$ rules out $\mathrm S$, and its $S_1$ determines $(\operatorname{col}B_0,\operatorname{row}A_0)$, of
dimensions $(r/2,r/2)$, against $(0,r)$ in $\mathrm T_A$ and $(r,0)$ in $\mathrm T_B$. The dimension of $S_1$ alone does
not separate it, since at $(m,n,r)=(4,5,2)$ it equals $mr$. This placement holds in exact arithmetic; HF casts the two
factors to the base dtype separately, and in bf16 the result is a generic pointing with a tiny defect. For every
pointing, the image determines $P=B_0A_0$ (its vertex is $\theta_0-P$), and $S_1=\{XA_0+B_0Y\}$ determines
$(\operatorname{row}A_0,\operatorname{col}B_0)$ whenever these are proper subspaces. The classification is set-level: in
the space of pointings $\mathrm T_A$ lies in the closure of $\mathrm S$, since $(\varepsilon B',A_0)\to(0,A_0)$, so the
quotient is not a topological disjoint union.
\end{restate}

\begin{proof}
Absorb the scale $s$ into $B$; this changes neither the strata nor the images. The space of pointings is
$\R^{m\times r}\times\R^{r\times n}$, and the residual is determined by the pointing.

\emph{Openness.} $\rk B_0A_0\le\min(\rk B_0,\rk A_0)$, and if both factors have rank $r$ then $B_0$ is injective and
$A_0$ surjective, so $\rk B_0A_0=r$. Hence $\mathrm S=\{\rk B_0=\rk A_0=r\}$ is the set where some $r\times r$ minor of
$B_0A_0$ is nonzero: it is open, and dense because its complement is a proper algebraic subset (\auxref{B.10}(a)). The
stratum $\mathrm T_A$ lies in the proper linear subspace $\{B_0=0\}$, which is closed with empty interior, so
$\mathrm T_A$ is nowhere dense, and likewise $\mathrm T_B$.

\emph{A translation argument.} We use the following fact here and in Theorem~II.11 and Proposition~II.16. Let
$1\le r<\min(m,n)$ and $D\ne0$, of rank $k$. If $k\ge r+1$ put $X=0$. Otherwise choose $X$ of rank $r+1-k\le r$ with
$\operatorname{col}X\cap\operatorname{col}D=0$ and $\operatorname{row}X\cap\operatorname{row}D=0$, possible since
$r+1\le\min(m,n)$. Factoring $D=U_DV_D\T$ and $X=U_XV_X\T$ with factors of full column rank, $D+X=[U_D\ U_X][V_D\ V_X]\T$
with both factors of full column rank, so $\rk(D+X)=r+1$. In either case $X\in\Mr$ and $D+X\notin\Mr$. Consequently,
\begin{equation}\label{eq:translation-stabiliser}
  c+\Mr=c'+\Mr\quad\Longrightarrow\quad c=c',
\end{equation}
since otherwise, with $D=c-c'$, the point $c+X=c'+(D+X)$ lies in the left side and not in the right side. The
translation stabiliser of $\Mr$ is $\{0\}$.

\emph{(a), ``if''.} In $\mathrm T_A$, let $\operatorname{row}A_0=\operatorname{row}A_0'$ with both of rank $r$. Then
$A_0'=A_0'A_0^+A_0=gA_0$ with $g=A_0'A_0^+\in\R^{r\times r}$, invertible because $\rk A_0'=r$. The map
$\varphi_g(B,A)=(Bg^{-1},gA)$ is a diffeomorphism of the parameter space with $\rho\circ\varphi_g=\rho$ (both residuals
are $\theta_0$) and $\varphi_g(0,A_0)=(0,A_0')$, a pointed isomorphism over $\Theta$. In $\mathrm T_B$,
$\operatorname{col}B_0=\operatorname{col}B_0'$ gives $B_0'=B_0k$ with $k\in\GL_r$, and $\varphi_{k^{-1}}(B_0,0)=(B_0',0)$.
In $\mathrm S$, let $B_0A_0=B_0'A_0'=P$ of rank $r$. Then
$\operatorname{col}B_0=\operatorname{col}P=\operatorname{col}B_0'$, so $B_0'=B_0k$ for a unique $k\in\GL_r$, and
$B_0kA_0'=B_0A_0$ gives $kA_0'=A_0$ because $B_0$ is injective. So $(B_0',A_0')=\varphi_{k^{-1}}(B_0,A_0)$, and the
residuals $\theta_0-P$ agree.

\emph{(a), ``only if''.} In $\mathrm T_A$, $d\rho_{(0,A_0)}(\dot B,\dot A)=\dot BA_0$, so
$S_1=\{XA_0\}=\mathcal R(\operatorname{row}A_0)$, which determines $\operatorname{row}A_0$ (proof of
\thmref{Observation}{I.7}); pointed isomorphisms preserve $S_1$ (\thmref{Observation}{I.6}). In $\mathrm T_B$,
$S_1=\mathcal C(\operatorname{col}B_0)$ determines $\operatorname{col}B_0$. In $\mathrm S$ the image is
$(\theta_0-P)+\Mr$ by (c), it is an isomorphism invariant, and by \eqref{eq:translation-stabiliser} it determines $P$.

\emph{Different strata.} $S_1(\mathrm T_A)=\mathcal R(V)$ with $\dim V=r\ge1$ and $S_1(\mathrm T_B)=\mathcal C(U)$ with
$\dim U=r<m$. They differ by the second fact in the proof of \thmref{Observation}{I.7}. (When $m\ne n$ their dimensions $mr$
and $nr$ already differ.) $\mathrm S$ differs from both because $\dim S_1(\mathrm S)=r(m+n-r)$ by (b), which exceeds $mr$
by $r(n-r)>0$ and $nr$ by $r(m-r)>0$. Isomorphism classes in $\mathrm T_A$ correspond to $\operatorname{row}A_0$, and every
$r$-dimensional subspace occurs; in $\mathrm T_B$ to $\operatorname{col}B_0$; in $\mathrm S$ to $P$, and every rank-$r$
matrix has a full-rank factorisation. This gives the stated disjoint union of sets.

\emph{(b)} At $(B_0,A_0)$, $S_1=\{XA_0+B_0Y\}$, of dimension $m\rk A_0+n\rk B_0-\rk A_0\rk B_0$ by
\thmref{Proposition}{II.3}: $mr$ in $\mathrm T_A$, $nr$ in $\mathrm T_B$ and $r(m+n-r)$ in $\mathrm S$. In every case the
image is a translate of $\Mr$ by (c), of dimension $r(m+n-r)$ (\auxref{B.5}(a)), so $d=r(m+n-r)$, and the deficits are
$r(n-r)$, $r(m-r)$ and $0$.

\emph{(c)} In $\mathrm T_A$ and $\mathrm T_B$ the residual is $\theta_0$ and the image is $\theta_0+\{BA\}=\theta_0+\Mr$.
This is a cone with vertex $\theta_0$, and it is not an affine subspace, because $\sum_{i\le r}E_{ii}$ and
$E_{r+1,r+1}$ lie in $\Mr$ while their sum has rank $r+1\le\min(m,n)$. By the cone part of \thmref{Definition}{I.5}
(\auxref{B.6}), $\theta_0$ is a singular point, and by (b) the deficit is positive, so $\theta_0$ is an apex. In
$\mathrm S$ the image is $(\theta_0-P)+\Mr$, and $\theta_0=(\theta_0-P)+P$ with $P\in\MM_r$. The stratum $\MM_r$ is open in
$\Mr$ and an embedded submanifold of dimension $r(m+n-r)$ (\auxref{B.5}(b)), so the image is a $C^1$ submanifold near
$\theta_0$, which is a smooth point; its deficit is $0$.

\emph{(d)} The $\mathrm T$-image is $\theta_0+\Mr$ and the $\mathrm S$-image $(\theta_0-P)+\Mr$ with $P\ne0$. The
translation argument with $D=P$ gives $X\in\Mr$ with $\theta_0+X$ outside the $\mathrm S$-image, and with $D=-P$ it gives
$X'\in\Mr$ with $\theta_0-P+X'$ outside the $\mathrm T$-image. (For a generic rank-one $X$, $\rk(P+X)=\rk(X-P)=r+1$
already.)

\emph{(e)} Every pointing has a rank triple $(\rk B_0,\rk A_0,\rk B_0A_0)$. The strata are $(0,r,0)$, $(r,0,0)$ and
$(r,r,r)$. If $B_0=0$ and $A_0\ne0$ is not of rank $r$, or $A_0=0$ and $B_0\ne0$ is not of rank $r$, or $B_0=A_0=0$, the
pointing is in the first three listed families. Otherwise both factors are nonzero, and either $\rk B_0A_0=r$, which
forces both ranks to be $r$ and gives $\mathrm S$, or $0<\rk B_0A_0<r$, or $B_0A_0=0$. So the list is exhaustive.

For every pointing, $S_1=\{XA_0+B_0Y\}=\{Z:Z(\operatorname{row}A_0)^\perp\subseteq\operatorname{col}B_0\}$. Indeed,
$(XA_0+B_0Y)w=B_0Yw$ for $w\in\ker A_0=(\operatorname{row}A_0)^\perp$; conversely, if
$Z\ker A_0\subseteq\operatorname{col}B_0$, then $Z=ZA_0^+A_0+B_0B_0^+Z(I-A_0^+A_0)$. Its annihilator in the Frobenius inner
product is $\{W:WA_0\T=0,\ B_0\T W=0\}$, the matrices with column space in $(\operatorname{col}B_0)^\perp$ and row space
in $(\operatorname{row}A_0)^\perp$, that is $(\operatorname{col}B_0)^\perp\otimes(\operatorname{row}A_0)^\perp$. When both
subspaces are proper, both complements are nonzero, and the column spaces of the elements of the annihilator span
$(\operatorname{col}B_0)^\perp$ while their row spaces span $(\operatorname{row}A_0)^\perp$. So $S_1$ determines
$(\operatorname{col}B_0,\operatorname{row}A_0)$ and their dimensions, and so does every pointed isomorphism class
(\thmref{Observation}{I.6}). The image is $(\theta_0-P)+\Mr$ with $P=B_0A_0$ (residual $\theta_0-P$), and it determines
$P$ by \eqref{eq:translation-stabiliser}.

HF's \texttt{orthogonal} initialisation draws one $Q\in\Ort(r)$ and sets $A_0=(G_AQ_{[0::2]})\T/10$ and
$B_0=G_BQ_{[1::2]}/10$, with $G_A\in\R^{n\times r/2}$ and $G_B\in\R^{m\times r/2}$ random and $Q_{[0::2]}$, $Q_{[1::2]}$
the even- and odd-indexed rows of $Q$. Distinct rows of $Q$ are orthogonal, so $B_0A_0=G_BQ_{[1::2]}Q_{[0::2]}\T
G_A\T/100=0$, and both factors have rank $r/2$ for generic $G_A,G_B$. The residual is $\theta_0$ and the image is
$\theta_0+\Mr$, whose vertex is singular. By \thmref{Proposition}{II.3}, $\dim S_1=m\tfrac r2+n\tfrac r2-\tfrac{r^2}4=
\tfrac r2(m+n-\tfrac r2)$, and the deficit $r(m+n-r)-\tfrac r2(m+n-\tfrac r2)=\tfrac r2(m+n-\tfrac{3r}2)$ is positive, so
the base point is an apex. Its image differs from every $\mathrm S$-image by \eqref{eq:translation-stabiliser}, and its
pair of subspace dimensions $(r/2,r/2)$, determined by $S_1$ since both subspaces are proper, differs from $(0,r)$ and
$(r,0)$; so it is isomorphic to no object of the three strata. At $(m,n,r)=(4,5,2)$, $\tfrac r2(m+n-\tfrac r2)=8=mr$.

Finally, for $B'$ of rank $r$ and $\varepsilon\ne0$, $(\varepsilon B',A_0)\in\mathrm S$, and $(\varepsilon B',A_0)\to
(0,A_0)\in\mathrm T_A$ as $\varepsilon\to0$, so $\mathrm T_A$ lies in the closure of $\mathrm S$.
\end{proof}

% ======================================================================================================================
\subsection{Proposition II.3}\label{proof:II.3}

\begin{restate}{Proposition}{II.3}{first-order rank formula: smoothness, alignment and scale}
For $\kappa(B,A)=W_{\rm base}+sBA$ with $s\ne0$,
\begin{equation*}
  \rk d\kappa_{(B,A)}=m\operatorname{rk}A+n\operatorname{rk}B-\operatorname{rk}A\operatorname{rk}B\qquad\text{for all }(B,A),
\end{equation*}
and $\{XA\}\cap\{BY\}=\{BRA:R\in\R^{r\times r}\}$. For full-rank frames this overlap is the tangent space of the product
$\mathrm{FA}(A)\times\mathrm{FB}(B)$ of \thmref{Proposition}{I.8}(a), re-based at $\kappa(B,A)$, whose image is the meet
of the two one-sided images: a generalised LoRA-XS. It is LoRA-XS proper when $(B,A)$ span the top-$r$ singular subspaces
of $W_0$ ($\sigma_r>0$), as at a PiSSA init; bottom frames (MiLoRA) give a generalised LoRA-XS.

Under gradient flow $\dot B=-\eta_B\nabla_BL$, $\dot A=-\eta_A\nabla_AL$, with $G=\nabla_WL$,
\begin{equation*}
  \tfrac{d}{dt}L=-s^2\big(\eta_B\|GA\T\|_F^2+\eta_A\|B\T G\|_F^2\big)
\end{equation*}
at every time, in particular at $(B_0,A_0)$. To first order in the step sizes the same holds for full-batch gradient
descent, and in expectation for unbiased minibatch SGD. Under Adam ($\varepsilon\to0$, with bias correction) the first
step changes the loss by
\begin{equation*}
  -\lvert s\rvert\big(\eta_B\|GA_0\T\|_1+\eta_A\|B_0\T G\|_1\big)+O(\eta^2),
\end{equation*}
with entrywise $\ell_1$ norms. This first-order term is linear in $\lvert s\rvert$, against quadratic under gradient flow.
\end{restate}

\begin{proof}
\emph{Rank.} $d\kappa_{(B,A)}(\dot B,\dot A)=s(\dot BA+B\dot A)$, and $s\ne0$, so the image of the differential is
$\{XA\}+\{BY\}=\mathcal R(\operatorname{row}A)+\mathcal C(\operatorname{col}B)$. With $a=\rk A$ and $b=\rk B$, the two
summands have dimensions $ma$ and $nb$ (proof of \thmref{Observation}{I.7}). A matrix $D$ lies in both iff
$\operatorname{row}D\subseteq\operatorname{row}A$ and $\operatorname{col}D\subseteq\operatorname{col}B$. The space of such
$D$ is $\operatorname{col}B\otimes\operatorname{row}A$, of dimension $ab$, and every such $D$ satisfies
$D=BB^+DA^+A=B(B^+DA^+)A$, since $BB^+$ and $A^+A$ are the projectors onto $\operatorname{col}B$ and
$\operatorname{row}A$. Conversely $BRA$ lies in both summands. So $\{XA\}\cap\{BY\}=\{BRA\}$ and
$\rk d\kappa=ma+nb-ab$, for every $(B,A)$.

\emph{The overlap.} Let $B$ have full column rank and $A$ full row rank. Re-based at $W=\kappa(B,A)$, the one-sided methods
$X\mapsto W+XA$ and $Y\mapsto W+sBY$ are those of \thmref{Proposition}{I.8}(a) with frames $A$ and $sB$, so their product is
$R\mapsto W+sBRA$, a linear method whose image $W+\{BRA\}$ is its own tangent space at $W$. That image is the meet
$(W+\{XA\})\cap(W+\{BY\})$ by the identity just proved. If $\operatorname{col}B=\operatorname{col}U_r$ and
$\operatorname{row}A=\operatorname{row}V_r\T$ for the top-$r$ singular frames of $W_0$ with $\sigma_r>0$, then $B=U_rk_1$ and
$A=k_2V_r\T$ with $k_1,k_2\in\GL_r$, and $\{BRA\}=\{U_rR'\Sigma_rV_r\T\}$ with $R'=k_1Rk_2\Sigma_r^{-1}$ ranging over
$\R^{r\times r}$, which is LoRA-XS's update space (\thmref{Proposition}{I.8}(a)). For other frames, such as the bottom
singular frames of MiLoRA, the overlap is the same kind of sandwich in those frames.

\emph{Gradient flow.} With $W=\kappa(B,A)$, the chain rule gives $\nabla_B(L\circ\kappa)=sGA\T$ and
$\nabla_A(L\circ\kappa)=sB\T G$. Along the flow,
\begin{equation*}
  \tfrac{d}{dt}L=\langle G,s(\dot BA+B\dot A)\rangle=s\langle GA\T,\dot B\rangle+s\langle B\T G,\dot A\rangle
  =-s^2\big(\eta_B\|GA\T\|_F^2+\eta_A\|B\T G\|_F^2\big).
\end{equation*}
One step of full-batch gradient descent moves $W$ by $s(\Delta B\,A+B\,\Delta A)+s\,\Delta B\,\Delta A$ with
$\Delta B=-\eta_BsGA\T$ and $\Delta A=-\eta_AsB\T G$. For $L$ of class $C^2$ near $W$, the change of loss is
$\langle G,s(\Delta B\,A+B\,\Delta A)\rangle$ plus terms of second order in the step sizes, and the first term is the
expression above with $\eta$ in place of $\eta\,dt$. With an unbiased minibatch gradient $\hat G$, $\mathbb E\hat G=G$, the
first-order change is $-s^2(\eta_B\langle GA\T,\hat GA\T\rangle+\eta_A\langle B\T G,B\T\hat G\rangle)$, whose expectation is
the same expression.

\emph{Adam.} With bias correction, Adam's first moment and second moment after one step are $\hat m_1=g$ and
$\hat v_1=g\odot g$, where $g$ is the first gradient, so the first update is $-\eta\,g/(\lvert g\rvert+\varepsilon)$
entrywise, which tends to $-\eta\operatorname{sign}(g)$ as $\varepsilon\to0$ (entries with $g=0$ do not move). For $B$,
$g=sGA_0\T$, so $\delta B=-\eta_B\operatorname{sign}(s)\operatorname{sign}(GA_0\T)$, and its first-order effect on the loss is
$\langle G,s\,\delta B\,A_0\rangle=s\langle GA_0\T,\delta B\rangle=-\eta_B\lvert s\rvert\|GA_0\T\|_1$. The $A$-term is the same
with $B_0\T G$, and the cross term $s\,\delta B\,\delta A$ and the Taylor remainder are $O(\eta^2)$.

\emph{The alignment remark after the proposition.} Among $A_0$ with $A_0A_0\T=cI$, write $A_0=\sqrt c\,P$ with $P$ of
orthonormal rows. The $B$-term of the gradient-flow rate is $s^2\eta_B\|GA_0\T\|_F^2=s^2\eta_Bc\,\|GP\T\|_F^2$, and by Ky Fan's
maximum principle (\auxref{B.17}(a)) its largest value over $P$ is $s^2\eta_Bc\sum_{i\le r}\sigma_i(G)^2$, attained when
$\operatorname{row}A_0$ is a top-$r$ right singular subspace of $G$. Under Adam the first-step term
$\|GA_0\T\|_1$ is not invariant under $A_0\mapsto OA_0$ with $O\in\Ort(r)$, so it depends on the basis of
$\operatorname{row}A_0$ and not only on the subspace.
\end{proof}

% ======================================================================================================================
\subsection{Lemma II.4}\label{proof:II.4}

\begin{restate}{Lemma}{II.4}{zero gate}
Let $Q=U\times V$ with $U$ a vector space, let $\rho=\theta_0+\beta(u,v)$ with $\beta$ smooth, and suppose
$\beta(u_0,v)=0$ for all $v$; for example $\beta(\cdot,v)$ linear and $u_0=0$. Then $D_v\rho_{q_0}=0$. Hence
$\nabla_v(L\circ\rho)(q_0)=0$, $S_1=D_u\beta(u_0,v_0)(U)$ ($=\beta(U,v_0)$ in the linear case), and the cometric at
$q_0$ is $K_{q_0}=D_u\beta\,\Lambda_{uu}\,D_u\beta\T$, where $\Lambda_{uu}$ is the $uu$-block of $\Lambda$; no
block-diagonality is needed, since the $v$-columns of $D\rho$ vanish. The weight trajectory $\theta(t)$ agrees with that
of the method with $v$ frozen at $v_0$ to first order in $t$ (gradient flow), or in the step size (one or more Adam
steps). The parameter $v$ itself may still move: under terms not of the form $L\circ\rho$ (regularisers, decoupled weight
decay), and under Adam, whose step-2 move of $v$ is $O(\eta)$ although its effect on $\theta$ is $O(\eta^2)$.

\emph{Function-level variant.} For an extension with realised function $F(u,v,x)$, if $F(u_0,v,x)=F_0(x)$ for all $v$
and $x$, then $D_vF=0$ at $q_0$, with the same conclusions at the level of functions. No linearity is needed.
\end{restate}

\begin{proof}
Since $\beta(u_0,v)=0$ for every $v$, differentiating in $v$ gives $D_v\beta(u_0,v)=0$ for every $v$, in particular at
$v_0$. So $D\rho_{q_0}=[\,D_u\beta(u_0,v_0)\ \ 0\,]$. Then $\nabla_v(L\circ\rho)(q_0)=D_v\rho_{q_0}\T\nabla L(\theta_0)=0$,
$S_1=D_u\beta(u_0,v_0)(U)$, and in the linear case $D_u\beta(0,v_0)u=\beta(u,v_0)$. For any symmetric $\Lambda$ written in
blocks,
\begin{equation*}
  K_{q_0}=\begin{bmatrix}D_u\beta&0\end{bmatrix}\begin{bmatrix}\Lambda_{uu}&\Lambda_{uv}\\ \Lambda_{vu}&\Lambda_{vv}
  \end{bmatrix}\begin{bmatrix}D_u\beta\T\\ 0\end{bmatrix}=D_u\beta\,\Lambda_{uu}\,D_u\beta\T .
\end{equation*}

\emph{Gradient flow.} For $\dot q=-\Lambda\nabla(L\circ\rho)(q)$, the weights $\theta(t)=\rho(q(t))$ satisfy
$\dot\theta(0)=-K_{q_0}\nabla L(\theta_0)$. The method with $v$ frozen, $u\mapsto\rho(u,v_0)$ with the metric
$\Lambda_{uu}$, has the same velocity at $t=0$, so the two weight trajectories agree to first order in $t$.

\emph{Adam.} Let $\eta$ be the step size and $\varepsilon>0$ fixed. At step one the $v$-gradient vanishes, so Adam's
first and second moments in $v$ vanish and $v$ does not move; the $u$-update is computed from the same gradient as for the
frozen method. So after one step the two methods have the same parameters and the same weights. At step two, the gradient
in $v$ is $D_v\beta(u_1,v_0)\T\nabla L$. Since $\beta$ is smooth and $D_v\beta(u_0,\cdot)\equiv0$, we have
$\|D_v\beta(u,v)\|=O(\|u-u_0\|)$ uniformly for $v$ near $v_0$, and $\|u_1-u_0\|=O(\eta)$, so this gradient is $O(\eta)$,
and the move of $v$ is at most of order $\eta$; it is of exact order $\eta$ when $\varepsilon$ is negligible against the
gradient, since Adam normalises it. Its effect on $\theta$ is
\begin{equation*}
  \rho(u,v)-\rho(u,v_0)=\int_0^1D_v\beta\big(u,v_0+\tau(v-v_0)\big)(v-v_0)\,d\tau=O(\|u-u_0\|\,\|v-v_0\|)=O(\eta^2).
\end{equation*}
For any fixed number $k$ of steps the same estimate propagates: the Adam update is a smooth function of the gradient
history when $\varepsilon>0$, and by induction on the step, $v_t-v_0=O(\eta)$, the $u$-iterates of the two methods
differ by $O(\eta^2)$ (their gradients differ by $O(\eta)$ and are multiplied by $\eta$), and the weights differ by
$O(\eta^2)$ by the display. For $\varepsilon\to0$ the same argument holds as long as no entry of the frozen method's
gradient vanishes in the first $k$ steps, so that the normalisation stays smooth. Terms that are not of the form
$L\circ\rho$, such as a regulariser $R(v)$ or decoupled weight decay $v\mapsto(1-\eta\lambda)v$, move $v$ already at step
one, which is why the statement concerns the weight trajectory.

\emph{Function-level variant.} If $F(u_0,v,x)=F_0(x)$ for all $v$ and $x$, then $D_vF(u_0,v_0,x)=0$ for every $x$, and the
argument above applies verbatim with the realised function in place of $\theta$.
\end{proof}

% ======================================================================================================================
\subsection{Theorem II.5}\label{proof:II.5}

\begin{restate}{Theorem}{II.5}{complete invariants of the basic actions; classical}
Let $W,W'\in\R^{m\times n}$.
\begin{enumerate}[label=(\arabic*)]
  \item \emph{Right orthogonal (input side).} $W'=WR$ for some $R\in\Ort(n)$ iff $WW\T=W'W'^{\top}$. For $\SO(n)$, when
    $\rk W=n$, also require that the then unique $R$ have $\det R=1$.
  \item \emph{Left orthogonal.} $W'=PW$ for some $P\in\Ort(m)$ iff $W\T W=W'^{\top}W'$.
  \item \emph{Two-sided.} $W'=PWR$ for some $P\in\Ort(m)$ and $R\in\Ort(n)$ iff $W$ and $W'$ have the same singular
    values, counted with multiplicity.
  \item \emph{Left torus.} If $W$ has no zero row, $W'=DW$ for some $D\in T_m^+$ iff $w_i'\in\R_{>0}\,w_i$ for every $i$.
    For the full torus $T_m$, whose signs are unconstrained, the complete invariant is the set of zero rows together with
    the row \emph{lines}.
  \item \emph{Full Hadamard torus.} $W'=E\odot W$ for some $E\in(\R^\times)^{m\times n}$ iff $W$ and $W'$ have the same
    zero pattern.
\end{enumerate}
\end{restate}

\begin{proof}
Write $w_i,w_i'\in\R^n$ for the rows of $W$ and $W'$.

\emph{(1)} If $W'=WR$ with $R$ orthogonal, then $W'W'^{\top}=WRR\T W\T=WW\T$. Conversely, let $WW\T=W'W'^{\top}$. For
$c\in\R^m$,
\begin{equation*}
  \Big\lVert\sum_ic_iw_i\Big\rVert^2=c\T WW\T c=c\T W'W'^{\top}c=\Big\lVert\sum_ic_iw_i'\Big\rVert^2 .
\end{equation*}
Hence $\varphi:\sum_ic_iw_i\mapsto\sum_ic_iw_i'$ is well defined on $\operatorname{row}W$ (if $\sum_ic_iw_i=\sum_ic_i'w_i$, the
vector $\sum_i(c_i-c_i')w_i'$ has norm $0$), linear, isometric and onto $\operatorname{row}W'$. In particular both row spaces
have dimension $\rk WW\T$. By Witt's extension theorem (\auxref{B.25}(b)), $\varphi$ extends to some $S\in\Ort(n)$, for
instance by an isometry between the orthogonal complements, which have equal dimension. Then $Sw_i=w_i'$ for every $i$,
that is $W'=WS\T$, and $R=S\T\in\Ort(n)$.

For $\SO(n)$: if $\rk W<n$, the complement $(\operatorname{row}W)^\perp$ is nonzero, and composing the extension with a
reflection of that complement changes $\det S$ without changing $Sw_i$; so some $R\in\SO(n)$ exists as soon as some
$R\in\Ort(n)$ does. If $\rk W=n$, then $W$ has full column rank and $WR=W'$ determines $R=W^+W'$ uniquely, so
$W'\in W\cdot\SO(n)$ iff this $R$ has determinant $1$.

\emph{(2)} $W'=PW$ iff $W'^{\top}=W\T P\T$ with $P\T\in\Ort(m)$, and by (1) applied to $W\T$ and $W'^{\top}$ this happens
iff $W\T W=W'^{\top}W'$.

\emph{(3)} If $W'=PWR$, then $W'^{\top}W'=R\T(W\T W)R$ has the eigenvalues of $W\T W$, so the singular values agree with
multiplicity. Conversely, if $W=U\Sigma V\T$ and $W'=U'\Sigma V'^{\top}$ are singular value decompositions with the same
$\Sigma$ (\auxref{B.11}(a)), then $W'=(U'U\T)\,W\,(VV'^{\top})$ with both factors orthogonal.

\emph{(4)} For $D=\diag(d)$ the rows of $DW$ are $d_iw_i$. If $d\in\R^m_{>0}$, then $w_i'=d_iw_i\in\R_{>0}w_i$. Conversely,
if $w_i'=d_iw_i$ with $d_i>0$ for every $i$ ($d_i$ is unique because $w_i\ne0$), then $W'=\diag(d)W$ with
$\diag(d)\in T_m^+$. For $T_m$, the rows $d_iw_i$ with $d_i\ne0$ vanish exactly where $w_i$ does and otherwise span the line
$\R w_i$. Conversely, if $W$ and $W'$ have the same zero rows and $\R w_i'=\R w_i$ for every nonzero row, then $w_i'=d_iw_i$
with $d_i\ne0$, and with $d_i=1$ on the zero rows $W'=\diag(d)W$, $\diag(d)\in T_m$.

\emph{(5)} $(E\odot W)_{ij}=E_{ij}W_{ij}$ vanishes iff $W_{ij}$ does, because $E_{ij}\ne0$. Conversely, if $W$ and $W'$
have the same zero pattern, put $E_{ij}=W'_{ij}/W_{ij}$ on the support and $E_{ij}=1$ off it; then $E$ has no zero entry
and $W'=E\odot W$.
\end{proof}

\noindent Item (1) is also the orbit-separation statement of the first fundamental theorem for $\Ort(n)$
(\auxref{B.25}(a)); the proof above uses only the extension of isometries.

% ======================================================================================================================
\subsection{Corollary II.6}\label{proof:II.6}

\begin{restate}{Corollary}{II.6}{what exactly orthogonal methods can never do}
Let a method have image in $W_0\cdot\Ort(n)$ (input side), and work in exact arithmetic. This covers:
\begin{itemize}
  \item OFT and block OFT with an exactly orthogonal parametrisation: the original paper's exact Cayley map, and HF PEFT
    from 0.16 with \texttt{use\_cayley\_\allowbreak neumann=False} on Linear and Conv targets;
  \item GOFT with strict Givens rotations and no scaler, as in the paper (\texttt{no\_scaling=True} in the official
    peft-givens code);
  \item HRA, and ETHER as defined in its paper (not ETHER+);
  \item any number of merge-and-restart cycles of these.
\end{itemize}
Then:
\begin{itemize}
  \item it preserves $K_{\mathrm{out}}$, hence every neuron norm and the angle between any two nonzero neurons, and so the
    hyperspherical energy, which depends only on these angles;
  \item it preserves the whole singular spectrum of $W_0$;
  \item block-diagonal OFT with input blocks $J$, with the same partition in every cycle, preserves each partial Gram
    matrix $W_{:,J}W_{:,J}^{\top}$.
\end{itemize}
\end{restate}

\begin{proof}
\emph{Coverage.} Each listed parametrisation produces $W=W_0R$ with $R\in\Ort(n)$, read in the convention $y=Wx$, where an
input-side rotation acts on the columns of $W_0$. The exact Cayley map $C(\Omega)=(I+\Omega)(I-\Omega)^{-1}$ of a skew
matrix is orthogonal (\auxref{B.23}(a)); HF PEFT's input-side OFT merges to $W_0R\T$, and $R\T$ is orthogonal too. A
block-diagonal matrix with orthogonal blocks is orthogonal. A Givens rotation is orthogonal, and so is a product of them.
HRA multiplies Householder reflections $H_u=I-2uu\T/\lVert u\rVert^2$, which are orthogonal, and ETHER as defined in its
paper multiplies $W_0$ on the input side by one reflection $I-2\hat u\hat u\T$ with $\lVert\hat u\rVert=1$. For merge and
restart, if every weight reachable in cycle $k$ is $W_{k-1}R_k$ with $R_k\in\Ort(n)$, then by induction
$W_k=W_0R_1\cdots R_k$, and $R_1\cdots R_k\in\Ort(n)$.

\emph{Gram matrix and spectrum.} For $W=W_0R$ with $R\in\Ort(n)$, $K_{\mathrm{out}}(W)=W_0RR\T W_0\T=K_{\mathrm{out}}(W_0)$
(\thmref{Theorem}{II.5}(1)). The neuron norms are $\lVert w_i\rVert^2=K_{\mathrm{out}}(W)_{ii}$, and the cosine between two
nonzero neurons is $K_{\mathrm{out}}(W)_{ij}/(K_{\mathrm{out}}(W)_{ii}K_{\mathrm{out}}(W)_{jj})^{1/2}$, so both are
preserved. The hyperspherical energy of \citet{liu2018learning} is a function of the distances between the normalised
neurons, $\lVert\hat w_i-\hat w_j\rVert^2=2-2\cos_{ij}$, so it is preserved as well. The singular values of $W$ are the
square roots of the eigenvalues of $K_{\mathrm{out}}(W)=K_{\mathrm{out}}(W_0)$ (equivalently, \thmref{Theorem}{II.5}(3)
with $P=I$).

\emph{Partial Grams.} Let $R=\operatorname{blockdiag}(R_J)_J$ be block-diagonal for a partition of the input coordinates
into blocks $J$, with orthogonal blocks. The columns of $W_0R$ indexed by $J$ are $(W_0R)_{:,J}=(W_0)_{:,J}R_J$, so
$W_{:,J}W_{:,J}^{\top}=(W_0)_{:,J}R_JR_J\T(W_0)_{:,J}^{\top}=(W_0)_{:,J}(W_0)_{:,J}^{\top}$. With the same partition in
every cycle, $R_1\cdots R_k$ is again block-diagonal with orthogonal blocks, so the identity survives merging.
\end{proof}

The remarks after the corollary in \S\ref{sec:erlangen-orthogonal} make several claims that do not follow from its
hypothesis. We prove them in the same order. Throughout, $K_{\mathrm{out}}(W)=WW\T$, and for an input-side factor $G$ we have
$K_{\mathrm{out}}(W_0G)=W_0\,GG\T\,W_0\T$.

\emph{BOFT and HF OFT 0.14--0.15.} For $W=\diag(s)W_0R\T$ with $R$ orthogonal,
\begin{equation*}
  K_{\mathrm{out}}(W)=\diag(s)\,W_0R\T RW_0\T\diag(s)=\diag(s)\,K_{\mathrm{out}}(W_0)\,\diag(s),
\end{equation*}
with entries $s_is_jK_{\mathrm{out}}(W_0)_{ij}$. If $s_is_j\ne0$ and the neurons $i,j$ of $W_0$ are nonzero,
\begin{equation*}
  \cos'_{ij}=\frac{s_is_jK_{\mathrm{out}}(W_0)_{ij}}{\lvert s_i\rvert\lvert s_j\rvert\,(K_{\mathrm{out}}(W_0)_{ii}
  K_{\mathrm{out}}(W_0)_{jj})^{1/2}}=\operatorname{sign}(s_is_j)\cos_{ij} .
\end{equation*}
So $\lvert\cos_{ij}\rvert$ is preserved whenever $s_is_j\ne0$, and $\cos_{ij}$ itself whenever $s_is_j>0$. If every $s_i$ is
nonzero and all share a sign, every cosine is preserved. Conversely, if every cosine of $W_0$ is nonzero and every cosine is
preserved, then every $s_i$ is nonzero (otherwise neuron $i$ of $W$ vanishes and its cosines are undefined) and
$\operatorname{sign}(s_is_j)=1$ for all $i\ne j$, so all $s_i$ share a sign. With the butterfly permutation equal to the
identity, $R$ is block-diagonal with orthogonal blocks $R_J$ over a partition of the input coordinates, and
$W_{:,J}=\diag(s)(W_0)_{:,J}R_J\T$, so $W_{:,J}W_{:,J}^{\top}=\diag(s)(W_0)_{:,J}(W_0)_{:,J}^{\top}\diag(s)$. HF PEFT OFT
0.14--0.15 has the same form with its output scale \texttt{oft\_s} in place of $s$, which gives the same conclusions.

\emph{GOFT with its scaler.} For $W=W_0\diag(g)G$ with $G$ orthogonal, whatever $G$ is,
\begin{equation*}
  K_{\mathrm{out}}(W)=W_0\diag(g)\,GG\T\diag(g)\,W_0\T=W_0\diag(g)^2W_0\T .
\end{equation*} This differs from
$K_{\mathrm{out}}(W_0)$ in general (for $W_0=I_n$ it is $\diag(g)^2$), so the image is not contained in $W_0\cdot\Ort(n)$
and the corollary does not apply.

\emph{Output-side implementations.} If $W=PW_0$ with $P\in\Ort(m)$, then $K_{\mathrm{in}}(W)=W_0\T P\T PW_0=K_{\mathrm{in}}(W_0)$,
and the singular values are preserved by \thmref{Theorem}{II.5}(3). If $P$ is block-diagonal with orthogonal blocks $P_I$
over a partition of the output coordinates, the rows indexed by $I$ satisfy $W_{I,:}=P_I(W_0)_{I,:}$, so
$W_{I,:}^{\top}W_{I,:}=(W_0)_{I,:}^{\top}(W_0)_{I,:}$. In general $K_{\mathrm{out}}(PW_0)=PK_{\mathrm{out}}(W_0)P\T$
differs from $K_{\mathrm{out}}(W_0)$, for instance for $W_0=E_{11}$ and $P$ a rotation moving $e_1$. ETHER's code with
\texttt{ether\_nb}${}>1$ multiplies the rows of output group $g$ on the right by a reflection $H_g$, $W_{I_g,:}=
(W_0)_{I_g,:}H_g$. Within a group, $W_{I_g,:}W_{I_g,:}^{\top}=(W_0)_{I_g,:}(W_0)_{I_g,:}^{\top}$, which keeps neuron norms
and within-group angles. Across groups the inner products become $w_{0,i}\T H_gH_{g'}w_{0,j}$. For two groups of one row
each, with $w_{0,1}=w_{0,2}=e_1$ in $\R^2$, $H_1=I$ and $H_2$ the reflection that swaps $e_1$ and $e_2$, the two neurons
become $e_1$ and $e_2$: the cosine drops from $1$ to $0$, and the singular values change from $(\sqrt2,0)$ to $(1,1)$.

\emph{HF PEFT's default OFT.} Let $Q$ be skew. Then $(I-Q)(I+Q+Q^2+Q^3)=I-Q^4$ and $I-Q$ is invertible
(\auxref{B.23}(a),(e)), so
\begin{equation*}
  R:=I+2Q+2Q^2+2Q^3+Q^4=(I+Q)(I+Q+Q^2+Q^3)=(I+Q)(I-Q)^{-1}(I-Q^4)=C(Q)\,(I-Q^4),
\end{equation*}
where all factors are polynomials in $Q$ or their inverses and commute. The matrix $Q^4$ is symmetric, since
$(Q^4)\T=(-Q)^4$, and $C(Q)$ is orthogonal, so $R\T R=(I-Q^4)C(Q)\T C(Q)(I-Q^4)=(I-Q^4)^2$. If $\pm i\theta_k$ are the
eigenvalues of $Q$, then $Q^4$ has eigenvalues $\theta_k^4$ (and $0$), so the singular values of $R$ are
$\lvert1-\theta_k^4\rvert$ (and $1$). Thus $R\T R=I$ only if every $\theta_k$ is $0$ or $2^{1/4}$; in particular
$R\T R\ne I$ for every $Q\ne0$ with $\lVert Q\rVert_{\mathrm{op}}<2^{1/4}$, and $R$ is singular when some $\theta_k=1$.
The adapted weight is $W=W_0R\T$, with $K_{\mathrm{out}}(W)=W_0R\T RW_0\T=W_0(I-Q^4)^2W_0\T$, and
$(I-Q^4)^2-I=Q^8-2Q^4$ is of order $\lVert Q\rVert^4$. When $\lVert Q\rVert_{\mathrm{op}}\le2^{1/4}$, every
$\lvert1-\theta_k^4\rvert\le1$, so $R$ is a contraction, and each merge $W\mapsto WR\T$ gives
$WR\T RW\T\preceq WW\T$ and lowers every singular value (\auxref{B.18}(b),(c)); \thmref{Theorem}{V.3}, corollary 4, records
the iteration. In PEFT 0.16 and 0.17 the last coefficient was $2$, $R=p(Q)$ with $p(x)=1+2x+2x^2+2x^3+2x^4$. Then
$R\T=p(-Q)$, and writing $p=E+O$ with even part $E=1+2x^2+2x^4$ and odd part $O=2x+2x^3$,
\begin{equation*}
  p(x)p(-x)=E^2-O^2=(1+4x^2+8x^4+8x^6+4x^8)-(4x^2+8x^4+4x^6)=1+4x^6+4x^8,
\end{equation*}
so $R\T R=I+4Q^6+4Q^8$. Its eigenvalues are $1+4\theta_k^6(\theta_k^2-1)$, so it differs from $I$ unless every $\theta_k$ is $0$
or $1$, and the deviation is of order $\lVert Q\rVert^6$.

\emph{One non-orthogonal plane changes a cosine.} We prove the following, which covers qGOFT, ETHER+ and HF's default OFT.
Let $n>2$ and let $G\in\R^{n\times n}$ be an input-side factor such that $S=GG\T$ acts as the identity on the orthogonal
complement of a plane $\Pi$, maps $\Pi$ into itself, and $S|_\Pi\ne I_\Pi$ (that is, $G$ is not orthogonal on $\Pi$).
Then for $W_0$ outside a nowhere-dense null set some cosine of $W_0G$ differs from the corresponding cosine of $W_0$.
Indeed, for two rows $w_1,w_2$ of $W_0$ put
\begin{equation*}
  \phi(W_0)=(w_1\T Sw_2)^2\,\lVert w_1\rVert^2\lVert w_2\rVert^2-(w_1\T w_2)^2\,(w_1\T Sw_1)(w_2\T Sw_2),
\end{equation*}
a polynomial in $W_0$. Where both cosines are defined, $\cos'_{12}=\cos_{12}$ forces $\phi(W_0)=0$. Choose a unit vector
$c\perp\Pi$ (possible since $n>2$) and $x\in\Pi$ with $x\T Sx\ne\lVert x\rVert^2$ (possible since $S|_\Pi-I_\Pi$ is a
nonzero symmetric map), and take $w_1=c$, $w_2=c+x$. Then $w_1\T Sw_2=1$, $w_1\T Sw_1=1$, $w_2\T Sw_2=1+x\T Sx$,
$\lVert w_2\rVert^2=1+\lVert x\rVert^2$ and $w_1\T w_2=1$, so $\phi=\lVert x\rVert^2-x\T Sx\ne0$. By \auxref{B.10}(b), $\phi\ne0$
on an open dense set of full measure, which proves the claim. A conformal block $\sqrt c\,\mathrm{Rot}$ with $c\ne1$ on two
input coordinates has $S|_\Pi=cI$, so it is covered. For qGOFT, take one block non-orthogonal and the others equal to $I$. For
ETHER+, take the output factor equal to $I$ (its two vectors equal) and the input factor $I-uu\T+vv\T$ with unit $u\ne\pm v$:
it is symmetric, equal to $I$ on $\{u,v\}^\perp$, preserves $\Pi=\spanop\{u,v\}$, and on $\Pi$ has eigenvalues
$1\pm(1-(u\T v)^2)^{1/2}$, so it is not orthogonal there; for $u\perp v$ it kills $u$ and is singular. For HF's default
OFT, take $Q=\theta(e_1e_2\T-e_2e_1\T)$ with $\theta\notin\{0,2^{1/4}\}$; then $R\T R=(I-Q^4)^2$ is $(1-\theta^4)^2I$ on
$\Pi=\spanop\{e_1,e_2\}$ and $I$ elsewhere.

\emph{Finite precision.} The last remark concerns floating-point arithmetic, and Appendix~\ref{app:repro} reports the
measured drift; in exact arithmetic the corollary holds as proved.

% ======================================================================================================================
\subsection{Theorem II.7}\label{proof:II.7}

\begin{restate}{Theorem}{II.7}{HRA is the meet of LoRA with the rotation orbit, and it starts at an apex}
Let $W_0\in\R^{m\times n}$ be injective ($\rk W_0=n$), and let $r$ be even with $2\le r\le n-2$. Define
$\mathrm{HRA}_r:(u_1,\dots,u_r)\mapsto W_0H_{u_1}\cdots H_{u_r}$ with $u_i\ne0$, where
$H_u=I-2uu\T/\lVert u\rVert^2$. Write $\mathrm{OFT}_{\SO}:=W_0\cdot\SO(n)$ for the rotation orbit. (The $\Ort(n)$-orbit
would add the odd component: $W_0H_{e_1}$ lies in $W_0+\Mr$ but not in $\Im\mathrm{HRA}_r$.)
\begin{enumerate}[label=(\alph*)]
  \item \textbf{Meet.} $\Im\mathrm{HRA}_r=(W_0+\Mr)\cap W_0\cdot\SO(n)$. Further, $d(\mathrm{HRA}_r)=rn-\tfrac{r(r+1)}2$,
    so the generic fibre has dimension $\tfrac{r(r+1)}2=rn-d$. This gauge is non-linear. The evident continuous global
    symmetries are the $r$ scalings $u_i\mapsto c_iu_i$; the braid group also acts globally by Hurwitz moves
    $(u_i,u_{i+1})\mapsto(H_{u_i}u_{i+1},u_i)$, since $H_aH_b=H_{H_ab}H_a$.
  \item \textbf{First-order collapse.} At the paired initialisation $u_{2i-1}=u_{2i}=w_i$ ($i\le k=r/2$, with the $w_i$
    independent and $U=\spanop\{w_i\}$),
    \begin{equation*}
      S_1=W_0\cdot\{\Omega\in\so(n):P_{U^\perp}\Omega P_{U^\perp}=0\},\qquad\dim S_1=kn-\tfrac{k(k+1)}2<d(\mathrm{HRA}_r).
    \end{equation*}
  \item \textbf{Vertex.} The germ of $\Im\mathrm{HRA}_r$ at $W_0$ is diffeomorphic to the germ at $0$ of the cone
    $\Sigma_r=\{\Omega\in\so(n):\rk\Omega\le r\}$ (the skew-determinantal cone), which is not linear. So $W_0$ is a
    singular point (the vertex) of the image. Together with (b), the paired initialisation is an \textbf{apex}.
  \item \textbf{Scope.}
    \begin{itemize}
      \item \emph{Not a product.} (a) is a meet in the poset of subsets of $\Theta$ containing $W_0$
        (\thmref{Observation}{I.4}(c)). HRA is not the
        categorical product of $\mathrm{LoRA}_r$ and $\mathrm{OFT}_{\SO}$: the set-level fibre product
        $\{(B,A,R):W_0+BA=W_0R\}$ has local dimension $d+r^2=rn+\tfrac{r(r-1)}2\ne rn$ over a generic point of the
        meet. (It is not pure-dimensional: over $W_0$ alone it contains $\{(B,0,I)\}$, of dimension $mr$.)
      \item \emph{Exactly orthogonal OFT is smaller.} The Cayley map never produces a rotation with eigenvalue $-1$.
        For exact full-block Cayley OFT (the original paper; HF with \texttt{use\_cayley\_\allowbreak neumann=False}),
        $(W_0+\Mr)\cap\Im\mathrm{OFT}$ is dense in $\Im\mathrm{HRA}_r$ but not equal to it: $W_0H_{e_1}H_{e_2}$, with
        $H_{e_1}H_{e_2}=\diag(-1,-1,1,\dots,1)$, is missing. For exact block Cayley OFT with $b<n$ the intersection is a
        proper subset of $\Im\mathrm{HRA}_r$.
      \item \emph{HF's default OFT is incomparable.} HF PEFT's default OFT is not orthogonal
        (\thmref{Corollary}{II.6}); its flag is \texttt{use\_cayley\_\allowbreak neumann=\allowbreak True}. Since $R\T R=(I-Q^4)^2=I$ forces $Q=0$ near $Q=0$, its image meets
        $\Im\mathrm{HRA}_r$ near $W_0$ only at $W_0$, while its intersection with $W_0+\Mr$ contains non-orthogonal
        points such as $W_0(I+2Q+2Q^2+2Q^3+Q^4)$ with $\rk Q=2$.
      \item \emph{Implementation.} The theorem concerns HRA with $\lambda<\infty$ (HF \texttt{apply\_GS=False}). The
        strictly orthogonal Gram--Schmidt variant has the smaller image $W_0(I-2P_U)$, of dimension $r(n-r)$. The paired
        initialisation is HF's default (\texttt{repeat\_interleave}) and is not stated in the HRA paper. HF HRA acts on
        \texttt{in\_features}, so injectivity requires \texttt{out\_features}${}\ge{}$\texttt{in\_features} (and full
        column rank); the theorem does not cover \texttt{down\_proj} or GQA key/value projections. The Gram--Schmidt
        variant never contains $W_0$, since $I-2P_U\ne I$, and at HF's default paired initialisation it runs
        Gram--Schmidt on duplicated columns, which is numerically degenerate.
    \end{itemize}
\end{enumerate}
\end{restate}

\begin{proof}
For $u\ne0$ write $\hat u=u/\lVert u\rVert$. The reflection $H_u=I-2\hat u\hat u\T$ is symmetric and orthogonal, has
determinant $-1$, fixes $u^\perp$ pointwise, and $H_u-I$ has rank one. For $g\in\Ort(n)$, $gH_ug^{-1}=H_{gu}$. For
$x,y\in\R^n$ put $x\wedge y=xy\T-yx\T\in\so(n)$, and for a subspace $A\subseteq\R^n$ let
$\R^n\wedge A=\spanop\{x\wedge a:x\in\R^n,\ a\in A\}$. Because $W_0$ is injective, $X\mapsto W_0X$ is an injective linear map
$\R^{n\times n}\to\R^{m\times n}$; in particular $W_0R=W_0R'$ implies $R=R'$.

\emph{Two lemmas.} (i) For every subspace $A$ of dimension $a$,
\begin{equation*}
  \R^n\wedge A=\{\Omega\in\so(n):P_{A^\perp}\Omega P_{A^\perp}=0\},\qquad\dim(\R^n\wedge A)=an-\tfrac{a(a+1)}2 .
\end{equation*}
For $x,y\in A^\perp$ and $z\in\R^n$, $a'\in A$, $y\T(z\wedge a')x=(y\T z)(a'^{\top}x)-(y\T a')(z\T x)=0$, which gives
``$\subseteq$''. Take an orthonormal basis $e_1,\dots,e_a$ of $A$ and $e_{a+1},\dots,e_n$ of $A^\perp$. The matrices
$e_p\wedge e_q$ ($p<q$) form a basis of $\so(n)$, and $P_{A^\perp}\Omega P_{A^\perp}=0$ says exactly that the coefficients
with $p,q>a$ vanish, so the right side is spanned by the $e_p\wedge e_q$ with $p\le a$, which lie in $\R^n\wedge A$. Its
dimension is $\dim\so(n)-\dim\so(n-a)=an-a(a+1)/2$.

(ii) Let $R(u)=H_{u_1}\cdots H_{u_r}$, put $g_k=H_{u_1}\cdots H_{u_{k-1}}$ (so $g_1=I$) and $a_k=g_k\hat u_k$, and let
$A(u)=\spanop\{a_1,\dots,a_r\}$. Then the right-translated differential satisfies
\begin{equation*}
  \{dR_u(\dot u)\,R(u)^{-1}:\dot u\in(\R^n)^r\}=\R^n\wedge A(u).
\end{equation*}
Indeed, varying only $u_k$ in the direction $\dot u_k$ moves $\hat u_k$ by $\delta=(I-\hat u_k\hat u_k\T)\dot u_k/\lVert
u_k\rVert$, which ranges over $\hat u_k^\perp$, and $H_{u_k}$ by $\dot H=-2(\delta\hat u_k\T+\hat u_k\delta\T)$. Since
$H_{u_k}\hat u_k=-\hat u_k$ and $H_{u_k}\delta=\delta$, we get $\dot HH_{u_k}=2(\delta\hat u_k\T-\hat u_k\delta\T)=2\,\delta\wedge
\hat u_k$. Hence
\begin{equation*}
  \dot R\,R^{-1}=g_k\dot H\,H_{u_{k+1}}\cdots H_{u_r}\,R^{-1}=g_k(\dot HH_{u_k})g_k^{-1}=2\,(g_k\delta)\wedge a_k ,
\end{equation*}
using $H_{u_{k+1}}\cdots H_{u_r}R^{-1}=(g_kH_{u_k})^{-1}$ and $g(x\wedge y)g\T=gx\wedge gy$ for orthogonal $g$. As
$\delta$ ranges over $\hat u_k^\perp$, $g_k\delta$ ranges over $a_k^\perp$, and $a_k^\perp\wedge a_k=\R^n\wedge a_k$
because $a_k\wedge a_k=0$. Summing over $k$ gives the claim.

\emph{(a), ``$\subseteq$''.} $R=H_{u_1}\cdots H_{u_r}$ is orthogonal with $\det R=(-1)^r=1$, and it fixes
$\bigcap_iu_i^\perp$, of dimension at least $n-r$, pointwise. So $\rk(R-I)\le r$ and $W_0R-W_0=W_0(R-I)\in\Mr$.

\emph{(a), ``$\supseteq$''.} Let $W_0R=W_0+X$ with $R\in\SO(n)$ and $\rk X\le r$. Since $W_0$ is injective,
$k':=\rk(R-I)=\rk X\le r$. By the Cartan--Dieudonn\'e theorem in Scherk's sharp form (\auxref{B.22}), $R$ is a product of
exactly $k'$ reflections and $\det R=(-1)^{k'}$, so $k'$ is even. Append $(r-k')/2$ pairs $H_vH_v=I$; then $R$ is a product
of exactly $r$ reflections and $W_0R\in\Im\mathrm{HRA}_r$.

\emph{Leaves.} For $V\in\Gr(r,n)$ let $\SO(V)=\{S\in\SO(n):Sx=x\text{ for }x\in V^\perp\}$. Then
\begin{equation*}
  (W_0+\Mr)\cap W_0\cdot\SO(n)=\bigcup_{V\in\Gr(r,n)}W_0\cdot\SO(V).
\end{equation*}
If $S\in\SO(V)$, then $S-I$ vanishes on $V^\perp$, so $\rk(S-I)\le r$. Conversely, if $\rk(R-I)\le r$, the fixed space
$F=\ker(R-I)$ has dimension at least $n-r$, and any $r$-dimensional $V\supseteq F^\perp$ has $V^\perp\subseteq F$, so
$R\in\SO(V)$. A point with $\rk(R-I)=r$ lies on exactly one leaf, $V=\operatorname{im}(R-I)$.

\emph{The $\Ort(n)$-orbit.} $W_0H_{e_1}-W_0=W_0(H_{e_1}-I)$ has rank one, so $W_0H_{e_1}\in W_0+\Mr$. If
$W_0H_{e_1}=W_0R$ with $R=H_{u_1}\cdots H_{u_r}$, injectivity gives $R=H_{e_1}$, whose determinant is $-1\ne(-1)^r$.

\emph{Dimension and fibre.} $\rho=\mathrm{HRA}_r$ is a rational map on the connected open set
$Q=(\R^n\setminus\{0\})^r$, and $d\rho_u(\dot u)=W_0\,dR_u(\dot u)$ has the rank of $dR_u$, which by lemma (ii) equals
$\dim(\R^n\wedge A(u))$. By lemma (i) this is $an-a(a+1)/2$ with $a=\dim A(u)\le r$, a non-decreasing function of $a$ for
$a\le n$. For $u_k=e_k$ we get $g_k=H_{e_1}\cdots H_{e_{k-1}}$ and $a_k=e_k$, so $a=r$ is attained. By \auxref{B.9}(d),
\begin{equation*}
  d(\mathrm{HRA}_r)=\max_u\rk d\rho_u=rn-\tfrac{r(r+1)}2 .
\end{equation*}
The computation uses neither the parity of $r$ nor $r\le n-2$, so the same value holds for every $1\le r\le n$; for
$r\in\{n-1,n\}$ it equals $\dim\SO(n)$. At a point of maximal rank the fibre is locally a submanifold of dimension
$rn-d=r(r+1)/2$ (\auxref{B.1}(b)), and the points of maximal rank, those with $a_1,\dots,a_r$ independent, are open and
dense in $Q$ (\auxref{B.10}). As a check, a generic point of the meet is fixed by its moved space
$\operatorname{im}(R-I)\in\Gr(r,n)$, of dimension $r(n-r)$, and by a fixed-point-free element of $\SO(r)$ on it, an open
subset of dimension $r(r-1)/2$; the sum is $d$.

\emph{Symmetries.} $H_{cu}=H_u$ for $c\ne0$, so the scalings $(\R^\times)^r$ preserve $\rho$. From $gH_bg^{-1}=H_{gb}$
with $g=H_a$ we get $H_aH_b=H_{H_ab}H_a$, so the Hurwitz move $\sigma_i:(u_i,u_{i+1})\mapsto(H_{u_i}u_{i+1},u_i)$ preserves the
product, and $H_{u_i}u_{i+1}\ne0$. It is invertible, with inverse $(u_i,u_{i+1})\mapsto(u_{i+1},H_{u_{i+1}}u_i)$, the moves
$\sigma_i,\sigma_j$ with $\lvert i-j\rvert\ge2$ commute, and $\sigma_1\sigma_2\sigma_1=\sigma_2\sigma_1\sigma_2$: on a triple
both composites give $(a,b,c)\mapsto(H_aH_bc,\,H_ab,\,a)$, using $H_{H_ab}H_ac=H_aH_bc$. So the braid group acts. The
scalings have $r$-dimensional orbits, while the generic fibre has dimension $r(r+1)/2>r$ for $r\ge2$, and the braid moves are
discrete. So the exhibited symmetries are not locally transitive on generic fibres; the remaining fibre directions come from
refactorising $R$ inside its moved space, and this is what we mean by a non-linear gauge.

\emph{(b)} At the paired initialisation, $g_{2i-1}=(H_{w_1}H_{w_1})\cdots(H_{w_{i-1}}H_{w_{i-1}})=I$, so $a_{2i-1}=\hat w_i$
and $a_{2i}=H_{w_i}\hat w_i=-\hat w_i$. Hence $A(u)=U$, $R(u)=I$, and by lemmas (ii) and (i)
\begin{equation*}
  S_1=W_0\cdot(\R^n\wedge U)=W_0\cdot\{\Omega\in\so(n):P_{U^\perp}\Omega P_{U^\perp}=0\},\qquad\dim S_1=kn-\tfrac{k(k+1)}2,
\end{equation*}
since $W_0$ is injective. Equivalently, a direct computation for one pair gives
\begin{equation*}
  \tfrac{\dd}{\dd\varepsilon}\big(H_{w+\varepsilon a}H_{w+\varepsilon b}\big)\big|_{\varepsilon=0}
  =\tfrac2{\lVert w\rVert}\big(v\hat w\T-\hat wv\T\big),\qquad v=(I-\hat w\hat w\T)(a-b),
\end{equation*}
and because every pair equals $I$ at the base point the differential of the product is the
sum of the pair differentials. Finally, with $r=2k$,
\begin{equation*}
  d-\dim S_1=2kn-k(2k+1)-kn+\tfrac{k(k+1)}2=\tfrac{k(2n-3k-1)}2 ,
\end{equation*}
which is positive because $n\ge2k+2$ gives $2n-3k-1\ge k+3$.

\emph{(c)} By \auxref{B.23}(b),(c), the Cayley map $C(\Omega)=(I+\Omega)(I-\Omega)^{-1}$ is a diffeomorphism from $\so(n)$
onto the open set $\mathcal O=\{R\in\SO(n):-1\notin\operatorname{spec}R\}$, which contains $I$, and
$\rk(C(\Omega)-I)=\rk\Omega$. By (a), $\Im\mathrm{HRA}_r\cap W_0\cdot\mathcal O=W_0\cdot C(\Sigma_r)$. The map
$\Omega\mapsto W_0C(\Omega)$ is a diffeomorphism of $\so(n)$ onto an open subset of the embedded submanifold
$W_0\cdot\SO(n)$ (a linear injection of $\SO(n)$), and it carries $\Sigma_r$ onto $\Im\mathrm{HRA}_r$ near $W_0$. So the
germs correspond. $\Sigma_r$ is a closed cone and is not linear (\auxref{B.7}(c)): with $r=2k$, the matrices
$\sum_{j\le k}e_{2j-1}\wedge e_{2j}$ and $e_{r+1}\wedge e_{r+2}$ lie in $\Sigma_r$ (here $r+2\le n$), and their sum has rank
$r+2$. By \auxref{B.6}, $\Sigma_r$ is not a $C^1$ submanifold near $0$. If $\Im\mathrm{HRA}_r$ were a $C^1$ submanifold near
$W_0$, it would be one of $W_0\cdot\SO(n)$, and its preimage $\Sigma_r$ under the diffeomorphism would be $C^1$ near $0$.
So $W_0$ is a singular point, the vertex of the germ. A singular image alone does not give a positive deficit; the apex
statement is (b).

\emph{(d), not a product.} Let $\mathcal F=\{(B,A,R)\in\R^{m\times r}\times\R^{r\times n}\times\SO(n):W_0+BA=W_0R\}$ be the
set-level fibre product. The points of the meet of the form $W_0C(\Omega)$ with $\rk\Omega=r$ form a $d$-dimensional
submanifold $\mathcal G$, the image under the diffeomorphism of (c) of the rank-$r$ part of $\Sigma_r$, which is open and
dense in $\Sigma_r$ and of dimension $d$ (\auxref{B.7}(c)). $\mathcal G$ is open in the meet, since there $\rk(R-I)\ge r$ and
$-1\notin\operatorname{spec}R$ are open conditions, and it is dense: an $R$ in the meet with eigenvalue $-1$ is a limit of
rotations with the same moved planes and the angles $\pi$ replaced by $\pi-\varepsilon$ (\auxref{B.7}(b)). Over a point
$W_0R\in\mathcal G$, the update $W_0(R-I)$ has rank exactly $r$, so near such a point $\mathcal F$ consists of the pairs
$(B,A)$ of full rank with $BA\in W_0(R-I)$ for $W_0R\in\mathcal G$, with $R$ recovered from $BA$ by injectivity. The map
$(B,A)\mapsto BA$ is a submersion from pairs of full rank onto $\MM_r$ with fibres the $\GL_r$-orbits, of dimension $r^2$
(\auxref{B.5}(b)). A submersion is transverse to every submanifold, so by \auxref{B.2}(a), applied with target $\MM_r$,
$\mathcal F$ is there a submanifold of dimension $d+r^2=rn+r(r-1)/2$. If HRA, or any method with an
$rn$-dimensional parameter space, were the product, its underlying set would be in bijection with $\mathcal F$ through its
two smooth projections (\thmref{Observation}{I.4}(b)), so $\mathcal F$ would be the image of an $rn$-manifold $P$ under
a smooth map $\iota$ into the ambient space $E=\R^{m\times r}\times\R^{r\times n}\times\R^{n\times n}$. This is
impossible, because $\mathcal F$ contains an embedded submanifold $\mathcal N$ of dimension $D=rn+r(r-1)/2>rn$ (the part
over $\mathcal G$). Fix $p\in\mathcal N$ and let $\pi:E\to T_p\mathcal N\cong\R^D$ be the orthogonal projection, which is
$1$-Lipschitz. Its restriction to $\mathcal N$ has invertible differential at $p$, so by the inverse function theorem
$\pi(N')$ contains a closed ball $\bar B$ of positive volume for a compact neighbourhood $N'$ of $p$ in $\mathcal N$. Cover
$P$ by countably many compact sets $K_j=\varphi_j([0,1]^{rn})$, images of cubes under charts, on which $\iota\circ\varphi_j$
is Lipschitz with a constant $L_j$. Cutting the cube into $k^{rn}$ subcubes of side $1/k$ covers $\pi(\iota(K_j))$ by
$k^{rn}$ sets of diameter at most $L_j\sqrt{rn}/k$, so its outer Lebesgue measure is at most
$k^{rn}\,\omega_D\,(L_j\sqrt{rn}/k)^D$, which tends to $0$ as $k\to\infty$ because $D>rn$. The countably many null sets
$\pi(\iota(K_j))$ would then cover $\pi(N')\supseteq\bar B$, which is impossible. For the parenthetical remark: over $W_0$ itself $\mathcal F$ contains $\{(B,0,I)\}$, of dimension $mr$,
while over points of $\mathcal G$ its fibre has dimension $r^2$. So the fibres of $\mathcal F$ over the meet jump, and
whenever $mr>d+r^2$ (for instance when $m\ge n+r$, as at $(m,n,r)=(12,9,4)$, where $mr=48>42$) the local dimension of
$\mathcal F$ at these points exceeds $d+r^2$ and $\mathcal F$ is not pure-dimensional.

\emph{(d), exactly orthogonal OFT.} If $C(\Omega)x=-x$, then $y=(I-\Omega)^{-1}x$ satisfies $(I+\Omega)y=-(I-\Omega)y$,
so $2y=0$ and $x=0$; this is the statement that $C$ misses eigenvalue $-1$. Exact full-block Cayley OFT has image
$W_0\cdot C(\so(n))=W_0\cdot\mathcal O$ (input-side forms $W_0C(\Omega)\T=W_0C(-\Omega)$ give the same set). Its
intersection with $W_0+\Mr$ is $W_0\cdot C(\Sigma_r)\subseteq\Im\mathrm{HRA}_r$, which contains $\mathcal G$ and is
therefore dense in $\Im\mathrm{HRA}_r$. The point $W_0H_{e_1}H_{e_2}$ lies in $\Im\mathrm{HRA}_r$ (pad with pairs), and by
injectivity it equals $W_0R$ only for $R=\diag(-1,-1,1,\dots,1)$, which has eigenvalue $-1$; so it is missing. For block
Cayley OFT with block size $b<n$ the image consists of $W_0R$ with $R$ block-diagonal for a partition with at least two
blocks. Choose coordinates $p,q$ in different blocks and $v=(e_p+e_q)/\sqrt2$. Then $R=H_{e_p}H_v\in\SO(n)$ has
$\rk(R-I)\le2\le r$ and $Re_p=H_{e_p}e_q=e_q$, so $R$ is not block-diagonal. Hence $W_0R\in\Im\mathrm{HRA}_r$ lies outside
the image of block OFT, by injectivity.

\emph{(d), HF's default OFT.} Its image consists of the $W_0R(Q)$, $R(Q)=I+2Q+2Q^2+2Q^3+Q^4$, with $Q$ skew and
block-diagonal (the input-side form uses $R(Q)\T=R(-Q)$, of the same kind). By the proof of \thmref{Corollary}{II.6},
$R(Q)\T R(Q)=(I-Q^4)^2$ has eigenvalues $(1-\theta_k^4)^2$, so $R(Q)$ is orthogonal only if every $\theta_k\in\{0,2^{1/4}\}$.
If $\lVert Q\rVert_{\mathrm{op}}<2^{1/4}$ this forces every $\theta_k=0$, hence $Q=0$ since a skew matrix is normal. As
$\Im\mathrm{HRA}_r\subseteq W_0\cdot\Ort(n)$ and $W_0$ is injective, the parameters near $Q=0$ reach $\Im\mathrm{HRA}_r$
only at $W_0$. For $Q=\theta(e_1e_2\T-e_2e_1\T)$ with $\theta\notin\{0,2^{1/4}\}$ (one plane, inside one block),
$R(Q)-I=Q(2I+2Q+2Q^2+Q^3)$ has rank at most $2\le r$ while $R(Q)$ is not orthogonal, so $W_0R(Q)\in W_0+\Mr$ lies outside
$W_0\cdot\Ort(n)\supseteq\Im\mathrm{HRA}_r$. Conversely, an orthogonal $R(Q)$ acts on each plane of $Q$ with
$\theta_k=2^{1/4}$ as $-C(2^{1/4}J)$, a rotation by $\pi+2\arctan2^{1/4}$, and as the identity on $\ker Q$; so a rotation by a
small angle $\varphi\ne0$ in one plane, which lies in $\Im\mathrm{HRA}_r$ by (a), is not of the form $R(Q)$. The two images
are therefore incomparable.

\emph{(d), implementation.} For orthonormal $\hat u_1,\dots,\hat u_r$,
$(I-2\hat u_1\hat u_1\T)(I-2\hat u_2\hat u_2\T)=I-2\hat u_1\hat u_1\T-2\hat u_2\hat u_2\T$ because $\hat u_1\T\hat u_2=0$, and
by induction $\prod_i(I-2\hat u_i\hat u_i\T)=I-2P_U$ with $U=\spanop\{\hat u_i\}$. So the Gram--Schmidt variant has image
$\{W_0(I-2P_U):U\in\Gr(r,n)\}$. The map $U\mapsto W_0(I-2P_U)$ is injective ($P_U$ determines $U$, and $W_0$ is injective) and
is an immersion of the compact manifold $\Gr(r,n)$, so the image has dimension $r(n-r)$. It never contains $W_0$, since
$I-2P_U\ne I$ for $U\ne0$. With duplicated columns, Gram--Schmidt meets a column that lies in the span of the previous
ones, whose residual vanishes, and normalising it divides by zero up to rounding. Injectivity of
$W_0\in\R^{\texttt{out}\times\texttt{in}}$ needs $\texttt{out}\ge\texttt{in}$ and full column rank, which fails for
\texttt{down\_proj} and for GQA key and value projections.
\end{proof}

% ======================================================================================================================
\subsection{The LOFT example of Exposé II.B}\label{proof:loft}

\begin{restate}{Example}{\unskip}{LOFT: a leaf, and how to choose it; \S\ref{sec:erlangen-loft}}
Let $P_r\in\R^{r\times n}$ have orthonormal rows spanning $V$, and let $\Ort(V)=\{I+P_r\T(T-I)P_r:T\in\Ort(r)\}$ be the
orthogonal maps that fix $V^\perp$ pointwise. LOFT trains $\rho(E)=W_0\big(I+P_r\T(Q(E)-I)P_r\big)$ with $E\in\so(r)$,
$E_0=0$, so $\lvert M\rvert=r(r-1)/2$. We read it with an exactly orthogonal $Q$ (the exponential, or the Cayley map
$C$); the released code defaults to a five-term Cayley--Neumann series with $R\T R=(I-X^6)^2$, which is only
approximately orthogonal, so like HF's default OFT (\thmref{Corollary}{II.6}) that version leaves its leaf.
\begin{enumerate}[label=(\alph*)]
  \item \emph{One leaf.} If $W_0P_r\T$ has rank $r$, then
    $\big(W_0+W_0P_r\T\,\R^{r\times r}P_r\big)\cap W_0\cdot\Ort(n)=W_0\cdot\Ort(V)$. For the principal support
    $P_r=V_r\T$ the first set is the image of LoRA-XS (this needs $\sigma_r(W_0)>0$), so the leaf is
    LoRA-XS${}\cap W_0\cdot\Ort(n)$. Principal-support LOFT reaches an open dense part of its component $W_0\cdot\SO(V)$;
    for injective $W_0$ and a Cayley $Q$, its image is exactly LoRA-XS${}\cap{}$the image of exact full-block Cayley
    OFT.
  \item \emph{Leaves over the Grassmannian.} If $W_0$ is injective,
    $(W_0+\Mr)\cap W_0\cdot\SO(n)=\bigcup_{V\in\Gr(r,n)}W_0\cdot\SO(V)$, which for even $r$ is $\Im\mathrm{HRA}_r$
    (\thmref{Theorem}{II.7}(a)). LOFT's image is the whole leaf $W_0\cdot\SO(V)$ with the exponential, and an open dense
    part of it with a Cayley map, which never reaches eigenvalue $-1$.
  \item \emph{A regular base point.} If $W_0P_r\T$ has rank $r$, then $S_1=W_0P_r\T\so(r)P_r$ has dimension
    $r(r-1)/2=d(M)$: LOFT starts at a smooth point of its leaf and has no gauge, while HRA's paired initialisation is an
    apex (\thmref{Theorem}{II.7}(b)). For $r\ge4$ there is no smooth pointed arrow LOFT${}\to{}$HRA at that
    initialisation, although the images are nested (obstruction test O3 of \S\ref{sec:atlas}).
  \item \emph{Choosing the leaf.} If $Q'(0)=\id$ (the exponential, or $E\mapsto C(E/2)$), the gradient pulled back to $E$
    at $\theta_0$ is $P_r\skewop(W_0\T G)P_r\T$ with $G=\nabla L(W_0)$; the Cayley map $C$ itself doubles it, which does
    not change the best support. Its squared norm is at most $2\sum_{k\le r/2}\mu_k^2$, where $\pm i\mu_k$ are the
    eigenvalues of $\skewop(W_0\T G)$, with equality when $V$ spans the invariant subspace of the top $\lfloor r/2\rfloor$
    pairs.
\end{enumerate}
\end{restate}

\begin{proof}
(a) A rotation $T\in\Ort(r)$ gives $W_0(I+P_r\T(T-I)P_r)=W_0+W_0P_r\T(T-I)P_r$, so ``$\supseteq$'' holds. Conversely let
$W=W_0+BXP_r$ with $B=W_0P_r\T$, and suppose $W=W_0S$ with $S\in\Ort(n)$. Then $WW\T=W_0W_0\T$, and expanding with
$P_rP_r\T=I_r$ gives $B\,(X+X\T+XX\T)\,B\T=0$. Since $B$ has full column rank, $X+X\T+XX\T=0$, that is
$(I+X)(I+X)\T=I$, so $T=I+X\in\Ort(r)$ and $W=W_0(I+P_r\T(T-I)P_r)\in W_0\cdot\Ort(V)$. For $P_r=V_r\T$ we have
$B=U_r\Sigma_r$, and $\{U_r\Sigma_rXV_r\T\}=\{U_rR\,\Sigma_rV_r\T\}$ through $R=\Sigma_rX\Sigma_r^{-1}$, which is
LoRA-XS's image. With the exponential, $T$ ranges over $\SO(r)$ and LOFT reaches the whole component $W_0\cdot\SO(V)$;
with a Cayley map it reaches the $S(T)=I+P_r\T(T-I)P_r$ with $T\in\SO(r)$ and no eigenvalue $-1$, an open dense part of
it. For injective $W_0$, $W_0S$ determines $S$, so meeting with exact full-block Cayley OFT keeps exactly the $S(T)$ with
$T\in\SO(r)$ and no eigenvalue $-1$ (the eigenvalues of $S(T)$ are those of $T$ together with $1$), which is the Cayley
image.

(b) Every $S\in\SO(V)$ has $\rk(S-I)\le r$, which gives ``$\supseteq$''. Conversely, if $W_0S-W_0$ has rank at most $r$
and $W_0$ is injective, then $\rk(S-I)\le r$. The fixed space $F=\ker(S-I)$ has dimension at least $n-r$, and $F^\perp$
is invariant. Any $r$-dimensional $V\supseteq F^\perp$ works: $S$ is the identity on $V^\perp\subseteq F$ and has
determinant $1$ on $V$. The statements about LOFT's image are those of (a).

(c) The differential at $0$ is $D\rho_0(E)=W_0P_r\T Q'(0)(E)P_r$, with $Q'(0)$ invertible. If this vanishes, then $W_0P_r\T Q'(0)(E)=0$ because $P_r$
has full row rank, and $E=0$ because $W_0P_r\T$ has full column rank. So $\dim S_1=\lvert M\rvert\ge d(M)\ge\dim S_1$. A
smooth pointed arrow would give $S_1(\mathrm{LOFT})\subseteq S_1(\mathrm{HRA})$ (\thmref{Observation}{I.6}), and by
\thmref{Theorem}{II.7}(b) and injectivity this asks every $\Omega\in\so(V)$ to satisfy $P_{U^\perp}\Omega P_{U^\perp}=0$.
Since $\dim(V\cap U^\perp)\ge r-r/2\ge2$, any two orthonormal $a,b\in V\cap U^\perp$ give $a\wedge b=ab\T-ba\T$
violating it.

(d) With $Q'(0)=\id$,
\begin{equation*}
  \tfrac{\dd}{\dd t}L\big(\rho(tE)\big)\big|_{t=0}=\langle G,W_0P_r\T EP_r\rangle=\langle P_rW_0\T GP_r\T,E\rangle
  =\langle P_r\skewop(W_0\T G)P_r\T,E\rangle,
\end{equation*}
since $E$ is skew. The squared norm of this gradient is $\lVert P_r\skewop(W_0\T G)P_r\T\rVert^2$, the compression of the
skew matrix $\skewop(W_0\T G)$ to $V$, so the bound and its equality case are \auxref{B.17}(c) with $k=r$.
\end{proof}

% ====================================================================================================
\subsection{Proposition II.8}\label{proof:II.8}

\begin{restate}{Proposition}{II.8}{DoRA is LoRA followed by an output scale, at the level of images}
Assume $1\le r\le\min(m,n)$ and $s\ne0$. Let $N$ normalise rows; DoRA's magnitude $\mu\in\R^m$ is per output neuron, as
in HF PEFT and the official code (the paper's notation says ``column-wise'', in the transposed convention). Images refer
to the merged (eval-mode) weight; HF's train-mode forward with \texttt{lora\_dropout}${}>0$ is not of the form $Wx$. Let
$W_0$ have no zero row, let $\mathcal U$ be the open set of matrices with no zero row, and let $\mathrm{Diag}_m$ be the
monoid of all diagonal matrices. Then
\begin{equation*}
  \Im\mathrm{DoRA}_r=\{\diag(\mu)\,N(V):V\in(W_0+\Mr)\cap\mathcal U,\ \mu\in\R^m\}=\mathrm{Diag}_m\cdot(W_0+\Mr),
\end{equation*}
which is exactly the image of ``$\mathrm{LoRA}_r$ followed by a trainable output scale $\ell\in\R^m$''. With $\mu$
restricted to $(\R^\times)^m$, or equivalently after intersecting with $\mathcal U$:
$\Im\mathrm{DoRA}_r\cap\mathcal U=T_m\cdot(W_0+\Mr)\cap\mathcal U$. Moreover:
\begin{itemize}
  \item $d(\mathrm{DoRA}_r)\le\min\big(mn,\ r(m+n-r)+m\big)$. At a parameter point with $\mu\in(\R^\times)^m$ and
    $X=sBA\in\MM_r$, with $U=\operatorname{col}X$ and $R=\operatorname{row}X$, the Jacobian rank is
    $r(m+n-r)+m-\dim\ker\big(\delta\mapsto P_{U^\perp}\diag(\delta)W_0P_{R^\perp}\big)$; where some $\mu_i=0$ it is
    smaller when $r\ge2$, and for $r=1$ when some row $w_{0,i}$ with $\mu_i=0$ lies outside $R$. For generic $W_0$, $d=\min\big(mn,\ r(m+n-r)+m\big)$; in particular $d=r(m+n-r)+m$ iff
    $(m-r)(n-r)\ge m$, which holds for all practical shapes. For special $W_0$ it is smaller: for $W_0=ac\T$ with all
    $a_i\ne0$ and $r<n$, $d=r(m+n-r)+m-r$, strictly below the generic value iff $r<m$ and $r\le n-2$.
  \item The image covariance group (\thmref{Definition}{II.14}) contains $\mathrm{Mon}_m\times\GL_n$.
  \item At $B=0$, in HF's magnitude coordinates, $\Delta W=\big(\diag(\mu_i/\lVert w_{0,i}\rVert)-I\big)W_0$, which can
    have rank up to $\rk W_0$, so DoRA escapes LoRA's rank bound.
\end{itemize}
\end{restate}

\begin{proof}
DoRA is the map $\rho(\mu,B,A)=\diag(\mu)N(V)$ with $V=W_0+sBA$, defined on the open set
$\mathcal D=\{(\mu,B,A):V\in\mathcal U\}$, where $N(V)$ divides row $i$ by $\lVert v_i\rVert$. Since $s\ne0$ and every matrix
of rank at most $r$ factors with inner dimension $r$ (\auxref{B.11}(a)), $\{sBA\}=\Mr$.

\emph{The image.} $\diag(\mu)N(V)=\diag(\mu_i/\lVert v_i\rVert)V\in\mathrm{Diag}_m\cdot(W_0+\Mr)$, which gives the first
inclusion; the first equality is the definition of $\rho$. Conversely, take $\diag(\ell)V$ with $V=W_0+X$, $\rk X\le r$, and
let $Z=\{i:v_i=0\}$. For $0<\varepsilon\le1$ put
\begin{equation*}
  V'=W_0+\Big(I-\varepsilon\sum_{i\in Z}E_{ii}\Big)X .
\end{equation*}
Its update still has rank at most $r$. Its rows outside $Z$ are those of $V$, and for $i\in Z$ we have $x_i=-w_{0,i}$, so
row $i$ of $V'$ is $w_{0,i}-(1-\varepsilon)w_{0,i}=\varepsilon w_{0,i}\ne0$. Hence $V'\in\mathcal U$. With $\ell_i'=0$ on $Z$
and $\ell'_i=\ell_i$ elsewhere, $\diag(\ell')V'=\diag(\ell)V$, since both vanish on the rows in $Z$. Taking
$\mu_i=\ell_i'\lVert v_i'\rVert$ gives $\diag(\mu)N(V')=\diag(\ell)V$. The set $\mathrm{Diag}_m\cdot(W_0+\Mr)$ is by
definition the image of $(\ell,B,A)\mapsto\diag(\ell)(W_0+sBA)$.

\emph{The torus form.} If $\mu\in(\R^\times)^m$, then $\diag(\mu_i/\lVert v_i\rVert)\in T_m$ and the result has no zero
row. Conversely, if $DV\in\mathcal U$ with $D\in T_m$ and $V\in W_0+\Mr$, then $V\in\mathcal U$ and $\mu_i=D_{ii}\lVert
v_i\rVert\ne0$ realises $DV$. Finally, $\diag(\ell)V\in\mathcal U$ forces every $\ell_i\ne0$ and $V\in\mathcal U$, so
$\Im\mathrm{DoRA}_r\cap\mathcal U=T_m\cdot(W_0+\Mr)\cap\mathcal U$, and these are exactly the values at
$\mu\in(\R^\times)^m$.

\emph{A change of coordinates.} Let $\rho'(\ell,B,A)=\diag(\ell)(W_0+sBA)$ be LoRA followed by an output scale. The map
$\varphi(\mu,B,A)=(\mu_i/\lVert v_i\rVert,B,A)$ is a diffeomorphism of $\mathcal D$ onto
$\{(\ell,B,A):W_0+sBA\in\mathcal U\}$, with inverse $(\ell,B,A)\mapsto(\ell_i\lVert v_i\rVert,B,A)$, and
$\rho=\rho'\circ\varphi$. So the Jacobian rank of $\rho$ at $p$ equals that of $\rho'$ at $\varphi(p)$, and $\mu_i=0$ iff
$\ell_i=0$.

\emph{Upper bound.} $\rho'$ factors through $(\ell,B,A)\mapsto(\ell,W_0+sBA)\in\R^m\times(W_0+\Mr)$, whose image has
dimension $m+r(m+n-r)$ (\auxref{B.5}(a); for $r=\min(m,n)$, $\Mr=\R^{m\times n}$ has dimension $mn=r(m+n-r)$). A polynomial
map does not raise dimension (\auxref{B.9}(b)), and trivially $d\le mn$.

\emph{The Jacobian rank.} The differential of $\rho'$ at $(\ell,B,A)$ is
$(\delta,\dot B,\dot A)\mapsto\diag(\delta)V+\diag(\ell)\,s(\dot BA+B\dot A)$. Let $\ell$ be invertible and $X=sBA\in\MM_r$, so
that $B$ and $A$ have full rank and $\{s(\dot BA+B\dot A)\}=T_X\MM_r=\{D:P_{U^\perp}DP_{R^\perp}=0\}$, of dimension
$r(m+n-r)$ (\auxref{B.5}(b)). Substituting $\delta=\diag(\ell)\delta'$, the image of the differential is
$\diag(\ell)\big[\{\diag(\delta')V\}+T_X\MM_r\big]$, whose dimension is that of the bracket. The first summand has dimension
$m$, because $V$ has no zero row. A matrix $\diag(\delta')V$ lies in $T_X\MM_r$ iff
$P_{U^\perp}\diag(\delta')VP_{R^\perp}=0$, and $VP_{R^\perp}=W_0P_{R^\perp}$ because $XP_{R^\perp}=0$. So the intersection
has the dimension of the kernel $K$ of $\delta\mapsto P_{U^\perp}\diag(\delta)W_0P_{R^\perp}$, and the rank is
$r(m+n-r)+m-\dim K$.

\emph{Zero magnitudes.} Let $Z_0=\{i:\ell_i=0\}\ne\emptyset$, keep $X\in\MM_r$, and write $\mathcal S=\{\diag(\delta)V\}$
and $T=T_X\MM_r$. Let $\Pi_0$ keep the rows in $Z_0$ and $\Pi_1$ the others. In the image of the differential, the rows in
$Z_0$ are $\delta_iv_i$ and involve only $\delta_{Z_0}$, while the other rows do not involve $\delta_{Z_0}$; so the image has
dimension $\lvert Z_0\rvert+\dim\Pi_1(\mathcal S+T)$, after the same substitution as above on the rows outside $Z_0$. On the
other hand, the value $F=\dim(\mathcal S+T)$ of the formula satisfies $F=\dim\Pi_1(\mathcal S+T)+\dim\big((\mathcal S+T)\cap
\ker\Pi_1\big)$, and $(\mathcal S+T)\cap\ker\Pi_1$ contains every $\sum_{i\in Z_0}e_i(c_iv_i+y_i)\T$ with $y_i\in R$, since
$e_iv_i\T\in\mathcal S$ and $e_iy\T\in T$ for $y\in R$. As $v_i-w_{0,i}\in R$, the rank at the point is at most
\begin{equation*}
  F-\sum_{i\in Z_0}\big(\dim\spanop(w_{0,i},R)-1\big)\le F-\lvert Z_0\rvert(r-1),
\end{equation*}
which is smaller than $F$ whenever $r\ge2$ or some $w_{0,i}$ with $i\in Z_0$ lies outside $R$, in particular at generic
points with some $\mu_i=0$.

\emph{The generic value.} If $r=\min(m,n)$, the image is $\mathrm{Diag}_m\cdot\R^{m\times n}=\R^{m\times n}$ and $d=mn$,
which is the claimed value. Let $r<\min(m,n)$, and let $E\in\R^{m\times(m-r)}$ and $F\in\R^{n\times(n-r)}$ have orthonormal
columns spanning $U^\perp$ and $R^\perp$. Then $P_{U^\perp}DP_{R^\perp}=0$ iff $E\T DF=0$, and
\begin{equation*}
  E\T\diag(\delta)W_0F=\sum_i\delta_i\,(E\T e_i)(F\T w_{0,i})\T ,
\end{equation*}
so $m-\dim K$ is the dimension $k$ of the span of the $m$ rank-one matrices $t_i=(E\T e_i)(F\T w_{0,i})\T$ in
$\R^{(m-r)\times(n-r)}$, and the rank is $r(m+n-r)+k$ with $k\le\min\big(m,(m-r)(n-r)\big)$. Since
$r(m+n-r)+(m-r)(n-r)=mn$, it suffices to reach $k=\min(m,(m-r)(n-r))$. A witness: choose $E$ whose rows
$p_i=E\T e_i$ are in general position (every $m-r$ of them independent; orthonormalising the columns of a generic matrix keeps
this), let $U=(\operatorname{col}E)^\perp$, $R=\spanop\{e_1,\dots,e_r\}$ and $F=[e_{r+1}\cdots e_n]$. Assign to each $i$ a
column index $c(i)\in\{1,\dots,n-r\}$: if $m\le(m-r)(n-r)$, use each index at most $m-r$ times; otherwise give each index to
exactly $m-r$ of the first $(m-r)(n-r)$ values of $i$, and arbitrary indices to the rest. Put $w_{0,i}=e_{r+c(i)}$, so $F\T w_{0,i}=e_{c(i)}$ and $t_i=p_ie_{c(i)}\T$ has
$p_i$ in column $c(i)$. Matrices supported in different columns are independent, and in one column at most $m-r$ vectors in
general position are independent, while exactly $m-r$ of them form a basis; so the $t_i$ span $\min(m,(m-r)(n-r))$
dimensions. This $W_0$ has no zero row, and with
$B,A$ of full rank spanning $U$ and $R$ and $\mu_i=\lVert v_i\rVert$ the point lies in $\mathcal D$, since $v_i=w_{0,i}+x_i$
with $x_i\in R\perp w_{0,i}$. Now write $P_{U^\perp}=I-B(B\T B)^{-1}B\T$ and $P_{R^\perp}=I-A\T(AA\T)^{-1}A$. After clearing
the denominators $\det(B\T B)$ and $\det(AA\T)$, the $k\times k$ minors of the matrix of $\delta\mapsto
P_{U^\perp}\diag(\delta)W_0P_{R^\perp}$ are polynomials in $(W_0,B,A)$, and one of them is nonzero at the witness. Viewed as a
polynomial in $(B,A)$, its coefficients are polynomials in $W_0$, not all identically zero; for $W_0$ outside their common
zero set it is nonzero on an open dense set of $(B,A)$ (\auxref{B.10}), which meets the open dense set where $B,A$ have full
rank and $W_0+sBA\in\mathcal U$. At such a point with $\ell=\mathbf 1$ the rank is $r(m+n-r)+\min(m,(m-r)(n-r))$. With the
upper bound, $d=\min(mn,r(m+n-r)+m)$ for $W_0$ outside a proper algebraic subset, and $d=r(m+n-r)+m$ iff $(m-r)(n-r)\ge m$.

\emph{Rank-one $W_0$.} Let $W_0=ac\T$ with all $a_i\ne0$ and $r<n$. At a point with $\ell$ invertible and $X\in\MM_r$,
$P_{U^\perp}\diag(\delta)ac\T P_{R^\perp}=(P_{U^\perp}(\delta\odot a))(P_{R^\perp}c)\T$. If $c\notin R$, this vanishes iff
$\delta\odot a\in U$, so $K=\diag(a)^{-1}U$ has dimension $r$ and the rank is $r(m+n-r)+m-r$. If $c\in R$, then $K=\R^m$ and
the rank is $r(m+n-r)$. The points with $\ell$ invertible, $X\in\MM_r$ and $c\notin R$ form an open dense subset of the
domain (since $r<n$), and the points of maximal rank form an open nonempty subset, so the two meet and
$d=r(m+n-r)+m-r$. This is at most $mn$, since $(m-r)(n-r-1)\ge0$. It is strictly below the generic value
$r(m+n-r)+\min(m,(m-r)(n-r))$ iff $m-r<(m-r)(n-r)$, that is iff $r<m$ and $r\le n-2$.

\emph{Covariance.} Let $g=P_\pi D$ be monomial, with $P_\pi$ a permutation matrix and $D\in T_m$, and let $h\in\GL_n$.
Then $gW_0h^{-1}$ has no zero row, and $g\,\diag(x)\,g^{-1}=\diag(P_\pi x)$, so $g\,\mathrm{Diag}_m\,g^{-1}=\mathrm{Diag}_m$;
also $g\Mr h^{-1}=\Mr$. Hence
\begin{equation*}
  g\,\big(\mathrm{Diag}_m\cdot(W_0+\Mr)\big)\,h^{-1}=\mathrm{Diag}_m\cdot(gW_0h^{-1}+\Mr)=\Im\mathrm{DoRA}_r(gW_0h^{-1}).
\end{equation*}
$\mathrm{Mon}_m\times\GL_n$ is a subgroup with this property, so it lies in the covariance group (proof of
\thmref{Definition}{II.14}).

\emph{The update at $B=0$.} HF initialises $\mu_i=\lVert w_{0,i}\rVert$ and then trains $\mu$; at $B=0$ the weight is
$\diag(\mu_i/\lVert w_{0,i}\rVert)W_0$, so $\Delta W=(\diag(\mu_i/\lVert w_{0,i}\rVert)-I)W_0$. With
$\mu_i=2\lVert w_{0,i}\rVert$ for every $i$, $\Delta W=W_0$, of rank $\rk W_0$, which exceeds $2r$ whenever $\rk W_0>2r$,
beyond the bound of \thmref{Theorem}{II.1}(c).
\end{proof}

% ======================================================================================================================
\subsection{Proposition II.9}\label{proof:II.9}

\begin{restate}{Proposition}{II.9}{the torus ladder}
\begin{enumerate}[label=(\alph*)]
  \item \textbf{HiRA.} $\rho(B,A)=W_0\odot(J+BA)$, where $J$ is the all-ones matrix; HF PEFT has $B_0=0$ and no scale
    factor on Linear and Conv layers (on embeddings it uses $A_0=0$). For every $W_0$,
    $\Im\mathrm{HiRA}_r=W_0+D_{W_0}(\Mr)$ with $D_{W_0}:X\mapsto W_0\odot X$. If $W_0$ has no zero entries,
    $D_{W_0}\in\GL(\Theta)$, so inside the slice at $\theta_0$ HiRA is LoRA transported by a diagonal automorphism: it has
    the same $d=r(m+n-r)$ (for $r\le\min(m,n)$) and the same gauge. In all cases $\rk\Delta W\le r\cdot\rk W_0$. The
    transport depends on $\theta_0$ (\thmref{Proposition}{II.13}).
  \item \textbf{Scalings.} Output-side scalings, $\diag(\ell)W_0=W_0\odot(J+(\ell-\mathbf 1)\mathbf 1\T)$, and
    input-side scalings, $W_0\diag(\ell)=W_0\odot(J+\mathbf 1(\ell-\mathbf 1)\T)$, have images in $\Im\mathrm{HiRA}_1$.
    Output-side examples: (IA)$^3$ on keys and values, the post-linear scale of SSF, DoRA's magnitude. Input-side
    examples: (IA)$^3$'s feed-forward vector, which scales the input of $W_{\mathrm{down}}$, and a LayerNorm gain
    $\gamma$ folded into the next linear map. The fold is exact when the norm's output is consumed only by linear maps:
    the input norms of pre-norm blocks, and a final norm before an untied head (in post-norm BERT, the last block's norm
    feeds only the head). It does not apply to Gemma-style post-sublayer norms, which HF's default LN-tuning targets for
    Gemma~2 and~3, to QK-norms, or to post-norm residual paths. Where it applies, $\gamma$ is tied across all consumers of
    the norm (Q, K and V, or gate and up) rather than being a free point per box, and with tied embeddings a fold into
    the head breaks the tie. LayerNorm biases and SSF's shift fold into the next bias, not into $W$. Two-sided scalings
    $\diag(a)W_0\diag(b)=W_0\odot(ab\T)$ lie in $\Im\mathrm{HiRA}_2$. As a morphism, $h(\ell)=(\ell-\mathbf 1,\mathbf
    1\T)$ lands in $\mathrm{HiRA}_1$ pointed at $(0,\mathbf 1\T)$, which is not HiRA's random-$A$ default pointing.
  \item \textbf{Orbit dimension.} The two-sided torus orbit $\{\diag(a)W_0\diag(b)\}$ has dimension $m+n-c(W_0)$, where
    $c(W_0)$ is the number of connected components of the bipartite support graph of $W_0$, an isolated vertex counting
    as its own component.
  \item \textbf{RoAd.} RoAd$_1$ (the paper's main variant, HF \texttt{road\_1}) is $W=RW_0$ with $R$ block-diagonal in
    scaled rotations $\alpha_i\mathrm{Rot}(\theta_i)$ of output pairs. It is the left action of
    $\CO^+(2)^{m/2}\cong(\mathbb C^\times)^{m/2}$: it is \textbf{(IA)$^{\mathbf 3}$ over $\mathbb C$}. It multiplies the
    input Gram matrix of each output pair by $\alpha_i^2$. Since $\alpha_i$ may be $0$, its image is the orbit closure,
    a real-linear space of dimension exactly $m$ whenever no output-pair block of $W_0$ is zero. RoAd$_2$ and RoAd$_4$
    (HF \texttt{road\_2}, \texttt{road\_4}) let each block range over all of $M_2(\R)$. They realise the block-diagonal
    action of the monoid $M_2(\R)^{m/2}$ (its group part $\GL_2^{m/2}$ is non-abelian), with linear image
    $\{\mathrm{blockdiag}(R_i)W_0\}$ of dimension $2m$ when every pair block has rank $2$. For them the Gram statement
    and ``over $\mathbb C$'' fail, and RoAd$_4$ carries a positive-dimensional gauge (two dimensions per row). HF pairs
    output $j$ with $j+g/2$ inside each group of size $g$ (default $64$) rather than adjacent outputs as in the paper,
    and also rotates the bias; the statements above do not depend on the pairing.
  \item \textbf{Zeros and covariance.} Every $W$ reachable by HiRA, (IA)$^3$ or SSF vanishes wherever $W_0$ does,
    $Z(W_0)\subseteq Z(W)$, and this persists under merging. New zeros can appear ($\ell_i=0$, or $(J+BA)_{ij}=0$), so
    the zero pattern itself is not invariant. The image covariance group of output-scaled boxes ((IA)$^3$ on $k,v$;
    post-linear SSF) is $\mathrm{Mon}_m\times\GL_n$, and that of input-scaled boxes ((IA)$^3$ on $W_{\mathrm{down}}$) is
    $\GL_m\times\mathrm{Mon}_n$. HiRA's is exactly $\mathrm{Mon}_m\times\mathrm{Mon}_n$ ($=\GL_1\times\GL_1$ when
    $m=n=1$).
\end{enumerate}
\end{restate}

\begin{proof}
\emph{(a)} $W_0\odot(J+BA)=W_0+W_0\odot(BA)$, and $\{BA\}=\Mr$, so $\Im\mathrm{HiRA}_r=W_0+D_{W_0}(\Mr)$. In the standard
basis $E_{ij}$ of $\Theta=\R^{m\times n}$, $D_{W_0}$ is diagonal with entries $(W_0)_{ij}$, so it is invertible iff no entry of
$W_0$ vanishes. Then, with $\tau(\theta)=W_0+D_{W_0}(\theta-W_0)$, HiRA is $\tau\circ\rho_{\mathrm{LoRA}}$ for LoRA with residual
$W_0$ and scale $1$. Since $\tau$ is an affine bijection, $D(\tau\circ\rho)=D_{W_0}\circ D\rho$ has the same rank as $D\rho$ at
every point, so $d$ is unchanged (\auxref{B.9}(d)), and $\tau\circ\rho\circ\varphi=\tau\circ\rho$ iff $\rho\circ\varphi=\rho$,
so the gauge group is unchanged; this is the transport $(D_{W_0})_*$ of \thmref{Proposition}{II.13}. For
$r\le\min(m,n)$, $d(\mathrm{LoRA}_r)=r(m+n-r)$ (\auxref{B.5}). For the rank, with columns $b_k$ of $B$ and rows $a_k$ of
$A$, $W_0\odot(BA)=\sum_{k\le r}\diag(b_k)\,W_0\,\diag(a_k)$ (\auxref{B.12}(c)), a sum of $r$ matrices of rank at most
$\rk W_0$. The map $D_{W_0}$ is built from $W_0=\theta_0$, which is the dependence on the base point.

\emph{(b)} Entrywise, $(\diag(\ell)W_0)_{ij}=\ell_i(W_0)_{ij}=(W_0)_{ij}\big(1+(\ell_i-1)\big)$, so
$\diag(\ell)W_0=W_0\odot(J+(\ell-\mathbf 1)\mathbf 1\T)$, and $(\ell-\mathbf 1)\mathbf 1\T=BA$ with $B=\ell-\mathbf 1$ and
$A=\mathbf 1\T$, a point of $\mathrm{HiRA}_1$. The input side is the transpose computation. For two-sided scalings,
$\diag(a)W_0\diag(b)=W_0\odot(ab\T)$ and $ab\T-J=[a,\,-\mathbf 1]\,[b,\,\mathbf 1]\T$ has rank at most $2$. The map
$h(\ell)=(\ell-\mathbf 1,\mathbf 1\T)$ is smooth, sends $\mathbf 1$ to $(0,\mathbf 1\T)$ and satisfies
$\rho_{\mathrm{HiRA}_1}\circ h(\ell)=\diag(\ell)W_0$, so it is a pointed morphism into $\mathrm{HiRA}_1$ pointed there.
For a norm whose output $\gamma\odot\hat x$ (with $\hat x$ the normalised input) feeds only linear maps $W$,
$W(\gamma\odot\hat x)=(W\diag(\gamma))\hat x$, so the gain is an input-side scaling of each consumer, the same $\gamma$ for
all of them; a bias $\beta$ gives $W\beta$, a change of the next bias. When the norm's output is added to the residual
stream (post-sublayer norms, post-norm residual paths) or is used inside attention scores (QK-norms), no linear map of the
adapted box absorbs $\gamma$. With tied embeddings the head is $E\T$, and replacing it by $E\T\diag(\gamma)$ unties it.

\emph{(c)} The orbit map $\Psi(a,b)=\diag(a)W_0\diag(b)$ of $T_m\times T_n$ has differential at $(\mathbf 1,\mathbf 1)$
\begin{equation*}
  (\alpha,\beta)\mapsto\diag(\alpha)W_0+W_0\diag(\beta)=\big((\alpha_i+\beta_j)(W_0)_{ij}\big)_{ij}.
\end{equation*}
Its kernel is cut out by $\alpha_i+\beta_j=0$ along every edge $(i,j)$ of the bipartite support graph. On a connected
component the value at one vertex determines all the others (with $\beta=-\alpha$ across each edge), and an isolated
vertex carries one free coordinate; so the kernel has dimension $c(W_0)$ and the rank is $m+n-c(W_0)$. Because
$\Psi(g g')=g\cdot\Psi(g')$ for the commutative group $T_m\times T_n$, the differential has this rank at every point of the
group, so the orbit, an immersed submanifold of dimension $\dim(T_m\times T_n)$ minus that of the stabiliser, has
dimension $m+n-c(W_0)$.

\emph{(d)} Identify an output pair, a $2\times n$ block $X$ with rows $x_1,x_2$, with the complex row $z=x_1+\mathrm ix_2\in
\mathbb C^n$. For $\mathrm{Rot}(\theta)=\left(\begin{smallmatrix}\cos\theta&-\sin\theta\\ \sin\theta&\cos\theta\end{smallmatrix}
\right)$ the rows of $\alpha\mathrm{Rot}(\theta)X$ are $\alpha(\cos\theta\,x_1-\sin\theta\,x_2)$ and
$\alpha(\sin\theta\,x_1+\cos\theta\,x_2)$, which correspond to $\alpha e^{\mathrm i\theta}z$. So RoAd$_1$ multiplies each pair
by a complex scalar, the left action of $\CO^+(2)^{m/2}\cong(\mathbb C^\times)^{m/2}$, the complex analogue of (IA)$^3$'s
real output scaling. The input Gram matrix of the pair becomes
$(\alpha\mathrm{Rot}X)\T(\alpha\mathrm{Rot}X)=\alpha^2X\T X$. As $(\alpha,\theta)$ ranges over $\R\times\R$, $\alpha e^{\mathrm i\theta}$
ranges over all of $\mathbb C$, so the image is $\prod_i\mathbb C\cdot z_i$, the closure of the orbit, a real-linear space of
dimension $2\,\#\{i:z_i\ne0\}$, equal to $m$ when no pair block vanishes. In RoAd$_2$ each row of a block has its own
polar pair, so each row of the $2\times2$ block is an arbitrary vector of $\R^2$ and the block is an arbitrary element of
$M_2(\R)$; RoAd$_4$ also reaches every element of $M_2(\R)$, with eight parameters for four entries, so its fibres have
dimension $4$ per block, two per row. The image $\{R_iX_i\}$ of one block is the set of $2\times n$ matrices whose rows lie
in $\operatorname{row}X_i$, of dimension $2\rk X_i$, so the total is $2m$ when every block has rank $2$. For $R=\diag(1,2)$
and $X=I_2$, $(RX)\T(RX)=\diag(1,4)$ is not a multiple of $X\T X$, and $R$ does not commute with $\mathrm{Rot}(\pi/2)$, so it
is not multiplication by a complex number. The pairing of outputs and the rotated bias do not enter these computations.

\emph{(e) Zeros.} $W_{ij}=(W_0)_{ij}(1+(BA)_{ij})$ for HiRA, $\ell_i(W_0)_{ij}$ or $(W_0)_{ij}\ell_j$ for (IA)$^3$, and
$\gamma_i(W_0)_{ij}$ for SSF's post-linear scale (its shift changes the bias), so $(W_0)_{ij}=0$ implies $W_{ij}=0$. After a
merge the new base $W_1$ has $Z(W_0)\subseteq Z(W_1)$, and induction gives the claim for any number of cycles. New zeros:
$\ell_i=0$ kills row $i$, and $B=-e_i$, $A=e_j\T$ give $(J+BA)_{ij}=0$.

\emph{(e) Output and input scalings.} Here $\Im M(W_0)=\mathrm{Diag}_m\cdot W_0$. For $g$ monomial and $h\in\GL_n$,
$g\,\diag(\ell)\,W_0h^{-1}=\diag(\ell')\,gW_0h^{-1}$ with $\diag(\ell')=g\diag(\ell)g^{-1}$, and $\ell'$ ranges over $\R^m$
with $\ell$, so $\mathrm{Mon}_m\times\GL_n$ is contained in the covariance group. Conversely, let $(g,h)$ be in it and take
$W_0=e_jc\T$ with $c\ne0$. Then $\Im M(e_jc\T)=\R\,e_jc\T$, and $g\,\Im M(W_0)h^{-1}=\R\,(ge_j)(h^{-\top}c)\T$ is a line,
while $\Im M(gW_0h^{-1})=\{y(h^{-\top}c)\T:\operatorname{supp}y\subseteq\operatorname{supp}(ge_j)\}$ has dimension
$\lvert\operatorname{supp}(ge_j)\rvert$. Equality for every $j$ forces every column of $g$ to have one nonzero entry, so $g$ is
monomial. The input side is the transpose.

\emph{(e) HiRA.} For permutation matrices $P,P'$ and diagonal $D,D'$ we have $P(X\odot Y)P'=(PXP')\odot(PYP')$ and
$D(X\odot Y)D'=(DXD')\odot Y$. For $g=PD$ and $h^{-1}=D'P'$ this gives
\begin{equation*}
  g\big(W_0\odot(J+Z)\big)h^{-1}=(gW_0h^{-1})\odot(J+PZP'),
\end{equation*}
and $P\Mr P'=\Mr$, so $\mathrm{Mon}_m\times\mathrm{Mon}_n$ is contained in the covariance group. Conversely, let $(g,h)$ be in
it and take $W_0=E_{ij}$, whose image $E_{ij}\odot(J+\Mr)=\R E_{ij}$ is a line. With $u=ge_i$ and $v=h^{-\top}e_j$,
covariance requires $\Im M(uv\T)=\R\,uv\T$. Suppose $uv\T$ has two nonzero entries $(k,l)\ne(k',l')$. Taking $Z=E_{kl}\in\MM_{\le1}$
gives $Y=uv\T+u_kv_lE_{kl}\in\Im M(uv\T)$. If $Y=\lambda uv\T$, the entry $(k',l')$ forces $\lambda=1$ and the entry $(k,l)$
forces $u_kv_l=0$, a contradiction. So $\lvert\operatorname{supp}u\rvert=\lvert\operatorname{supp}v\rvert=1$ for all $i,j$,
which makes $g$ and $h^{-\top}$, hence $h$, monomial. For $m=n=1$, $\mathrm{Mon}_1=\GL_1$.
\end{proof}

% ======================================================================================================================
\subsection{Theorem II.10}\label{proof:II.10}

\begin{restate}{Theorem}{II.10}{cut-rank bound; standard}
Let $\Gamma$ be a tensor network, evaluated in the symmetric monoidal category $(\FinVect,\otimes)$, with input legs $V$
and output legs $U$. Every wire and open leg is assumed incident to a vertex: insert an identity box on any bare strand,
or equivalently count each bare input--output wire as crossing every cut. A crossing hyperedge (copy spider) contributes
its dimension once. Partition the vertices into $\mathcal I\sqcup\mathcal O$. The weight $w$ of the cut is the product
of the dimensions of the internal wires between the two sides and of the open legs attached on the wrong side (input legs
at $\mathcal O$, output legs at $\mathcal I$). Then
\begin{equation*}
  \rk\llbracket\Gamma\rrbracket\le\min_{\mathrm{cuts}}w,
\end{equation*}
and rank is subadditive over formal sums. The bound refers to a chosen presentation of the network.
\end{restate}

\begin{proof}
Fix a cut $\mathcal I\sqcup\mathcal O$. Write $V_{\mathcal I},V_{\mathcal O}$ for the tensor products of the input legs
attached to $\mathcal I$ and to $\mathcal O$, $U_{\mathcal I},U_{\mathcal O}$ likewise for the output legs, and $C$ for the
tensor product of the spaces carried by the crossing wires, each crossing hyperedge counted once. Contracting every wire
internal to $\mathcal I$ produces a tensor whose free indices are those of $V_{\mathcal I}$, $U_{\mathcal I}$ and $C$; read
it as a linear map $T_{\mathcal I}:V_{\mathcal I}\to U_{\mathcal I}\otimes C$. Likewise the $\mathcal O$ side gives
$T_{\mathcal O}:C\otimes V_{\mathcal O}\to U_{\mathcal O}$. A crossing hyperedge is a single summation index shared by
tensors on both sides, so it is summed once, over a basis of its space, when the two sides are joined. Contracting the
crossing indices joins the two halves, and up to the symmetry isomorphisms that reorder tensor factors,
\begin{equation*}
  \llbracket\Gamma\rrbracket=(\id_{U_{\mathcal I}}\otimes T_{\mathcal O})\circ(T_{\mathcal I}\otimes\id_{V_{\mathcal O}}):
  V_{\mathcal I}\otimes V_{\mathcal O}\to U_{\mathcal I}\otimes C\otimes V_{\mathcal O}\to U_{\mathcal I}\otimes U_{\mathcal O}.
\end{equation*}
A composite factors through its middle space, so its rank is at most
$\dim U_{\mathcal I}\cdot\dim C\cdot\dim V_{\mathcal O}=w$. Only the symmetric monoidal structure was used: no cups or caps.
A bare strand from an input to an output, carried by an identity box, has one of its two legs on the wrong side of every
cut, which is the convention of the statement; without it, $I_d\otimes BA$ would violate the bound. Minimising over cuts
gives the bound, and subadditivity over formal sums is $\rk(X+Y)\le\rk X+\rk Y$ (\auxref{B.12}(a)). The weight depends on
the drawing, which is why the bound refers to a presentation; compare \auxref{B.14}.
\end{proof}

\noindent\emph{The bounds of Table~\ref{tab:cutbounds}.} Each follows from a cut in the stated presentation or from
subadditivity. The LoRA family $sBA$ factors through the bond of dimension $r$. A split initialisation gives
$\Delta W=sBA-sB_0A_0$, a sum of two such terms (\thmref{Theorem}{II.1}(c)). For HRA, the telescoping identity
\begin{equation*}
  R-I=\sum_{k=1}^r\big(H_k\cdots H_r-H_{k+1}\cdots H_r\big)=\sum_{k=1}^r(H_k-I)H_{k+1}\cdots H_r
\end{equation*}
writes $R-I$ as $r$ terms of rank at most $\rk(H_k-I)=1$, so $\rk(W_0R-W_0)\le r$; a Givens rotation differs from $I$ in a
$2\times2$ block, so the same identity gives $2g$ for $g$ gates. ReLoRA after $K$ phases is a sum of $K$ rank-$r$ terms, and
GraLoRA factors through $\R^{kr}$ (proof of \thmref{Proposition}{II.12}). LoHa's Hadamard product, drawn with copy spiders on
the two legs, factors through the face-splitting bond of dimension $r_1r_2$ (\auxref{B.12}(c)). HiRA's update
$W_0\odot BA=\sum_{k\le r}\diag(b_k)W_0\diag(a_k)$ factors through $r$ copies of a bond of dimension $\rk W_0$. LoKr gives
$\rk(C\otimes BA)=\rk C\cdot\rk(BA)\le\rk C\cdot r'$ (\auxref{B.12}(b)). FourierFT's update is a sum of $n_c$ real cosine
patterns $\cos\big(2\pi(uj/m+vl/n)\big)=\operatorname{Re}(x_uy_v\T)$ with $(x_u)_j=e^{2\pi\mathrm iuj/m}$ and
$(y_v)_l=e^{2\pi\mathrm ivl/n}$ (\auxref{B.16}(c)), and $\operatorname{Re}(xy\T)=\operatorname{Re}x\operatorname{Re}y\T-
\operatorname{Im}x\operatorname{Im}y\T$ has rank at most $2$. MoRA's reshape, C3A's circulant blocks and MPO adapters have
no narrow cut in their presentation, so the theorem bounds them only by $\min(m,n)$.

% ======================================================================================================================
\subsection{Theorem II.11}\label{proof:II.11}

\begin{restate}{Theorem}{II.11}{Kronecker incomparability; the linear algebra is standard, the PEFT reading is new}
Write $\R^m=\R^{m_1}\otimes\R^{m_2}$ and $\R^n=\R^{n_1}\otimes\R^{n_2}$, and view $\Delta\in\R^{m\times n}$ as a 4-leg
tensor. Let $\operatorname{rk}$ be the rank across the cut $\{U_1U_2\}\mid\{V_1V_2\}$ and $\operatorname{rk}_\otimes$
the rank across $\{U_1V_1\}\mid\{U_2V_2\}$. Put
\begin{equation*}
  \kappa_1:=\min(m_1,m_2)\cdot\min(n_1,n_2),\qquad\nu:=\min(m_1,n_1)\cdot\min(m_2,n_2).
\end{equation*}

(a) $\operatorname{rk}_\otimes\Delta\le k$ iff $\Delta=\sum_{i\le k}C_i\otimes D_i$. Every $\Delta$ has
$\operatorname{rk}_\otimes\Delta\le\min(m_1n_1,m_2n_2)$.

(b) $\operatorname{rk}(C\otimes D)=\operatorname{rk}C\cdot\operatorname{rk}D$, while
$\operatorname{rk}_\otimes(C\otimes D)=1$ for $C,D\ne0$. So $\nu$ is the largest rank of a single Kronecker product.

(c) If $u\in\R^m$ and $v\in\R^n$ have Schmidt ranks $\rho_u,\rho_v$, then
$\operatorname{rk}_\otimes(uv\T)=\rho_u\rho_v$. A generic rank-one update therefore has Kronecker rank $\kappa_1$. For
aligned splits, $(m_1-m_2)(n_1-n_2)\ge0$, which include equal square splits and every split HF PEFT's factorisation
produces (it puts the smaller factor first on both sides), $\kappa_1=\min(m_1n_1,m_2n_2)$ is the maximal Kronecker rank.

(d) \emph{Incomparability.} For $1\le r<\nu$ and $1\le k<\kappa_1$, the sets $\Mr$ and
$\mathcal K_{\le k}:=\{\Delta:\operatorname{rk}_\otimes\Delta\le k\}$ of unconstrained $k$-term Kronecker sums are
incomparable, for every split. A generic rank-one update lies in the first and not the second, and $C\otimes D$ with
full-rank factors, of Kronecker rank $1$ and rank $\nu$, lies in the second and not the first. For equal square splits
($m=n=N$, all factors $s$), $\kappa_1=\nu=s^2=N$, and $I_N$ is such a product.

(e) \emph{Constrained Kronecker methods.} Let $\Im M\subseteq W_0+\mathcal K_{\le k}$. Then
$\Im\mathrm{LoRA}_r\not\subseteq\Im M$ whenever $k<\kappa_1$, for every pointing of $\mathrm{LoRA}_r$. Give
$\mathrm{LoRA}_r$ its zero-init pointing, with image $W_0+\Mr$. Then $\Im M\not\subseteq\Im\mathrm{LoRA}_r$ iff some
$\theta\in\Im M$ has $\rk(\theta-W_0)>r$. For LoKr, $\Delta W=C\otimes BA$ with $C\in\R^{m_1\times n_1}$,
$B\in\R^{m_2\times r'}$ and $A\in\R^{r'\times n_2}$, put $r'_{\mathrm{eff}}=\min(r',m_2,n_2)$. The maximum rank is
$\min(m_1,n_1)\cdot r'_{\mathrm{eff}}$, attained. So $\Im\mathrm{LoKr}_{r'}\subseteq\Im\mathrm{LoRA}_r$ iff
$r\ge\min(m_1,n_1)\,r'_{\mathrm{eff}}$, and otherwise the two images are incomparable, provided $\kappa_1>1$. In HF PEFT
the second LoKr factor is low-rank only when its rank parameter satisfies $r'<\max(m_2,n_2)/2$; otherwise it is full,
and LoKr is unconstrained KronA, with maximum rank $\nu=\min(m_1,n_1)\min(m_2,n_2)$. On a $768$-wide layer
($24\otimes32$) the maximum rank therefore jumps from $360$ at $r'=15$ to $768$ at $r'=16$. With
\texttt{decompose\_both}, $C$ is also capped at rank $r'$.
\end{restate}

\begin{proof}
Index $\R^m$ by pairs $(i_1,i_2)$ and $\R^n$ by $(j_1,j_2)$. The rank across $\{U_1U_2\}\mid\{V_1V_2\}$ is the ordinary
rank of $\Delta$. The rank across $\{U_1V_1\}\mid\{U_2V_2\}$ is the rank of the rearrangement
$\mathcal R(\Delta)\in\R^{m_1n_1\times m_2n_2}$ with entries $\mathcal R(\Delta)_{(i_1j_1),(i_2j_2)}=\Delta_{(i_1i_2),(j_1j_2)}$,
a linear bijection that permutes entries, with $\mathcal R(C\otimes D)=\vect C\,\vect D\T$ (\auxref{B.13}(a)). (In this
proof and the next, $\mathcal R$ denotes this rearrangement.)

\emph{(a)} $\rk\mathcal R(\Delta)\le k$ iff $\mathcal R(\Delta)=\sum_{i\le k}x_iy_i\T$, iff
$\Delta=\sum_{i\le k}\mathcal R^{-1}(x_iy_i\T)=\sum_{i\le k}C_i\otimes D_i$ with $\vect C_i=x_i$ and $\vect D_i=y_i$. The
matrix $\mathcal R(\Delta)$ has size $m_1n_1\times m_2n_2$, which gives the second claim (\auxref{B.13}(b)).

\emph{(b)} $\rk(C\otimes D)=\rk C\cdot\rk D$ (\auxref{B.12}(b)), and $\mathcal R(C\otimes D)=\vect C\vect D\T$ has rank $1$
for $C,D\ne0$. The largest value of $\rk C\cdot\rk D$ is $\min(m_1,n_1)\min(m_2,n_2)=\nu$, attained by factors of full rank.

\emph{(c)} Take Schmidt decompositions $u=\sum_{i\le\rho_u}\sigma_ix_i\otimes y_i$ and
$v=\sum_{j\le\rho_v}\tau_jz_j\otimes w_j$ with positive coefficients and orthonormal families (\auxref{B.11}(a)). By the
mixed-product rule,
\begin{equation*}
  uv\T=\sum_{i,j}\sigma_i\tau_j\,(x_iz_j\T)\otimes(y_iw_j\T).
\end{equation*}
The two families $(x_iz_j\T)_{ij}$ and $(y_iw_j\T)_{ij}$ are Frobenius-orthonormal, because
$\langle x_iz_j\T,\,x_{i'}z_{j'}\T\rangle=(x_i\T x_{i'})\,(z_j\T z_{j'})$ and similarly for the second family, and the
$\rho_u\rho_v$ coefficients are positive. By
\auxref{B.13}(c) this is an operator-Schmidt decomposition, so $\rk_\otimes(uv\T)=\rho_u\rho_v$. The Schmidt rank of $u$ is the
rank of its $m_1\times m_2$ reshape, which equals $\min(m_1,m_2)$ outside the zero set of the maximal minors, a proper
algebraic subset (\auxref{B.10}(b)); likewise for $v$. So for generic $u,v$, $\rk_\otimes(uv\T)=\min(m_1,m_2)\min(n_1,n_2)=
\kappa_1$. If $(m_1-m_2)(n_1-n_2)\ge0$, then either $m_1\le m_2$ and $n_1\le n_2$, so $\kappa_1=m_1n_1=\min(m_1n_1,m_2n_2)$,
or $m_1\ge m_2$ and $n_1\ge n_2$, so $\kappa_1=m_2n_2=\min(m_1n_1,m_2n_2)$; in both cases $\kappa_1$ is the bound of (a).
Equal square splits are aligned, and so is any split with the smaller factor first on both sides.

\emph{(d)} A generic rank-one $uv\T$ lies in $\MM_{\le1}\subseteq\Mr$ and has $\rk_\otimes=\kappa_1>k$ by (c), so it is not in
$\mathcal K_{\le k}$. A product $C\otimes D$ with factors of full rank has $\rk_\otimes=1\le k$ and rank $\nu>r$ by (b). Neither
argument uses alignment. For equal square splits $\kappa_1=\nu=s\cdot s=N$ and $I_N=I_s\otimes I_s$.

\emph{(e)} Any pointing of $\mathrm{LoRA}_r$ has image $W_{\mathrm{res}}+\Mr$ for its frozen residual $W_{\mathrm{res}}$.
Suppose $W_{\mathrm{res}}+\Mr\subseteq W_0+\mathcal K_{\le k}$, and put $c=W_{\mathrm{res}}-W_0$, so $c+\Mr\subseteq\mathcal
K_{\le k}$. The set $\mathcal K_{\le k}$ is closed (the zero set of the $(k+1)\times(k+1)$ minors of $\mathcal R(\cdot)$) and
closed under scaling. For $X\in\Mr$ and $t>0$, $c+tX\in\mathcal K_{\le k}$, hence $c/t+X\in\mathcal K_{\le k}$, and letting
$t\to\infty$ gives $X\in\mathcal K_{\le k}$. So $\Mr\subseteq\mathcal K_{\le k}$, which contradicts (d) when $k<\kappa_1$.
The second clause restates $\Im\mathrm{LoRA}_r=W_0+\Mr$ for the zero-init pointing.

For LoKr, $\rk(C\otimes BA)=\rk C\cdot\rk(BA)\le\min(m_1,n_1)\cdot\min(r',m_2,n_2)$ (\auxref{B.12}(a),(b)), with equality
for $C$ of full rank and $BA$ of rank $r'_{\mathrm{eff}}$, which exists by \auxref{B.11}(a). If
$r\ge\min(m_1,n_1)r'_{\mathrm{eff}}$, every update has rank at most $r$ and $\Im\mathrm{LoKr}_{r'}\subseteq W_0+\Mr$.
Otherwise the maximising update has rank above $r$, so $\Im\mathrm{LoKr}_{r'}\not\subseteq\Im\mathrm{LoRA}_r$; and since
$\Im\mathrm{LoKr}_{r'}\subseteq W_0+\mathcal K_{\le1}$, the first clause with $k=1<\kappa_1$ gives
$\Im\mathrm{LoRA}_r\not\subseteq\Im\mathrm{LoKr}_{r'}$. When HF's second factor is full, LoKr is $C\otimes D$ with both
factors free, whose maximum rank is $\nu$ by (b). On a $768$-wide layer HF factors $768=24\cdot32$ on both sides, so
$m_1=n_1=24$ and $m_2=n_2=32$: for $r'=15<16$ the maximum rank is $24\cdot15=360$, and for $r'=16$ the second factor is full and
the maximum is $24\cdot32=768$. The \texttt{decompose\_both} option factors $C$ through rank $r'$, which caps $\rk C$ at $r'$.
The thresholds and factorisations are read from the HF PEFT source \citep{xmangrulkar2022peft}.
\end{proof}

% ======================================================================================================================
\subsection{Proposition II.12}\label{proof:II.12}

\begin{restate}{Proposition}{II.12}{gauge census: rank is a cut, dimension is a volume}
\begin{enumerate}[label=(\alph*)]
  \item $d(M)\le\lvert M\rvert-\dim(\text{generic gauge orbit})$, with equality when the gauge acts locally
    transitively on generic fibres.
  \item The values in Table~\ref{tab:census} hold for one $m\times n$ box, with $r\le\min(m,n)$ throughout except for
    VeRA (whose published RoBERTa-base setting has $r=1024>768$). In the LoHa row,
    \begin{equation*}
      F=r_1(m+n-r_1)+r_2(m+n-r_2)-(m+n-1)
    \end{equation*}
    is a proved upper bound for $d$. The gauge column lists the exhibited Lie group. It acts locally transitively on
    generic fibres away from the $mn$ cap, but not where the cap binds; for LoHa at $(m,n,r_1,r_2)=(4,4,2,2)$ the orbit
    has dimension $15$ and the fibre $16$. BitFit does not touch the $m\times n$ box; it is linear in the biases. Hybrid
    GraLoRA (\texttt{hybrid\_r}${}>0$) is not covered by the table. FourierFT is injective only when its sampled
    frequencies contain no conjugate pair $(u,v),(-u,-v)$; otherwise $d=\lvert M\rvert-\#\{\text{pairs}\}$. At LLM sizes
    this is close to $\lvert M\rvert$, but at $768\times768$ with $1000$ frequencies a pair occurs in about half of the
    draws (HF's default seed $777$ gives exactly two pairs there).
  \item \emph{Equal-budget trades.} Assume $2r\le\min(m,n)$, $r^2\le\min(m,n)$ and $2r^2<m+n-1$. LoHa$_{r,r}$ and
    LoRA$_{2r}$ both cost $2r(m+n)$, but LoHa reaches at least $(m+n-1)-2r^2$ \emph{fewer} dimensions (exactly that
    many when $F$ is attained [Num]). For $r\ge3$ it gains rank up to $r^2>2r$ in exchange. For $r\le2$, $r^2\le2r$, so
    $\Im\mathrm{LoHa}_{r,r}\subsetneq\Im\mathrm{LoRA}_{2r}$ and LoHa is strictly dominated. When $2r^2\ge m+n-1$ (still with $2r\le\min(m,n)$) the
    bound no longer separates them, and wherever $F$ is attained LoHa$_{r,r}$ matches or exceeds LoRA$_{2r}$ in
    dimension as well as in rank, for instance at $(m,n,r)=(9,9,3)$, at $r=64$ on $4096\times4096$, and at
    LyCORIS-style $r=32$ on $320$- to $640$-wide layers. GraLoRA$_k$ costs the same as LoRA$_r$, reaches the same
    number of dimensions, and has maximum rank $\min(kr,m,n)$, but caps every block at rank $r/k$. A generic rank-$r$
    update has rank $r$ on every block, so for $k\ge2$ and $kr\le\min(m,n)$ the images of LoRA$_r$ and GraLoRA$_k$ are
    incomparable.
\end{enumerate}
\end{restate}

\begin{proof}
\emph{(a)} By \auxref{B.9}(d), $d(M)$ is the maximal rank of $d\rho$, attained on an open set, which is dense when $\rho$ is
real-analytic on a connected domain (\auxref{B.10}). At a point $q$ of maximal rank the fibre through $q$ is, near $q$, a
submanifold of dimension $\lvert M\rvert-d(M)$ (\auxref{B.1}(b)), and it contains the gauge orbit of $q$, an immersed
submanifold. At generic $q$ this gives $\dim(\text{orbit})\le\lvert M\rvert-d(M)$. If the gauge acts locally transitively on
generic fibres, the orbit is open in the fibre and the two dimensions agree.

\emph{(b), LoRA.} $\GL_r$ acts on pairs of full rank by $(Bg^{-1},gA)$, freely, and its orbits are the fibres of
$(B,A)\mapsto BA$ over $\MM_r$ (\auxref{B.5}(b)), so $d=r(m+n)-r^2=r(m+n-r)$. The maximum rank is $r$.

\emph{GraLoRA.} With $k\mid m,n,r$, the update is a $k\times k$ grid of blocks of size $p\times q=\tfrac mk\times\tfrac nk$, each
an independent LoRA of rank $t=r/k$, so the image is $W_0$ plus the product of $k^2$ copies of
$\MM_{\le t}(p\times q)$, of dimension $k^2\,t(p+q-t)=r(m+n-r)$ since $t\le\min(p,q)$. The gauge $\prod\GL_t$ acts freely on
generic parameters with orbits of total dimension $k^2t^2=r^2$, equal to the fibre dimension. For the rank, write block
$(i,j)$ as $U_{ij}V_{ij}\T$; then $\Delta W=\operatorname{blockdiag}(L_1,\dots,L_k)\,S$ with
$L_i=[U_{i1}\cdots U_{ik}]\in\R^{p\times r}$ and $S\in\R^{kr\times n}$ stacking the $V_{ij}\T$ placed in column block $j$, so
$\rk\Delta W\le\min(kr,m,n)$. (Placing rank-$t$ blocks only on the diagonal would give rank $r$, not $kr$.) To see that
$\min(kr,m,n)$ is attained, assume $p\le q$ (otherwise transpose) and put $a=\min(kt,p)$, so that
$\min(kr,m,n)=ka$. Write $a=k\alpha+\beta$ with $0\le\beta<k$, and let $N_{ij}=\alpha+[\,(j-i)\bmod k<\beta\,]$. Then every
row and column sum of $N$ is $a\le p\le q$, and $N_{ij}\le t$, because $\alpha\le t$ with $\alpha=t$ only when $\beta=0$. In
block row $i$ pick $N_{ij}$ distinct rows for each $j$ (in all $a\le p$ rows), in block column $j$ pick $N_{ij}$ distinct
columns for each $i$ (in all $a\le q$ columns), and put a $1$ at the matched pairs. The resulting $0/1$ matrix has at most
one $1$ in each row and column, so its rank is $\sum_{ij}N_{ij}=ka$, and its block $(i,j)$ has rank $N_{ij}\le t$, so it lies
in the image. Rank at least $ka$ is the nonvanishing of a minor, a polynomial in the block factors, so it holds for generic
parameters (\auxref{B.10}).

\emph{LoHa.} The update is $(B_1A_1)\odot(B_2A_2)$ with budget $(r_1+r_2)(m+n)$. The torus acts by
$(B_1,A_1,B_2,A_2)\mapsto(D_1B_1,A_1D_2,D_1^{-1}B_2,A_2D_2^{-1})$ and preserves the update, because
$(D_1XD_2)\odot(D_1^{-1}YD_2^{-1})=X\odot Y$ for diagonal $D_1,D_2$; together with $\GL_{r_1}\times\GL_{r_2}$ acting on the two
LoRA factors this gives an action of $G=\GL_{r_1}\times\GL_{r_2}\times T_m\times T_n$. We compute the stabiliser at generic
parameters. If $m=n=r_1$, then $r_1(m+n-r_1)=mn$ and $F\ge mn$ (since $r_2(m+n-r_2)-(m+n-1)=(r_2-1)(m+n-r_2-1)\ge0$), so the
bound $d\le\min(mn,F)$ is trivial; otherwise, by symmetry under transposition, assume $m\ge r_1+1$. Let
$(g_1,g_2,D_1,D_2)$ fix the point. From $D_1B_1g_1^{-1}=B_1$, every row of $B_1$ is a left eigenvector of $g_1$; for generic
$B_1$ any $r_1$ of its rows form a basis and the remaining rows have all coordinates nonzero in that basis, which forces a
single eigenvalue, so $g_1=cI$ and $D_1=cI$. From $g_1A_1D_2=A_1$ and nonzero columns of $A_1$, $D_2=c^{-1}I$. From
$D_1^{-1}B_2g_2^{-1}=B_2$ and $B_2$ of full column rank ($m\ge r_2$), $g_2=c^{-1}I$, and then the last equation holds. So the
stabiliser is $\{(cI,c^{-1}I,cI,c^{-1}I)\}$, of dimension $1$, and the generic orbit has dimension
$r_1^2+r_2^2+m+n-1$. By (a), $d\le(r_1+r_2)(m+n)-r_1^2-r_2^2-(m+n-1)=F$, and trivially $d\le mn$. Where $d=F$ the orbit
has the dimension of the fibre and the action is locally transitive; at $(4,4,2,2)$ the orbit has dimension $4+4+7=15$ while
the [Num] value $d=16$ gives a fibre of dimension $32-16=16$. For the rank, $\rk(X\odot Y)\le r_1r_2$ (\auxref{B.12}(c)) and
$\le\min(m,n)$. To attain $k=\min(r_1r_2,m,n)$, choose distinct pairs $(a_i,b_i)\in[r_1]\times[r_2]$ for $i\le k$, let the
first $k$ rows of $B_1$ be $e_{a_i}\T$ and the first $k$ columns of $A_1$ be $e_{a_i}$, likewise $B_2,A_2$ with the $b_i$, and
set the rest to zero. Then $(B_1A_1)_{ij}=[a_i=a_j]$ and $(B_2A_2)_{ij}=[b_i=b_j]$ for $i,j\le k$, so the Hadamard product is
$I_k$ in the top-left corner and zero elsewhere, of rank $k$; by \auxref{B.10} this rank is generic.

\emph{HRA.} The proof of \thmref{Theorem}{II.7}(a) computes $d=rn-r(r+1)/2$ and the fibre dimension $r(r+1)/2$ for every
$1\le r\le n$, with no use of parity or of $r\le n-2$; for $r\in\{n-1,n\}$ the value is $\dim\SO(n)$. The scalings
$(\R^\times)^r$ are the exhibited Lie gauge, and $W_0H_{u_1}\cdots H_{u_r}-W_0$ has rank at most $r$. Evenness of $r$ is used
only by HF's paired initialisation, which needs pairs to start at $W_0$.

\emph{DoRA.} The value of $d$ is \thmref{Proposition}{II.8}. $\GL_r$ acts on $(B,A)$ by $(Bg^{-1},gA)$, fixing $\mu$, freely on
pairs of full rank and preserving $\rho$. For generic $W_0$ with $(m-r)(n-r)\ge m$ the fibre has dimension
$r(m+n)+m-r(m+n-r)-m=r^2$, so $\GL_r$ is locally transitive on generic fibres. Otherwise $d=mn$ and the fibre has dimension
$r(m+n)+m-mn=r^2+m-(m-r)(n-r)>r^2$, which is $6$ at $(3,3,2)$; for rank-one $W_0$, $d=r(m+n-r)+m-r$ and the fibre has dimension
$r^2+r$. The update can have rank up to $\min(m,n)$ (\thmref{Proposition}{II.8}, last item).

\emph{VeRA.} The update is $\diag(b)\,B\diag(d)\,A$ with frozen random $B\in\R^{m\times r}$, $A\in\R^{r\times n}$ and trainable
$b\in\R^m$, $d\in\R^r$, so $\lvert M\rvert=m+r$, and $(b,d)\mapsto(cb,c^{-1}d)$ preserves it. The value $d(M)=m+r-1$ is
checked numerically [Num]. The rank is at most $\rk(B\diag(d)A)\le\min(r,m,n)$, and for invertible $\diag(b)$, $\diag(d)$ it
equals $\rk(BA')$ with $A'=\diag(d)A$, which is $\min(r,m,n)$ for generic frozen matrices (\auxref{B.10}). None of this uses
$r\le\min(m,n)$.

\emph{KronA.} The update $A\otimes B$ satisfies $\mathcal R(A\otimes B)=\vect A\vect B\T$, so the image is $W_0$ plus the
preimage under the linear bijection $\mathcal R$ (\auxref{B.13}(a)) of the Segre cone $\MM_{\le1}(m_1n_1\times m_2n_2)$, of
dimension $m_1n_1+m_2n_2-1$ (\auxref{B.8}). If $A\otimes B=A'\otimes B'\ne0$, then $\vect A\vect B\T=\vect A'\vect B'^{\top}$,
so $A'=cA$ and $B'=c^{-1}B$; the fibre over a nonzero update is the orbit of $\R^\times$. The maximum rank is $\nu$ by
\thmref{Theorem}{II.11}(b).

\emph{Tensor-train and Tucker adapters.} The update is a fixed reshape, a linear bijection, of a tensor in TT or Tucker
format, so $d$ is the dimension of the set of tensors of the given admissible minimal ranks, which is the parameter count
minus $\sum_er_e^2$, with the bond gauge $\prod_e\GL_{r_e}$ acting on the cores (\auxref{B.15}).

\emph{Linear methods.} For $\rho(q)=\theta_0+Lq$ with $L$ linear, $d=\rk L$, the fibres are the cosets $q+\ker L$, and
translation by $\ker L$ acts transitively on each. LoRA-FA, LoRA-XS, FourierFT, C3A and masks are of this form; so is MiSS
in each HF mode, including the BAT mode, whose update $\operatorname{blocks}(W_0)M+M$ is linear in the trained block $M$
(\thmref{Proposition}{II.13}).

\emph{FourierFT.} HF's update is $c\operatorname{Re}(\operatorname{ifft2}(S))$ with $S$ supported on a sampled set $\Omega$
of $\lvert M\rvert$ frequencies. By \auxref{B.16}(c), frequency $p=(u,v)$ contributes $s_p\kappa_p$ with
$\kappa_p(j,l)=\tfrac1{mn}\cos\big(2\pi(uj/m+vl/n)\big)$, and $\kappa_p=\kappa_{-p}$. Write $\chi_p(j,l)=e^{2\pi\mathrm
i(uj/m+vl/n)}$, so that $\langle\chi_p,\chi_{p'}\rangle=mn\,[p=p']$ and $\cos(\cdots)=\tfrac12(\chi_p+\chi_{-p})$. Hence
$\langle\kappa_p,\kappa_{p'}\rangle=0$ unless $p'=\pm p$, and every $\kappa_p\ne0$. So the patterns of distinct classes
$\{p,-p\}$ are linearly independent, and the rank of $s\mapsto\sum_ps_p\kappa_p$ is the number of classes met by $\Omega$, which
is $\lvert\Omega\rvert$ minus the number of conjugate pairs inside $\Omega$. Each pair $\{p,-p\}\subseteq\Omega$ contributes
the kernel direction $s_p=1$, $s_{-p}=-1$. The frequencies of pairs in random draws and for HF's default seed are numerical
observations.

\emph{(c)} Assume $2r\le\min(m,n)$, $r^2\le\min(m,n)$ and $2r^2<m+n-1$. Both methods cost $2r(m+n)$, and
$d(\mathrm{LoRA}_{2r})=2r(m+n-2r)$. By (b), $d(\mathrm{LoHa}_{r,r})\le F=2r(m+n-r)-(m+n-1)$, and
\begin{equation*}
  d(\mathrm{LoRA}_{2r})-F=(m+n-1)-2r^2>0 .
\end{equation*}
So LoHa reaches at least $(m+n-1)-2r^2$ fewer dimensions; since $F<d(\mathrm{LoRA}_{2r})\le mn$, the cap does not bind, and
the gap is exactly $(m+n-1)-2r^2$ when $d(\mathrm{LoHa}_{r,r})=F$. Its maximum rank is $\min(r^2,m,n)=r^2$ by (b), against
$2r$ for $\mathrm{LoRA}_{2r}$, and $r^2>2r$ iff $r\ge3$. For $r\le2$, every LoHa update has rank at most $r^2\le2r$, so
$\Im\mathrm{LoHa}_{r,r}\subseteq W_0+\MM_{\le2r}=\Im\mathrm{LoRA}_{2r}$ (both pointed at $W_0$, as HF zero-initialises one
LoHa factor), and the inclusion is strict because the dimensions differ. When $2r^2\ge m+n-1$, $F\ge d(\mathrm{LoRA}_{2r})$ and
the bound does not separate them; wherever $d(\mathrm{LoHa}_{r,r})=F$, LoHa reaches at least as many dimensions, and its
maximum rank $\min(r^2,m,n)$ is at least $\min(2r,m,n)=2r$, because $r\ge2$ here (for $r=1$, $2\ge m+n-1$ contradicts
$m,n\ge2$). The examples satisfy $2r^2\ge m+n-1$: $18\ge17$ at $(9,9,3)$, $8192\ge8191$ at $r=64$ on $4096\times4096$, and
$2048\ge m+n-1$ for $r=32$ and $m,n\le640$.

GraLoRA$_k$ has budget $k^2\cdot t(p+q)=r(m+n)$, dimension $r(m+n-r)$ and maximum rank $\min(kr,m,n)$ by (b). Let $k\ge2$ and
$kr\le\min(m,n)$. A rank-$r$ update $BA$ has blocks $B_iA_j$, with $B_i$ the rows of $B$ in block row $i$ and $A_j$ the
columns of $A$ in block column $j$; for generic $B,A$ every $B_i$ ($p\times r$, $r\le p$) and $A_j$ ($r\times q$, $r\le q$) has
rank $r$, so every block has rank $r>r/k$, and the update is not in GraLoRA's image. Conversely GraLoRA reaches rank $kr>r$,
outside $W_0+\Mr$. Both images contain $W_0$, so they are incomparable.
\end{proof}

% ======================================================================================================================
\subsection{Proposition II.13}\label{proof:II.13}

\begin{restate}{Proposition}{II.13}{transport of structure}
For $T\in\GL(\Theta)$ put $\tau_T(\theta)=\theta_0+T(\theta-\theta_0)$ and $T_*M:=(Q,q_0,\tau_T\circ\rho)$.
\begin{itemize}
  \item $T_*$ is a strict automorphism of $\Meth(\Theta,\theta_0)$, with inverse $(T^{-1})_*$. It maps $S_1$ to $T\,S_1$
    and preserves $d$, $g$, the gauge group, the first-order deficit, polynomial degree, the regular or apex type and
    the germ type of the image. It does not preserve rank, and, for the family extension below, it does not preserve the
    covariance group.
  \item $K^{T_*M}=T\,K^M\,T\T$ for every $T$. For any update rule that queries the loss only through
    $q\mapsto L(\rho(q))$, with its own state on $Q$, $T_*M$ trained on $L$ and $M$ trained on $L\circ\tau_T$ have
    identical $q$-trajectories, so their $\theta$-trajectories are conjugate by $\tau_T$. Rules that read
    $\nabla_\theta L$ and solve a least-squares problem on $\Theta$ (LoRA-Pro), and $\theta$-level weight decay or
    clipping, are covered only for orthogonal $T$. Orthogonality of $T$ is exactly the condition for $\tau_T$ to be a
    Frobenius isometry; for rules of the first kind it is needed only to keep the Frobenius spectrum of $K$, or an
    isometry-invariant loss class such as $\lVert\theta-\theta^\star\rVert^2$, unchanged.
  \item Coordinatewise optimizers (SGD, Adam) are also equivariant under coordinate permutations of $Q$. The Kronecker
    identification below needs this, because it composes $\mathcal R^{-1}_*$ with the vec reshape of $Q$.
  \item \emph{Families and rebasing.} For $T$ independent of the base point, define the family extension
    $(T_*M)(\theta):=(Q,q_0,\theta+T(\rho_\theta-\theta))$. For an additive family with fixed shape $C$
    (\thmref{Definition}{I.2}), $T_*M$ is additive with shape $TC$, so transport commutes with rebasing:
    $R_K(T_*M)=\tau_T(R_K(M))$ (\thmref{Theorem}{V.3}(a)). For action-type families this can fail even for $T=2I$,
    because the transported family acts by $I+T(g-I)$. For OFT on $\Theta=\R^{1\times2}$, two cycles of
    $\theta\mapsto\theta(2R-I)$ with $R=\operatorname{rot}(\pi)$ reach $9\theta_0$, outside the transported circle.
    Transports that depend on $\theta_0$ may fail (HiRA below) or may not (LoRA transported by the SVD frame of
    $\theta_0$, since $\Mr$ is $\Ort(m)\times\Ort(n)$-invariant). The covariance group (\thmref{Definition}{II.14}) is a
    property of a family and is compared through this extension.
\end{itemize}
Instances.
\begin{itemize}
  \item \emph{Kronecker family.} An isomorphism of slices over $\Theta'=\R^{m_1n_1\times m_2n_2}$ and over $\Theta$:
    $\sum_{i\le k}C_i\otimes D_i=\mathcal R^{-1}_*\mathrm{LoRA}_k$, with $Q$ reshaped by vec, dynamics included for
    coordinatewise optimizers. KronA with unconstrained factors is $\mathcal R^{-1}_*\mathrm{LoRA}_1$. LoKr is
    $\mathcal R^{-1}_*$ of $\mathrm{LoRA}_1$ with the second factor constrained to $\vect\MM_{\le r'}$ when
    $r'<\max(m_2,n_2)/2$ (HF's switch); otherwise HF's second factor is full and LoKr is KronA. LoKr's
    \texttt{decompose\_both} option also constrains the first factor.
  \item \emph{Linear methods as coordinate subspaces in transported frames.} FourierFT: sparse in the real cosine half of
    the real Fourier frame. With scaling $c$ (default $150$), HF's update $\Delta W=c\operatorname{Re}(\operatorname{ifft2}(S))$
    is the transport of a mask, with $g=0$, only when the sampled index set has no conjugate pair $(u,v),(-u,-v)$;
    otherwise $g$ is the number of such pairs. Its image lies in the point-symmetric subspace
    $\Delta W_{j,l}=\Delta W_{-j,-l}$, and the frame is orthogonal only up to scale (self-conjugate frequencies have twice
    the squared norm). WaveFT: sparse in a wavelet frame, when HF's pad--IDWT--crop map is invertible (plausible for Haar
    on even shapes; unverified for other wavelet families). C3A: block-circulant; over $\R$ the frame is given by the
    powers $P^k$ of the cyclic shift (Gram matrix $bI$); over $\mathbb C$, $\operatorname{circ}(w)=F^{-1}\diag(Fw)F$ for
    first column $w$, and the real form has $2\times2$ rotation--scaling blocks; HF adds a factor
    $1/\texttt{in\_features}$. LoRA-XS, SVFT, TinyLoRA: masks or random subspaces in the SVD frame of $W_0$, a
    $\theta_0$-dependent transport. SHiRA, FISH, Trainable Tokens: masks in the standard frame. MiSS/Bone in its default
    (balance) mode: $\Delta W=D(\mathbf 1_g\T\otimes I_r)$, which is LoRA-FA with a frozen \emph{summing} projection; its
    BAT mode, $\Delta W=\operatorname{blocks}(W_0)\,M+M$ blockwise, is linear in the trained block $M$ but depends on
    $\theta_0$, so it is not additive. These frame readings of FourierFT and C3A are for \texttt{init\_weights=True}; HF's
    defaults start off $\theta_0$ (FourierFT draws its spectrum from $\mathcal N(0,1)$ with scaling $150$, so at $768^2$
    with $1000$ frequencies $\lVert\Delta W_0\rVert_F\approx4.3$, about $14\%$ of a std-$0.04$ $W_0$; C3A is
    Xavier-uniform), and these are objects with a pointing defect.
  \item \emph{HiRA} $=(D_{W_0})_*\mathrm{LoRA}$ inside the slice at $\theta_0$, when $W_0$ has no zero entries, with
    $D_{W_0}:X\mapsto W_0\odot X$. Because $D_{W_0}$ depends on $\theta_0$, re-based HiRA is not transported ReLoRA. Two
    cycles reach $W_0\odot(J+Z_1)\odot(J+Z_2)$, whose relative update $Z_1+Z_2+Z_1\odot Z_2$ has rank up to
    $\min(2r+r^2,m,n)$, attained, where transported ReLoRA stays at rank $\le2r$.
\end{itemize}
\end{restate}

\begin{proof}
\emph{The automorphism.} $\tau_T$ is an affine bijection fixing $\theta_0$ with linear part $T$, so $T_*M$ is pointed. If
$h:M\to N$ satisfies $\rho_N\circ h=\rho_M$, then $\tau_T\circ\rho_N\circ h=\tau_T\circ\rho_M$, so $h$ is also an arrow
$T_*M\to T_*N$; let $T_*$ act as the identity on arrows. This is a functor, and $(T^{-1})_*T_*=\id$ on the nose because
$\tau_{T^{-1}}\circ\tau_T=\id$. So $T_*$ is a strict automorphism.

\emph{Invariants.} $D(\tau_T\circ\rho)_q=T\,D\rho_q$ for every $q$. So $S_1(T_*M)=T\,S_1(M)$, and the rank of the differential
is unchanged at every point; hence $d$ (\auxref{B.9}(d)), $\dim S_1$, the deficit, $g=\lvert M\rvert-d$, and the regular or
apex type are unchanged. Since $\tau_T$ is injective, $\tau_T\circ\rho\circ\varphi=\tau_T\circ\rho$ iff $\rho\circ\varphi=\rho$,
so the gauge groups coincide. Composition with an affine map does not raise polynomial degree, and $\tau_{T^{-1}}$ undoes
it, so the degree is unchanged. Finally $\tau_T$ is a diffeomorphism of $\Theta$ fixing $\theta_0$ that carries $\Im M$ onto
$\Im T_*M$, so the germs at $\theta_0$ are diffeomorphic. Rank is not preserved: any $T$ with $T(E_{11})=I_n$ ($n\ge2$, $m=n$)
sends a rank-one update to a full-rank one.

\emph{Covariance is not preserved.} Take the family extension of LoRA$_1$ by $T=D_V:X\mapsto V\odot X$, with $V$ fixed,
without zero entries and of rank at least $2$; its image at $W$ is $W+V\odot\MM_{\le1}$. If $(g,h)$ were in its covariance
group for every pair, comparing the images at $gWh^{-1}$ and $W$ would give $V\odot\MM_{\le1}=g\,(V\odot\MM_{\le1})\,h^{-1}$, so
$V\odot\MM_{\le1}$ would be a union of $\GL_m\times\GL_n$-orbits, that is of rank strata $\MM_k$. It has dimension $m+n-1$,
since $D_V$ is invertible, while $\dim\MM_k=k(m+n-k)>m+n-1$ for $2\le k\le\min(m,n)$. So it would lie in $\MM_{\le1}$, and in
particular $V\odot J=V$ would have rank at most one, a contradiction. Hence the covariance group of the transported family is
a proper subgroup of $\GL_m\times\GL_n$, the covariance group of LoRA.

\emph{Cometric and dynamics.} $K^{T_*M}=D(\tau_T\rho)\Lambda D(\tau_T\rho)\T=T\,D\rho\Lambda D\rho\T\,T\T=TK^MT\T$. The method $T_*M$
trained on $L$ minimises $L\circ\tau_T\circ\rho$ on $Q$, and $M$ trained on $L\circ\tau_T$ minimises $(L\circ\tau_T)\circ\rho$,
the same function. A rule that sees the loss only through this function, from the same initial state (and, for stochastic
rules, the same random stream), therefore produces the same $q$-trajectory, and the weights are $\tau_T(\rho(q(t)))$ against
$\rho(q(t))$. A rule that reads $\nabla_\theta L$ and takes the Frobenius least-squares step on $\Theta$ moves $q$ by
$-(D\rho\T T\T TD\rho)^+D\rho\T T\T\nabla L$ under $T_*M$, and by $-(D\rho\T D\rho)^+D\rho\T T\T\nabla L$ under $M$ with loss
$L\circ\tau_T$, since $\nabla_\theta(L\circ\tau_T)=T\T\nabla L$. They agree when $T\T T=I$ and differ in general: for
$\Theta=\R^2$, $\rho(q)=qe_1$ and $T=\diag(2,1)$ the steps are $-\tfrac12\partial_1L$ and $-2\partial_1L$. Weight decay and
clipping on $\theta$ use the Frobenius norm, which $\tau_T$ preserves iff $T$ is orthogonal. For orthogonal $T$,
$TKT\T$ is similar to $K$, so the Frobenius spectrum of the cometric is kept, and $\lVert\tau_T(\theta)-\tau_T(\theta^\star)
\rVert=\lVert\theta-\theta^\star\rVert$.

\emph{Permutations of $Q$.} A coordinatewise rule updates each coordinate from that coordinate's own gradient history with
shared hyperparameters. If $q(t)$ is its trajectory for an objective $f$ from $q_0$, then $\pi q(t)$ is its trajectory for
$f\circ\pi^{-1}$ from $\pi q_0$, for any coordinate permutation $\pi$.

\emph{Families and rebasing.} For an additive family, $\rho_\theta=\theta+\delta$ with $\delta$ independent of $\theta$ and
shape $C=\delta(Q)$, the extension is $\theta+T\delta$, additive with shape $TC$. After $K$ merge-and-restart cycles an
additive family reaches $R_K=\theta_0+C+\dots+C$ ($K$ summands), and since $T$ is linear,
$\theta_0+(TC)+\dots+(TC)=\theta_0+T(C+\dots+C)=\tau_T(R_K)$. For OFT on $\R^{1\times2}$ the family is
$\{\theta R:R\in\SO(2)\}$ and its extension by $T=2I$ is $\theta\mapsto\theta+2(\theta R-\theta)=\theta(2R-I)$. With
$R=\operatorname{rot}(\pi)=-I$, two cycles give $\theta_0(-3I)(-3I)=9\theta_0$. Every rebasing of OFT stays in
$\theta_0\cdot\SO(2)$, whose transport $\{\theta_0(2S-I):S\in\SO(2)\}$ consists of vectors of norm at most $3\lVert\theta_0\rVert$,
so $9\theta_0$ lies outside it when $\theta_0\ne0$. For LoRA transported by $X\mapsto U(\theta)XV(\theta)\T$ with orthogonal
$U(\theta),V(\theta)$, the transported shape is $U\Mr V\T=\Mr$, so the transported family equals LoRA's, and both sides of
the rebasing identity are $\theta_0+\MM_{\le\min(Kr,m,n)}$ (by the splitting argument of \thmref{Theorem}{II.1}(c)).

\emph{Kronecker family.} By \auxref{B.13}(a), $\mathcal R:\Theta\to\Theta'$ is a linear bijection with
\begin{equation*}
  \mathcal R\Big(\sum_{i\le k}C_i\otimes D_i\Big)=\sum_i\vect C_i\,\vect D_i\T=BA,
\end{equation*}
where $B$ has columns $\vect C_1,\dots,\vect C_k$ and $A$ has rows $\vect D_1\T,\dots,\vect D_k\T$. So the $k$-term Kronecker method is $\mathcal R^{-1}\circ\rho_{\mathrm{LoRA}_k}$ up to the
vec relabelling of $Q$, and the argument for $T_*$ applies verbatim to an isomorphism between two weight spaces. Its
dynamics agree with LoRA$_k$'s on the pulled-back loss for rules that query the loss through $\rho$, and the vec relabelling is a
coordinate permutation, harmless for coordinatewise rules. KronA is $k=1$. For LoKr,
$\mathcal R(C\otimes BA)=\vect C\vect(BA)\T$, LoRA$_1$ with second factor in $\vect\MM_{\le r'}$, when HF factors the second
factor; otherwise it is KronA. With \texttt{decompose\_both} the first factor is factored too.

\emph{FourierFT.} By the computation in the proof of \thmref{Proposition}{II.12}, the patterns $\kappa_p$ of distinct
classes $\{p,-p\}$ are orthogonal; with the sine patterns $\sin(2\pi(uj/m+vl/n))$ of the classes with $p\ne-p$ they form an
orthogonal basis of $\R^{m\times n}$. Let $T$ send the standard basis to this basis (cosines on the positions of the
coefficients, sines elsewhere). If $\Omega$ has no conjugate pair, FourierFT is $T_*$ of the mask on $\Omega$, with $g=0$;
otherwise the two coefficients of a pair hit the same pattern and $g$ is the number of pairs. Each $\kappa_p$ is even under
$(j,l)\mapsto(-j,-l)$ (indices modulo $m,n$), so the image lies in the point-symmetric subspace. A self-conjugate
frequency has $\kappa_p=\tfrac1{mn}\chi_p$ with $\chi_p$ real, of squared norm $\tfrac1{mn}$, while for $p\ne-p$ the squared
norm is $\tfrac1{2mn}$. HF's default draws $1000$ coefficients from $\mathcal N(0,1)$ and multiplies by $150$; for generic
frequencies $\mathbb E\lVert\Delta W_0\rVert_F^2=150^2\cdot1000/(2\cdot768^2)\approx19.1$, a root-mean-square norm of
$4.37$. A single draw deviates from it by about $2\%$, since the sum of $1000$ squared standard normals has relative
standard deviation $\sqrt{2/1000}$, which gives the value $\lVert\Delta W_0\rVert_F\approx4.3$ quoted. Against
$\lVert W_0\rVert_F\approx0.04\cdot768\approx30.7$ this is about $14\%$.

\emph{Other frames.} If HF's pad--IDWT--crop map $P$ is invertible, WaveFT is $P_*$ of a mask by definition. For C3A, the
powers $P^0,\dots,P^{b-1}$ of the cyclic shift span the $b\times b$ circulants and have Gram matrix $bI$, and
$\operatorname{circ}(w)=F^{-1}\diag(Fw)F$; the real form follows from conjugate eigenvalues (\auxref{B.16}(a),(b)). So a
block-circulant update is a mask in the frame $\{E^{(\mathrm{blk})}_{ij}\otimes P^k\}$. LoRA-XS's update
$U_rR\Sigma_rV_r\T=\sum_{i,j\le r}R_{ij}\sigma_ju_iv_j\T$ is a scaled mask in the orthonormal frame $(u_iv_j\T)$ built from the SVD
of $W_0$; SVFT's update $U M V\T$ with sparse $M$ is a mask there too, and TinyLoRA's is a random subspace of it. SHiRA,
FISH and Trainable Tokens update a fixed set of entries. In MiSS's balance mode HF computes the update on an input $x$ as
$D$ times the sum of the $g$ chunks of $x$ of size $r$, that is $D(\mathbf 1_g\T\otimes I_r)x$; this is LoRA-FA with the frozen
projection $\mathbf 1_g\T\otimes I_r$. Its BAT mode is linear in $M$ with coefficients from $W_0$, so it is linear but not additive.

\emph{HiRA.} By \thmref{Proposition}{II.9}(a), HiRA is $(D_{W_0})_*\mathrm{LoRA}$ when $W_0$ has no zero entries. A merge
gives $W_1=W_0\odot(J+Z_1)$, and the next cycle $W_1\odot(J+Z_2)=W_0\odot(J+Z_1)\odot(J+Z_2)=W_0\odot(J+Z_1+Z_2+Z_1\odot Z_2)$.
With $Z_k=B_kA_k$,
\begin{equation*}
  Z_1+Z_2+Z_1\odot Z_2=\big[B_1,\ B_2,\ B_1\bullet B_2\big]\,\big[A_1;\ A_2;\ (A_1\T\bullet A_2\T)\T\big]
\end{equation*}
(\auxref{B.12}(c)), which factors through $\R^p$ with $p=2r+r^2$. Call the two factors $L\in\R^{m\times p}$ and
$L'\in\R^{p\times n}$. The rows of $L$ and the columns of $L'$ are points $(x,y,x\otimes y)$ of a variety $\mathcal V$ that
spans $\R^r\oplus\R^r\oplus\R^{r^2}$: a linear form vanishing on $\mathcal V$ vanishes on $(x,0,0)$, on $(0,y,0)$ and then on
every $(0,0,x\otimes y)$. Choosing points of $\mathcal V$ one at a time outside the span of the previous ones, $m$ rows can be
chosen to span $\min(m,p)$ dimensions, and likewise $n$ columns. If $p\le\min(m,n)$, then $L$ has full column rank and $L'$
full row rank, so $\rk LL'=p$. If $p>\min(m,n)$, say $m\le n$, choose $L$ of rank $m$, that is $L:\R^p\to\R^m$ onto; then
$L(\mathcal V)$ spans $\R^m$, so $n\ge m$ columns of $L'$ can be chosen with images spanning $\R^m$, and $\rk LL'=m$; the case
$n<m$ is the transpose. In every case rank $\min(2r+r^2,m,n)$ is attained, hence generic (\auxref{B.10}). Transported ReLoRA instead reaches $D_{W_0}(Z_1+Z_2)$, whose preimage $Z_1+Z_2$ has rank at most $2r$.
\end{proof}

% ======================================================================================================================
\subsection{Definition II.14}\label{proof:II.14}

\begin{restate}{Definition}{II.14}{covariance group}
For a method defined for every $W_0$ (a natural family), the \textbf{covariance group}
$\mathcal G_M\subseteq\GL_m\times\GL_n$ is the largest subgroup with $\Im M(gW_0h^{-1})=g\,\Im M(W_0)\,h^{-1}$. When
the method depends on random frames or on data, the identity is required in distribution, with the data transformed.
$\mathcal G_M$ is the \emph{geometry the method lives in}, in Klein's sense.
\end{restate}

The definition contains one claim: a largest such subgroup exists.

\begin{proof}
Let $\mathcal S$ be the set of $(g,h)\in\GL_m\times\GL_n$ with $\Im M(gWh^{-1})=g\,\Im M(W)\,h^{-1}$ for every $W$ in the
domain of the family (which we assume stable under $\GL_m\times\GL_n$, or at least under the pairs considered). The pair
$(I,I)$ lies in $\mathcal S$. If $(g_1,h_1),(g_2,h_2)\in\mathcal S$, then for every $W$
\begin{equation*}
  \Im M(g_1g_2Wh_2^{-1}h_1^{-1})=g_1\,\Im M(g_2Wh_2^{-1})\,h_1^{-1}=g_1g_2\,\Im M(W)\,h_2^{-1}h_1^{-1}.
\end{equation*}
If $(g,h)\in\mathcal S$, apply its identity at $g^{-1}Wh$: $\Im M(W)=g\,\Im M(g^{-1}Wh)\,h^{-1}$, so
$\Im M(g^{-1}Wh)=g^{-1}\Im M(W)h$. So $\mathcal S$ is a subgroup, every subgroup with the property is contained in it, and
$\mathcal G_M=\mathcal S$. For random frames or data, read each identity as equality of the laws of the random sets, the
pushforwards of the frame and data laws under $F\mapsto\Im M(W;F)$, with the data transformed by the action of $(g,h)$
(for instance $x\mapsto hx$, so $C\mapsto hCh\T$). The same computation applies, with each equality an equality of laws and
with $X\mapsto gXh^{-1}$ acting on laws by pushforward; it uses only that the data transformations compose as a group
action.
\end{proof}

% ======================================================================================================================
\subsection{Observation II.15}\label{proof:II.15}

\begin{restate}{Observation}{II.15}{descent through function-preserving reparametrisations}
Let $M=\prod_bM_b$ be a box-wise method: each $M_b$ is a member of a natural family with covariance group
$\mathcal G_{M_b}$, and no trainable parameters or random frames are shared across boxes. Let $G_{\mathrm{arch}}$ act on
$\Theta=\prod_b\Theta_b\times\Theta_{\mathrm{frozen}}$, preserving this split and acting on each adapted box by
$W_b\mapsto g_bW_bh_b^{-1}$ with $(g_b,h_b)\in\mathcal G_{M_b}$. Assume $f(g\theta,-)=f(\theta,-)$ on a
$G_{\mathrm{arch}}$-stable set containing $\Im M(\theta_0)$. Then $\Im M(g\theta_0)=g\cdot\Im M(\theta_0)$, and the set
of \emph{functions} reachable by $M$ from $g\theta_0$ equals the set reachable from $\theta_0$. For methods with random
frames, the equality holds in distribution, provided the frames are drawn independently per box. Boxes that share a
frame need a pathwise condition: a law-preserving map $T_{(g,h)}$ on frame space with
$\Im M_b(g_bWh_b^{-1};TF)=g_b\,\Im M_b(W;F)\,h_b^{-1}$ for all $W$ and $F$, with the same $T$ on every box that shares
$F$ (marginal covariance in law, as in \thmref{Definition}{II.14}, is not enough). For data-dependent methods, require
covariance pathwise in the data, and that the box-$b$ data at $g\theta_0$ is the \thmref{Definition}{II.14} transform
under $(g_b,h_b)$ of the data at $\theta_0$. This holds for intertwiner gauges such as hidden-unit permutations,
query--key and value--output head gauges and residual rotations. The statement concerns reachable sets only; training
trajectories and outcomes can still differ, for instance under Adam or non-isotropic initialisations.

Examples of $G_{\mathrm{arch}}$:
\begin{itemize}
  \item hidden-neuron permutations;
  \item the query--key gauge $(W_Q,W_K)\mapsto(gW_Q,g^{-\top}W_K)$ per head. With no positional encoding or an absolute
    one, $g\in\GL_{d_k}$. Under RoPE, $g$ must commute with every RoPE rotation, which leaves the per-frequency $2\times2$
    rotation--scalings $(\mathbb C^\times)^{d_k/2}$, shared across a GQA group. Under partial RoPE on $r$ of the $d_k$
    dimensions the group is $(\mathbb C^\times)^{r/2}\times\GL_{d_k-r}$. Under QK-norm, $g$ must in addition be
    orthogonal when the norm's $\varepsilon>0$ (conformal, $g\in\CO(d_k)$, when $\varepsilon=0$) and satisfy
    $g\T\Gamma_qR_t\Gamma_kg=\Gamma_qR_t\Gamma_k$ for every RoPE rotation $R_t$, where $\Gamma_q,\Gamma_k$ are the gains.
    OLMo~2 normalises $q$ and $k$ across all heads jointly;
  \item a single global orthogonal rotation $Q$ of the residual stream of a pre-norm transformer whose norm gains can all
    be fused into the linear layers they feed, after LayerNorm is converted to RMSNorm (SliceGPT's computational
    invariance; QuaRot's offline rotation). SliceGPT's per-block rotations $Q_\ell$ need extra residual matrices and
    change the architecture. Sandwich and post-sublayer norms (Gemma~2 and~3, OLMo~2) are excluded, because their gains
    cannot be fused and $\diag(\gamma)$ does not commute with $Q$. The comparison is between the gain-fused model and the
    fused-and-rotated model: gain fusion is a diagonal input rescaling outside $G_{\mathrm{arch}}$ that already changes
    PiSSA's split and block OFT's image (\thmref{Proposition}{II.16}), and converting LayerNorm to RMSNorm changes LoRA's
    reachable set on the residual writers.
\end{itemize}
\end{restate}

\begin{proof}
\emph{The deterministic statement.} No parameters are shared, so $\Im M(\theta)=\prod_b\Im M_b(W_b)\times\{\theta_{\mathrm{frozen}}\}$.
For $g\in G_{\mathrm{arch}}$, which preserves the split,
\begin{align*}
  \Im M(g\theta_0)&=\prod_b\Im M_b(g_bW_{0,b}h_b^{-1})\times\{g\theta_{0,\mathrm{frozen}}\}\\
  &=\prod_bg_b\,\Im M_b(W_{0,b})\,h_b^{-1}\times\{g\theta_{0,\mathrm{frozen}}\}=g\cdot\Im M(\theta_0),
\end{align*}
using $(g_b,h_b)\in\mathcal G_{M_b}$ in each factor. The reachable functions are then
\begin{equation*}
  \{\Phi_f(\theta):\theta\in\Im M(g\theta_0)\}=\{\Phi_f(g\theta):\theta\in\Im M(\theta_0)\}=\{\Phi_f(\theta):\theta\in\Im M(\theta_0)\},
\end{equation*}
since $f(g\theta,-)=f(\theta,-)$ for $\theta\in\Im M(\theta_0)$. (Stability of the set under $G_{\mathrm{arch}}$ is not needed.)

\emph{Independent frames.} With frames $F_b$ drawn independently per box, the law of the tuple of sets
$\Im M_b(g_bW_{0,b}h_b^{-1};F_b)$ is the product of the marginal laws, and by covariance in law each marginal equals the law of
$g_b\Im M_b(W_{0,b};F_b)h_b^{-1}$. So the tuple has the law of the tuple of sets $g_b\Im M_b(W_{0,b};F_b)h_b^{-1}$. The reachable function set is a fixed measurable function of the
image, and the function-level identity above holds pathwise, so the laws of the reachable function sets agree.

\emph{Shared frames.} Under the pathwise condition, $\Im M(g\theta_0;TF)=g\cdot\Im M(\theta_0;F)$ for every $F$, with the same $T$
in every box that uses $F$, and $TF$ has the law of $F$. Hence $\Im M(g\theta_0;F)$ and $g\cdot\Im M(\theta_0;F)$ have the same
law. Marginal covariance does not suffice. Take two boxes in $\R^{2\times1}$ sharing a Haar frame $F\in\Ort(2)$, with the
natural family $\Im M_b(W;F)=W+\R\,Fa(W)$, where $a(W)=(w_1^2,w_2^2)/\lVert(w_1^2,w_2^2)\rVert$ for $W\ne0$. For $g\in\Ort(2)$,
both $Fa(gW)$ and $gFa(W)$ are uniform on the circle, so each box is covariant in law under $(g,1)$. At $\theta_0=(e_1,e_2)$ the
two lines have directions $Fe_1$ and $Fe_2$, which are orthogonal, and so are those of $g\cdot\Im M(\theta_0;F)$. For $g$ the
rotation by $\pi/4$, $g\theta_0=\big((e_1+e_2)/\sqrt2,(e_2-e_1)/\sqrt2\big)$ has $a(ge_1)=a(ge_2)=(e_1+e_2)/\sqrt2$, so the two
lines of $\Im M(g\theta_0;F)$ are parallel. The joint laws differ, and the conclusion fails; the numerical check reports a
function-level instance.

\emph{Data.} Suppose each $M_b$ is covariant pathwise in its data,
\begin{equation*}
  \Im M_b\big(g_bWh_b^{-1};(g_b,h_b)\cdot D\big)=g_b\,\Im M_b(W;D)\,h_b^{-1},
\end{equation*}
and the box-$b$ data at $g\theta_0$ is $(g_b,h_b)\cdot D_b$. Then the deterministic argument applies with the data fixed. For an
intertwiner gauge the input of box $b$ in the network with weights $g\theta$ is $h_b$ times its input with weights $\theta$, for
every network input, so calibration inputs transform by $x\mapsto h_bx$ and second moments by $C\mapsto h_bCh_b\T$, as
\thmref{Definition}{II.14} requires. The examples below are of this kind: a hidden-unit permutation permutes the inputs of the
next box; the query--key and value--output gauges leave the inputs of $W_Q,W_K,W_V$ unchanged ($h=I$) and transform the input of
$W_O$; a residual rotation $Q$ rotates every normalised residual input by $Q$. Training is not covered: \thmref{Theorem}{III.12}
shows that Adam is not equivariant under general rotations of LoRA's gauge, and a Kaiming-uniform $A_0$ is not
rotation-invariant in law.

\emph{Hidden-neuron permutations.} For a permutation matrix $P$, $\sigma(Pz)=P\sigma(z)$ for an entrywise activation and
$P(u\odot v)=Pu\odot Pv$, so $(W_{\mathrm{up}},W_{\mathrm{gate}},W_{\mathrm{down}})\mapsto(PW_{\mathrm{up}},PW_{\mathrm{gate}},
W_{\mathrm{down}}P\T)$ preserves the block.

\emph{The query--key gauge.} Without positional rotation the score is $x\T W_Q\T W_Ky$, and since
$(gW_Q)\T g^{-\top}W_K=W_Q\T W_K$, every $g\in\GL_{d_k}$ works. Under RoPE the score between positions $s$ and $t$ is
$x\T W_Q\T R_{t-s}W_Ky$ with commuting rotations $R_\tau$ acting by angle $\tau\omega_j$ on the $j$-th coordinate plane. After
the gauge it is $x\T W_Q\T g\T R_{t-s}g^{-\top}W_Ky$. For $W_Q,W_K$ of full row rank, invariance for all inputs and positions
is equivalent to $g\T R_\tau g^{-\top}=R_\tau$ for every $\tau$, that is, $g\T$ commutes with every $R_\tau$. RoPE's frequencies
$\omega_j=\mathrm{base}^{-2j/d_k}$ are distinct and lie in $(0,1]\subset(0,\pi)$, so $R_1$ has the distinct eigenvalues
$e^{\pm\mathrm i\omega_j}$; a matrix commuting with $R_1$ preserves each coordinate plane and commutes with a nontrivial rotation
on it, so it is a rotation--scaling $a I_2+bJ$ there. The commutant is $(\mathbb C^\times)^{d_k/2}$ in the invertible case, and
it is closed under transposition, so the condition on $g\T$ is the same condition on $g$. In a GQA group one $W_K$ is paired with
several $W_Q$, so all of them use the same $g$. Under partial RoPE, $R_1$ is the identity on the $d_k-r$ unrotated coordinates and
has no eigenvalue $1$ on the others, so commuting matrices preserve this splitting, and the group is
$(\mathbb C^\times)^{r/2}\times\GL_{d_k-r}$.

\emph{QK-norm.} The score becomes $(\Gamma_qN(q))\T R_t(\Gamma_kN(k))$ with $q=W_Qx$, $k=W_Ky$ and
$N(v)=v\,(\lVert v\rVert^2/d_k+\varepsilon)^{-1/2}$. If $\varepsilon>0$ and $g$ is orthogonal, then $g^{-\top}=g$ and $N(gv)=gN(v)$,
so the score becomes $N(q)\T g\T\Gamma_qR_t\Gamma_kg\,N(k)$, which is invariant under the stated relation. If $\varepsilon=0$, then
$N$ removes scale, and $g=cO$ with $c>0$, $O$ orthogonal, gives $N(gq)=ON(q)$ and $N(g^{-\top}k)=ON(k)$; the relation is then
required of $O=g/c$. Conversely, let $W_Q,W_K$ have full row rank, so that $q$ and $k$ range over $\R^{d_k}$. Invariance at $t=0$
gives $\langle\Gamma_qN(gq),\Gamma_kN(g^{-\top}k)\rangle=\langle\Gamma_qN(q),\Gamma_kN(k)\rangle$ for all $q,k$. The values of $N$ span
$\R^{d_k}$ (they fill a ball when $\varepsilon>0$ and a sphere when $\varepsilon=0$), and the gains are invertible, so we may
choose $k_1,\dots,k_{d_k}$ for which the vectors $\Gamma_kN(g^{-\top}k_i)$ form a basis. Solving the resulting linear
equations gives $\Gamma_qN(gq)=A'\,\Gamma_qN(q)$ for a fixed matrix $A'$, so $N(gq)=A\,N(q)$ with $A=\Gamma_q^{-1}A'\Gamma_q$, and
$A$ is invertible because the values $N(gq)$ span $\R^{d_k}$. If $\varepsilon>0$, letting $q\to0$ (where $N(v)=v/\sqrt\varepsilon+O(\lVert v\rVert^3)$) gives $A=g$, and
then $\lVert gq\rVert=\lVert q\rVert$ for all $q$, so $g$ is orthogonal. If $\varepsilon=0$, $\lVert N(v)\rVert=\sqrt{d_k}$ for all
$v\ne0$, so $A$ is orthogonal, and $gq$ is a positive multiple of $Aq$ for every $q$, so $A^{-1}g$ has every vector as an
eigenvector and $g=cA$ with $c>0$. With $g$ orthogonal (or $O=g/c$), invariance for all $t$ and all $N(q),N(k)$ forces the
stated relation.

\emph{Residual rotations.} In a pre-norm transformer with gains fused and RMSNorm $N(h)=h\,(\lVert h\rVert^2/d+\varepsilon)^{-1/2}$,
replace the residual stream $h$ by $Qh$: embeddings $E\mapsto QE$, the writers $W_O,W_{\mathrm{down}}\mapsto QW_O,QW_{\mathrm{down}}$,
and every reader (the $q,k,v$, gate and up projections and the head) $W\mapsto WQ\T$. Since $\lVert Qh\rVert=\lVert h\rVert$,
$N(Qh)=QN(h)$, so each reader sees $WQ\T QN(h)=WN(h)$, and the function is unchanged. A gain that is not fused would sit
between $N$ and the reader as $\diag(\gamma)$, and $\diag(\gamma)Q\ne Q\diag(\gamma)$ unless $\gamma$ is constant. A
post-sublayer norm writes $h+\diag(\gamma)N(W_Oz)$, and $\diag(\gamma)N(QW_Oz)=\diag(\gamma)QN(W_Oz)\ne Q\diag(\gamma)N(W_Oz)$,
so the rotation does not preserve the function. LayerNorm subtracts the mean, the projection $I-\mathbf 1\mathbf 1\T/d$, which
commutes with $Q$ only if $Q\mathbf 1=\pm\mathbf 1$; SliceGPT therefore folds the projection into the writers and uses RMSNorm.
In the original model a LayerNorm ignores the component of the residual stream along $\mathbf 1$, so LoRA's updates of a writer
reach functions only through $(I-\mathbf 1\mathbf 1\T/d)\Delta W$; in the converted model the writer is
$(I-\mathbf 1\mathbf 1\T/d)W_O$ and an update $\Delta W$ with a component along $\mathbf 1$ is seen by every later RMSNorm, so the
reachable function sets differ. Gain fusion replaces a reader $W$ by $W\diag(\gamma)$, a non-conformal input rescaling when
$\gamma$ is not constant; by \thmref{Proposition}{II.16} it changes PiSSA's split, and it changes block OFT's image
$\{W\diag(\gamma)R\}$ against $\{WR\diag(\gamma)\}$ unless $\diag(\gamma)$ commutes with the block-orthogonal group, that is
unless $\gamma$ is constant on each block. Consecutive SliceGPT rotations $Q_{\ell-1}\ne Q_\ell$ require the matrix
$Q_{\ell-1}\T Q_\ell$ on the residual path, a new operation.
\end{proof}

% ======================================================================================================================
\subsection{Proposition II.16}\label{proof:II.16}

\begin{restate}{Proposition}{II.16}{covariant data-aware splits}
Let $C=\mathbb E[xx\T]\succ0$ be the input second moment, so that under $x\mapsto hx$ we have $W\mapsto Wh^{-1}$ and
$C\mapsto hCh\T$. Let $LL\T=C$ and $p:=\min(m,n)$. Write $[M]_r$ for the rank-$r$ truncated SVD,
$\langle M\rangle_r:=M-[M]_{p-r}$ for the bottom $r$ singular components, and $\CO(n)=\R_{>0}\cdot\Ort(n)$. Assume
$1\le r<p$ and the spectral gap each split needs: $\sigma_r>\sigma_{r+1}$ of $W_0L$ in (a); in (b),
$\sigma_r>\sigma_{r+1}$ of $W_0C$ or $W_0$ for top-$r$ splits, and $\sigma_{p-r}>\sigma_{p-r+1}$ for bottom-$r$ splits.
Covariance groups are taken over the generic set where these gaps hold.

(a) The whitened split
\begin{equation*}
  X^*=\arg\min_{\rk X\le r}\lVert(W_0-X)L\rVert_F=[W_0L]_rL^{-1}=\Pi_rW_0,
\end{equation*}
where $\Pi_r$ is the top-$r$ eigenprojector of the output second moment $W_0CW_0\T$, is $\GL_n$-covariant on the input
side and $\CO(m)$-covariant on the output side; its joint group is $\CO(m)\times\GL_n$. This is the classical
reduced-rank-regression (SVD-LLM) solution.

(b) The formula-level CorDA splits, IPM ($X=[W_0C]_rC^{-1}$, the minimiser of $\lVert(W_0-X)C\rVert_F$) and KPM (the
bottom-$r$ components of $W_0C$ mapped back by $C^{-1}$, $X=\langle W_0C\rangle_rC^{-1}$), with or without CorDA's
trace-relative damping $C+\kappa(\tr C/n)I$, and the PiSSA and MiLoRA splits ($[W_0]_r$ and the bottom-$r$ part
$\langle W_0\rangle_r$), have input-side covariance group exactly $\CO(n)$; jointly, $\CO(m)\times\CO(n)$.

(c) \emph{Scope.} An absolute damping $C+\lambda I$ reduces the input-side group of (a) and of (b) to $\Ort(n)$. The
official and HF CorDA code divide each calibration sample by the absolute value of its largest entry before
accumulating $C$, which keeps only positive scalings and coordinate permutations, so (b) describes the CorDA formulas and
does not cover shipped CorDA under rotations.
\end{restate}

\begin{proof}
A split $X=X(W_0,C)$ of rank at most $r$ yields the method with frozen residual $W_0-X$ and image $(W_0-X)+\Mr$. Under $(g,h)$,
with $C'=hCh\T$, covariance of images means $(gW_0h^{-1}-X(gW_0h^{-1},C'))+\Mr=g(W_0-X(W_0,C))h^{-1}+\Mr$, using $g\Mr h^{-1}=\Mr$.
Since $r<p$, the translation stabiliser of $\Mr$ is $\{0\}$ (equation~\eqref{eq:translation-stabiliser}), so this is equivalent to
\begin{equation*}
  X(gW_0h^{-1},hCh\T)=g\,X(W_0,C)\,h^{-1}.
\end{equation*}

\emph{Lemma.} Let $1\le s<p$, $g\in\GL_m$ and $k\in\GL_n$. Then $[gMk]_s=g[M]_sk$ for every $M\in\R^{m\times n}$ for which both
truncations are unique iff $g\in\CO(m)$ and $k\in\CO(n)$. ``If'': for $g=aP$ and $k=bQ$ with $a,b>0$ and $P,Q$ orthogonal, an
SVD $M=\sum_i\sigma_iu_iv_i\T$ gives the SVD $gMk=\sum_i(ab\sigma_i)(Pu_i)(Q\T v_i)\T$ with the same ordering, so
$[gMk]_s=g[M]_sk$. ``Only if'': let $U_s\subseteq\R^m$ be any $s$-dimensional subspace and $y\perp U_s$ a unit vector, and
likewise $V_s\subseteq\R^n$ and $x\perp V_s$ (possible because $s<p$). Let $M_s$ have rank $s$ with column space $U_s$ and row
space $V_s$, and for small $t>0$ put $M=M_s+t\,yx\T$. Then $[M]_s=M_s$, and $gMk=gM_sk+t\,(gy)(k\T x)\T$, whose truncation is
unique for small $t$ because $gM_sk$ has rank $s$. The identity gives $[gMk]_s=gM_sk$, so the residual $t(gy)(k\T x)\T$ of the
truncation must have its column space orthogonal to that of $[gMk]_s$ and its row space orthogonal to that of $[gMk]_s$
(\auxref{B.11}(a): $[X]_s$ and $X-[X]_s$ are built from disjoint sets of singular vectors). So $gy\perp gU_s$ and
$k\T x\perp k\T V_s$. Since every pair of orthogonal unit vectors $u,y$ occurs with $u\in U_s$ and $y\perp U_s$, $g$ maps
orthogonal vectors to orthogonal vectors. Then for orthonormal eigenvectors $e,f$ of $g\T g$ with eigenvalues $\alpha,\beta$,
$0=\langle g(e+f),g(e-f)\rangle=\alpha-\beta$, so $g\T g$ is scalar and $g\in\CO(m)$; likewise $k\T\in\CO(n)$.

\emph{(a)} By reduced-rank regression (\auxref{B.11}(c)), $X^*=[W_0L]_rL^{-1}$, unique under the gap, and $[W_0L]_r=\Pi_rW_0L$
with $\Pi_r$ the top-$r$ eigenprojector of $W_0LL\T W_0\T=W_0CW_0\T$; so $X^*=\Pi_rW_0$, which does not depend on the choice of
$L$. Under $(g,h)$ the output second moment becomes $gW_0h^{-1}\,hCh\T\,h^{-\top}W_0\T g\T=gW_0CW_0\T g\T$, and with the square
root $hL$ of $hCh\T$,
\begin{equation*}
  X^*(gW_0h^{-1},hCh\T)=[gW_0L]_r\,L^{-1}h^{-1}.
\end{equation*}
This equals $gX^*h^{-1}=g[W_0L]_rL^{-1}h^{-1}$ for all $W_0$ (with gaps) iff $[gM]_r=g[M]_r$ for all $M=W_0L$, which range
over all matrices with a gap. By the lemma with $k=I$ this holds iff $g\in\CO(m)$, with no condition on $h$. So the joint group
is $\CO(m)\times\GL_n$, and its input and output parts are $\GL_n$ and $\CO(m)$.

\emph{(b), undamped.} For IPM, $X(W,C)=[WC]_rC^{-1}$ minimises $\lVert(W-X)C\rVert_F$ (\auxref{B.11}(c) with $L=C$). Under
$(g,h)$, with $M=W_0C$,
\begin{equation*}
  X(gW_0h^{-1},hCh\T)=[gW_0Ch\T]_rh^{-\top}C^{-1}h^{-1}=[gMh\T]_r\,h^{-\top}C^{-1}h^{-1},
\end{equation*}
and covariance asks this to equal $g[M]_rC^{-1}h^{-1}$, that is $[gMh\T]_r=g[M]_rh\T$. As $W_0$ varies, $M=W_0C$ ranges over all
matrices, so by the lemma covariance holds iff $g\in\CO(m)$ and $h\T\in\CO(n)$, that is $h\in\CO(n)$. KPM:
$X=\langle W_0C\rangle_rC^{-1}=W_0-[W_0C]_{p-r}C^{-1}$, and the same computation with $s=p-r\in[1,p-1]$ gives the condition
$[gMh\T]_{p-r}=g[M]_{p-r}h\T$, with the gap $\sigma_{p-r}>\sigma_{p-r+1}$. PiSSA: $[gW_0h^{-1}]_r=g[W_0]_rh^{-1}$ for all $W_0$ iff
$g\in\CO(m)$ and $h^{-1}\in\CO(n)$; MiLoRA: $\langle W\rangle_r=W-[W]_{p-r}$ reduces to the lemma at $s=p-r$. In each case the
input-side group (with $g=I$) is $\CO(n)$ and the joint group is $\CO(m)\times\CO(n)$.

\emph{(b), trace-relative damping.} Write $C_\kappa=C+\kappa(\tr C/n)I$ with $\kappa>0$ and consider $X_\kappa(W,C)=X(W,C_\kappa)$
for IPM (KPM is identical with $s=p-r$). For $h=cQ\in\CO(n)$,
$hC_\kappa h\T=c^2QCQ\T+\kappa(\tr C/n)c^2I=(hCh\T)_\kappa$, since $\tr(c^2QCQ\T)=c^2\tr C$; so damping commutes with $\CO(n)$ and
the damped splits inherit the covariance of the undamped ones. Conversely, let $(g,h)$ be in the joint group, put
$S_1=(hCh\T)_\kappa$, $S_2=C_\kappa$ and $N=W_0h^{-1}$. Covariance reads $[gNS_1]_rS_1^{-1}=g[NhS_2]_rS_2^{-1}h^{-1}$. With
$M=NS_1$, which ranges over all matrices, and $G=S_1^{-1}hS_2$, this becomes $[gM]_rG=g[MG]_r$, or $[gM'G^{-1}]_r=g[M']_rG^{-1}$
with $M'=MG$. By the lemma, $g\in\CO(m)$ and $G\in\CO(n)$ for every $C\succ0$, that is,
\begin{equation*}
  h\,S_2^2\,h\T=c(C)^2\,S_1^2\qquad\text{for every }C\succ0 .
\end{equation*}
Both sides are continuous in $C$, and $c(C)^2=\tr(hS_2^2h\T)/\tr(S_1^2)$ is too, so the identity extends to $C=vv\T$ with $v\ne0$.
There, with $w=hv$, $a=\kappa\lVert v\rVert^2/n>0$ and $b=\kappa\lVert w\rVert^2/n$,
$S_2^2=(\lVert v\rVert^2+2a)vv\T+a^2I$ and $S_1^2=(\lVert w\rVert^2+2b)ww\T+b^2I$, so
\begin{equation*}
  hh\T=\alpha_wI+\beta_w\,ww\T\qquad(\alpha_w=c^2b^2/a^2),
\end{equation*}
for every $w\in\R^n$, since $w=hv$ ranges over all vectors. So $hh\T-\alpha_wI$ has rank at most one for every $w$. If $hh\T$
were not scalar, it would have an eigenvalue $\alpha$ of multiplicity $n-1$ with $hh\T=\alpha I+\gamma ee\T$, $\gamma\ne0$; taking
$w$ neither parallel nor orthogonal to $e$, the matrix $(\alpha-\alpha_w)I+\gamma ee\T$ would have rank at most one with range
along $w$, which is impossible (for $n\ge3$ it has rank at least $n-1\ge2$ unless $\alpha_w=\alpha$, and then its range is
$\R e\ne\R w$; for $n=2$ its range is $\R e$ or $e^\perp$). So $hh\T$ is scalar and $h\in\CO(n)$. The input-side group of the damped
splits is therefore exactly $\CO(n)$, and the joint group $\CO(m)\times\CO(n)$.

\emph{(c), absolute damping.} $Q(C+\lambda I)Q\T=QCQ\T+\lambda I$ for $Q\in\Ort(n)$, so the damped splits of (a) and (b) are
$\Ort(n)$-covariant. We show that no other $h$ is in their input-side groups.

For (b): the undamped splits are invariant under $C\mapsto tC$, so the damped split at $(W,tC)$ equals the undamped split at
$(W,C+(\lambda/t)I)$, which tends to the undamped split at $(W,C)$ as $t\to\infty$ wherever the gap holds, because a truncation is
continuous where it is unique (it is the unique minimiser of a continuous function over the closed set $\MM_{\le s}$). The same
holds at the transformed data $(Wh^{-1},thCh\T)$. So every $h$ in the damped group lies in the undamped group $\CO(n)$.

For (a): with $S=(h\T h)^{-1}$, the transformed damped moment is $W_0h^{-1}(hCh\T+\lambda I)h^{-\top}W_0\T=W_0(C+\lambda S)W_0\T$, so
covariance reads $\Pi_r\big(W_0(C+\lambda S)W_0\T\big)W_0=\Pi_r\big(W_0(C+\lambda I)W_0\T\big)W_0$. Both projectors have range in
$\operatorname{col}W_0$, onto which $W_0$ maps, so this is equality of the projectors. Replacing $C$ by $tC$ and letting $t\to0$
gives $\Pi_r(W_0SW_0\T)=\Pi_r(W_0W_0\T)$ for every $W_0$ for which both have a gap at $r$. Suppose $S$ is not scalar, with
orthonormal eigenvectors $e,f$ for eigenvalues $\alpha>\beta$, and complete them to an orthonormal eigenbasis $e,f,g_3,\dots,g_n$ of
$S$. Choose orthonormal $z_1,\dots,z_p$ with $z_r=\cos\phi\,e+\sin\phi\,f$, $z_{r+1}=-\sin\phi\,e+\cos\phi\,f$ for some
$\phi\notin\tfrac\pi2\mathbb Z$, and the other $z_i$ among the $g_j$; put $W_0=\sum_{i\le p}d_ie_iz_i\T$ with $d_1>\dots>d_p>0$.
Then $W_0W_0\T=\diag(d_1^2,\dots,d_p^2)$ (padded), with top-$r$ eigenspace $\spanop\{e_1,\dots,e_r\}$, while $W_0SW_0\T$ is
diagonal except for the entry $d_rd_{r+1}(\beta-\alpha)\sin\phi\cos\phi\ne0$ in positions $(r,r+1)$ and $(r+1,r)$. Taking $d_i$ for
$i<r$ much larger and $d_i$ for $i>r+1$ much smaller than $d_r,d_{r+1}$, the top-$r$ eigenspace of $W_0SW_0\T$ is
$\spanop\{e_1,\dots,e_{r-1}\}$ plus the top eigenvector of the $2\times2$ block, which is not $e_r$; both matrices have a gap at
$r$, and the projectors differ. Hence $S$ is scalar and $h\in\CO(n)$.

It remains to exclude $h=cQ$ with $c\ne1$; since $Q$ is in the group, it suffices to exclude $h=cI$. For (a), $h=cI$ turns the
condition into $\Pi_r(W_0(C+\lambda c^{-2}I)W_0\T)=\Pi_r(W_0(C+\lambda I)W_0\T)$. Replacing $C$ by $tC$, the function
$P(\mu)=\Pi_r(W_0CW_0\T+\mu W_0W_0\T)$ would satisfy $P(\mu)=P(c^{-2}\mu)$ for all $\mu>0$, hence $P(\mu)=P(c^{-2j}\mu)$ for all
$j\in\mathbb Z$, and letting $j\to\pm\infty$ gives $P(0)=P(\infty)$, the top-$r$ eigenprojectors of $W_0CW_0\T$ and of $W_0W_0\T$
(both limits exist under gaps). Take $W_0$ rectangular-diagonal with entries $d_1>\dots>d_p>0$ and $C=\diag(c_i)$ with
$c_i=i/d_i^2$ for $i\le p$ and $c_i=1$ beyond. Then $W_0W_0\T$ has top-$r$ eigenspace $\spanop\{e_1,\dots,e_r\}$ and
$W_0CW_0\T=\diag(1,2,\dots,p)$ (padded) has top-$r$ eigenspace $\spanop\{e_{p-r+1},\dots,e_p\}$, which differ since $r<p$. For (b),
$h=cI$ gives, by scale invariance and linearity in $W$, $X(W_0,C+\lambda c^{-2}I)=X(W_0,C+\lambda I)$, and the same iteration
applied to $G(\mu)=X(W_0,C+\mu I)$ gives $G(0)=G(\infty)$ with $G(0)=X(W_0,C)$ and $G(\infty)=X(W_0,I)$. With the same $W_0$ and
$C$, $W_0C$ is rectangular-diagonal with the increasing entries $i/d_i$, so for IPM $G(0)=\sum_{i>p-r}d_iE_{ii}$ while
$G(\infty)=[W_0]_r=\sum_{i\le r}d_iE_{ii}$; for KPM the roles of top and bottom are exchanged, with the same conclusion. So
$cI$ is not in the group, and the damped input-side groups are exactly $\Ort(n)$.

\emph{(c), per-sample normalisation.} The shipped code divides a sample $x$ by $\nu(x)=\lvert\max_jx_j\rvert$. If $h=cP$ with $c>0$
and $P$ a permutation matrix, $\nu(hx)=c\,\nu(x)$, so the accumulated $C$ transforms as $hCh\T$ up to the factor $c^{-2}$, to which
the splits are insensitive, and $h\in\CO(n)$. Conversely, suppose $\nu(hx)=c\,\nu(x)$ for all $x$ and some $c>0$ ($n\ge2$). For
$x=-e_j$, $\nu(x)=0$, so $\max_i(-h_{ij})=0$ and $h\ge0$ entrywise; for $x=e_j$, $\max_ih_{ij}=c$; for $x=\mathbf 1$, every row sum of
$h$ is at most $c$. So a row $i_j$ with $h_{i_jj}=c$ has no other nonzero entry, the rows $i_j$ are distinct, and $h=cP$. Under a
rotation the normalisation does not commute with $h$, and the shipped split is not covariant. For instance, with $n=m=2$, $r=1$,
samples $x_1=e_1$ and $x_2=e_1+e_2$ (so $\nu=1$ for both and $C=x_1x_1\T+x_2x_2\T=\left(\begin{smallmatrix}2&1\\1&1\end{smallmatrix}
\right)$, up to a factor), and $h=Q$ the rotation by $\pi/4$: the rotated samples are $(1,1)/\sqrt2$ and $(0,\sqrt2)$, normalised to
$(1,1)$ and $(0,1)$, so the shipped moment is $C'=\left(\begin{smallmatrix}1&1\\1&2\end{smallmatrix}\right)$, while
$QCQ\T=\tfrac12\left(\begin{smallmatrix}1&1\\1&5\end{smallmatrix}\right)$. With $W_0=Q$, the shipped IPM split at the rotated data
is $X(I,C')=[C']_1C'^{-1}$, the top eigenprojector of $C'$, while covariance would require $X(Q,C)Q\T=Q[C]_1C^{-1}Q\T$, the top
eigenprojector of $QCQ\T$. The two symmetric matrices do not commute, so their top eigenvectors differ, and covariance fails.
\end{proof}
